%%%%%%%%%%%%%%%%%%%%%%%%%%%%%%%%%%%%%%%%%%%%%%%%%%%%%%%%%%%%%%%%%%%%%%
% Lớp 11 | Đạo hàm | Các quy tắc đạo hàm
% Dạng 1D7?2-1: Tính đạo hàm
% Số câu: 275  (sinh tự động bởi tach_cau.py; từ nay thầy sửa trực tiếp file này)
%%%%%%%%%%%%%%%%%%%%%%%%%%%%%%%%%%%%%%%%%%%%%%%%%%%%%%%%%%%%%%%%%%%%%%

%%%%%%%%%%%%%%%% TRẮC NGHIỆM 4 PHƯƠNG ÁN %%%%%%%%%%%%%%%%

% nguồn: Ok_Toan11_TracNghiem_Lan2
\begin{ex}%[1D7N2-1]%[ma:dada4760]
	Cho hàm số $f(x)$ là hàm số trên  $\mathbb{R}$ xác định bởi $f(x)=2x$. Chọn khẳng định đúng.
	 \choice
		{\True  $f'(x)=2$}
		 {$f'(x)=1$}
		 {$f'(x)=x$}
		{$f'(x)$ không  tồn tại}
	\loigiai{Ta có $f'(x)=(2x)'=2$.} 
\end{ex}

% nguồn: Ok_Toan11_TracNghiem_Lan2
\begin{ex}%[1D7N2-1]%[ma:056b0ee5]
	Tìm đạo hàm của hàm số  $y=2\, 022^{x}$.
	 \choice
		{$y'=x \cdot  2\, 022^{x-1}$}
		{$y'=\dfrac{2\,  022^{x}}{\ln~2\, 022}$}
		{\True $y'=2\, 022^{x} \cdot \ln~2\,  022$}
		{$2\,  022$}
	\loigiai{ 
		Ta có  $y'=2\, 022^{x} \cdot \ln~2\, 022$.
	 }
\end{ex}

% nguồn: Ok_Toan11_TracNghiem_Lan2
\begin{ex}%[1D7N2-1]%[ma:0a525524]
	Đạo  hàm của hàm số $y=6^{x}$ là
	\choice 
		{$y'=\dfrac{6^{x}}{\ln~6}$} 
		{\True  $y'=6^{x}\ln~6$}
		 {$y'=6^{x}$}
		{$y'=x \cdot  6^{x-1}$}
	\loigiai{ 
		Ta  có $y=6^{x} \Rightarrow y'=6^{x}\ln~6$.
	} 
\end{ex}

% nguồn: Ok_Toan11_TracNghiem_Lan2
\begin{ex}%[1D7N2-1]%[ma:11ec4bdf]
	Khẳng định  nào sau đây là đúng?
	 \choice
		 {$\left(\sqrt{x}\right)'=\dfrac{1}{\sqrt{x}}$}
		 {$\left(\dfrac{1}{x^2}\right)'=\dfrac{1}{x^2}$}
		{$(x)'=0$} 
		{\True $(kx)'=k$, với $k$ là  hằng số}
	 \loigiai{Ta có $(kx)'=k$, với $k$ là hằng số.
	 }
\end{ex}

% nguồn: Ok_Toan11_TracNghiem_Lan2
\begin{ex}%[1D7N2-1]%[ma:e8979699]
	Tính đạo hàm của hàm  số $y=2^{x}$.
	 \choice
		{$y'=x  \cdot 2^{x-1}$}
		 {\True $y'=2^{x}\ln~2$}
		 {$y'=2^{x}$}
		 {$y'=2^{x}\ln~x$}
	 \loigiai{
		 Ta có $y'=2^{x}\ln~2$.
	} 
\end{ex}

% nguồn: Ok_Toan11_TracNghiem_Lan2
\begin{ex}%[1D7N2-1]%[ma:ed27ca5e]
	Đạo hàm của hàm số $y=-2x^3+3$  là
	\choice 
		 {$y'=-6x$}
		 {$y'=-6x^2+3$}
		{\True $y'=-6x^2$} 
		 {$y'=-3x^2$}
	\loigiai{Ta có $y'=-3\cdot  2\cdot x^2=-6x^2$.
	 }
\end{ex}

% nguồn: Ok_Toan11_TracNghiem_Lan2
\begin{ex}%[1D7N2-1]%[ma:f54cf2fc]
	 Tìm đạo hàm của hàm số $y=\log x$.
	 \choice
		 {$y'=\dfrac{\ln 10}{x}$}
		{\True $y'=\dfrac{1}{x\ln 10}$} 
		{$y'=\dfrac{1}{10\ln x}$} 
		 {$y'=\dfrac{1}{x}$}
	 \loigiai{
		Áp dụng công thức $\left(\log_a x\right)'=\dfrac{1}{x\ln a}$,  ta được $y'=\dfrac{1}{x \ln 10}$.
	 }
\end{ex}

% nguồn: Ok_Toan11_TracNghiem_Lan2
\begin{ex}%[1D7N2-1]%[ma:02e61618]
	 Tính đạo hàm của hàm số $y=13^{x}$.
	\choice 
		{$y'=\dfrac{13^{x}}{\ln13}$} 
		{$y'=x  \cdot 13^{x-1}$}
		 {\True $y'=13^{x}\ln~13$}
		{$y'=13^{x}$} 
	\loigiai{ 
		 Ta có $y'=13^{x}\ln~13$.
	 }
\end{ex}

% nguồn: Ok_Toan11_TracNghiem_Lan2
\begin{ex}%[1D7N2-1]%[ma:e97329ea]
	Tìm đạo hàm của hàm số  $y=\pi^{x}$.
	\choice 
		 {$y'=x \cdot \pi^{x-1}\ln~\pi$}
		{\True $y'=\pi^{x}\ln~\pi$} 
		 {$y'=\dfrac{\pi^{x}}{\ln~\pi}$}
		{$y'=x \cdot \pi^{x-1}$} 
	\loigiai{ 
		Ta có áp dụng  $(a^{x})'=a^{x} \cdot \ln~a$ (với $a>0$, $a\ne1$).
	} 
\end{ex}

% nguồn: Ok_Toan11_TracNghiem_Lan2
\begin{ex}%[1D7N2-1]%[ma:e4d98d99]
	Trên khoảng $(0;+\infty)$, hàm số $y=\log_3  x$ có đạo hàm là
	 \choice
		{$y'=\dfrac{x}{\ln 3}$} 
		{$y'=x\ln  3$}
		{\True $y'=\dfrac{1}{x\ln 3}$} 
		 {$y'=\dfrac{\ln 3}{x}$}
	\loigiai{ 
		Ta có  $y'=\left(\log_3 x\right)'=\dfrac{1}{x\ln 3}$.
	} 
\end{ex}

% nguồn: Ok_Toan11_TracNghiem_Lan2
\begin{ex}%[1D7N2-1]%[ma:88438998]
	Trên khoảng $(0;+\infty)$, đạo hàm của hàm số $y=\log _{\sqrt{3}}  x$ là
	 \choice
		{$y'=\dfrac{\sqrt{3}}{x\ln 3}$} 
		{$y'=\dfrac{1}{2x\ln 3}$} 
		{\True  $y'=\dfrac{1}{x\ln \sqrt{3}}$}
		{$y'=\dfrac{\sqrt{3}}{x}$} 
	\loigiai{ 
		Ta có $y'=\dfrac{1}{x\ln \sqrt{3}}$. 
	} 
\end{ex}

% nguồn: Ok_Toan11_TracNghiem_Lan2
\begin{ex}%[1D7N2-1]%[ma:4c639c3e]
	Đạo hàm của hàm số  $y=5^{x}$ là
	\choice 
		{$y'=-5^{x}\ln~5$} 
		{\True  $y'=5^{x}\ln~5$}
		 {$y'=\dfrac{5^{x}}{\ln~5}$}
		 {$y'=-\dfrac{5^{x}}{\ln~5}$}
	\loigiai{ 
		Ta  có $y'=5^{x}\ln~5$.
	} 
\end{ex}

% nguồn: Ok_Toan11_TracNghiem_Lan2
\begin{ex}%[1D7N2-1]%[ma:d2b5f599]
	 Tính đạo hàm của hàm số $y=3^{x}$.
	\choice 
		{$y'=3^{x}$} 
		 {$y'=\dfrac{3^{x}}{\ln~3}$}
		{\True $y'=3^{x}\ln~3$} 
		{$y'=x \cdot 3^{x-1}$} 
	 \loigiai{
		Ta có áp dụng  công thức $\left(a^{x}\right)'=a^{x} \cdot \ln~a$.
	 }
\end{ex}

% nguồn: Ok_Toan11_TracNghiem_Lan2
\begin{ex}%[1D7N2-1]%[ma:0a97f02b]
	Trên  khoảng $(0;+\infty)$, đạo hàm của hàm số $y=\log_2 x$ là
	 \choice
		{\True $y'=\dfrac{1}{x\ln  2}$}
		{$y'=\dfrac{\ln  2}{x}$}
		{$y'=\dfrac{1}{x}$} 
		{$y'=\dfrac{1}{2x}$} 
	\loigiai{ 
		Áp dụng công thức $\left(\log_a x\right)'=\dfrac{1}{x\ln a}$. \\ 
		Ta có  $y'=\dfrac{1}{x\ln 2}$.
	 }
\end{ex}

% nguồn: Ok_Toan11_TracNghiem_Lan2
\begin{ex}%[1D7N2-1]%[ma:63b1d3dc]
	Hàm số $y=-\dfrac{1}{\sin^2x}$ là đạo hàm của hàm số  nào dưới đây?
	 \choice
		{$y=\tan  x$}
		 {\True $y=\cot x$}
		{$y=-\cot  x$}
		 {$y=\dfrac{1}{\sin x}$}
	\loigiai{Ta có  $(\cot x)'=-\dfrac{1}{\sin^2x}$.
	} 
\end{ex}

% nguồn: Ok_TracNghiemLop11_Lan1
\begin{ex}%[1D7N2-1]%[ma:315556db]
	Tính đạo hàm  của hàm số $y=\dfrac{x-1}{2x-1}$.
	\choice 
		 {\True $y'=\dfrac{1}{(2x-1)^2} $}
		{$y'=\dfrac{1}{2x-1}  $}
		 {$y'=\dfrac{1}{(2x+1)^2} $}
		 {$y'=-\dfrac{3}{(2x-1)^2} $}
	\loigiai{ 
		 Ta có $y'=\dfrac{(x-1)'(2x-1)-(2x-1)'(x-1)}{(2x-1)^2}=\dfrac{1}{(2x-1)^2}$.
	} 
\end{ex}

% nguồn: Ok_TracNghiemLop11_Lan1
\begin{ex}%[1D7N2-1]%[ma:4da3aa45]
	Đạo  hàm của hàm số $y=\log_3\left(x+1\right)$ là
	\choice 
		{$y'=-\dfrac{1}{\ln 3}$} 
		 {\True $y'=\dfrac{1}{\left(x+1\right)\ln 3}$}
		 {$y'=\dfrac{1}{\left(x+1\right)}$}
		{$y'=\dfrac{x+1}{\ln 3}$} 
	\loigiai{ 
	Đạo hàm của hàm số $y=\log_3\left(x+1\right)$  là $y'=\dfrac{1}{\left(x+1\right)\ln 3}$.}
\end{ex}

% nguồn: Ok_TracNghiemLop11_Lan1
\begin{ex}%[1D7N2-1]%[ma:d3911af5]
	Khẳng định nào sau đây  \textbf{sai}?
	\choice 
		{$y=x\Rightarrow  y'=1$}
		{$y=x^3  \Rightarrow y'=3x^2$}
		{\True  $y=x^5 \Rightarrow y'=5x$}
		{$y=x^4 \Rightarrow y'=4x^3$} 
	\loigiai{ 
		Ta có $y=x^n \Rightarrow y'=n\cdot  x^{n-1}$, $\forall n\in \mathbb{N}^{*}$.\\
		 Nên $y=x^5 \Rightarrow y'=5x^4$.
	} 
\end{ex}

% nguồn: Ok_TracNghiemLop11_Lan1
\begin{ex}%[1D7N2-1]%[ma:3c77097f]
	Đạo hàm của hàm số $y=x^4-3x^2-5x+2017$ là 
	\choice 
		 {$y'=x^3-6x-5$}
		{\True  $y'=4x^3-6x-5$}
		{$y'=4x^3-6x+2017$} 
		{$y'=4x^3+6x-5$} 
	\loigiai{ 
		 Ta có $y'=4x^3-6x-5$.
	 }
\end{ex}

% nguồn: Ok_TracNghiemLop11_Lan1
\begin{ex}%[1D7N2-1]%[ma:432bed97]
	Đạo hàm của hàm số $ y=\sin  2x-\cos x$ là
	\choice 
		{$y'=2\cos  x+\sin x$}
		{$y'=\cos 2x+\sin  x$}
		{\True  $y'=2\cos 2x+\sin x$}
		{$y'=2\cos x-\sin  x$}
	\loigiai{ 
		$y'=(\sin 2x)'-(\cos x)'  = 2\cos2x+\sin x$.
	 }
\end{ex}

% nguồn: Ok_TracNghiemLop11_Lan1
\begin{ex}%[1D7N2-1]%[ma:dd8fb73b]
	Cho  hàm số $y=x^{2016}$. Tính $y'(1)$.
	 \choice
		{$+\infty$} 
		{$2016!$} 
		{\True  $2016$}
		 {$0$}
	\loigiai{Ta có $y'=2016  x^{2015}\Rightarrow y'(1)=2016$.
	} 
\end{ex}

% nguồn: Ok_TracNghiemLop11_Lan1
\begin{ex}%[1D7N2-1]%[ma:76b0a36c]
	Cho các hàm số $y=\sin x,y=\cos x,y=\tan x,y=\cot x$ có đạo hàm trên tập xác định của nó.  Mệnh đề nào sau đây \textbf{sai}?
	 \choice
		{${\left(\tan  x\right)}'=\dfrac{1}{\cos^2x}$}
		{${\left(\sin x\right)}'=\cos  x$}
		 {${\left(\cos x\right)}'=-\sin x$}
		{\True ${\left(\cot x\right)}'=\dfrac{1}{\sin^2x}$} 
	 \loigiai{
	Ta có ${\left(\cot  x\right)}'=\dfrac{-1}{\sin^2x}$.}
\end{ex}

% nguồn: Ok_TracNghiemLop11_Lan1
\begin{ex}%[1D7N2-1]%[ma:26806564]
	Với mọi  $x\in\mathbb{R}$, đạo hàm của hàm số $y=2\sin x-\cos x$ là
	 \choice
		{$y'=2\cos x-\sin x$} 
		{$y'=-2\cos  x-\sin x$}
		{$y'=2\cos x+\sin  x$}
		{\True  $y'=2\cos x+\sin x$}
	\loigiai{ 
		 Ta có: $y'=\left(2\sin x\right)'-\left(\cos x\right)'=2\cos x+\sin x$.
	} 
\end{ex}

% nguồn: Ok_TracNghiemLop11_Lan1
\begin{ex}%[1D7N2-1]%[ma:e0c02356]
	Tính đạo hàm của  hàm số $y=x^2-2x+3$.
	 \choice
		{$y'=2x $} 
		{$y'=2x+2 $} 
		{$y'=x^2-2 $} 
		 {\True $y'=2x-2 $}
	\loigiai{Ta có  $y'=(x^2-2x+3)'=2x-2$.
	 }
\end{ex}

% nguồn: Ok_TracNghiemLop11_Lan1
\begin{ex}%[1D7N2-1]%[ma:25d4f541]
	Tính đạo hàm của hàm số $y=\tan 3x$. 
	\choice 
		{\True $y'=\dfrac{3}{\cos^{2}3x}$} 
		{$y'=\dfrac{-3}{\cos^{2}3x}$} 
		{$y'=\dfrac{-3}{\sin^{2}3x}$} 
		 {$y'=\dfrac{1}{\cos^{2}3x}$}
	 \loigiai{
		Ta có  $y'=\dfrac{3}{\cos^{2}3x}$.
	} 
\end{ex}

% nguồn: Ok_TracNghiemLop11_Lan1
\begin{ex}%[1D7N2-1]%[ma:c1541a5f]
	Đạo hàm của hàm số nào  dưới đây bằng $0$?
	\choice 
		{$y=\cos x$} 
		 {$y=x$}
		{\True $y=\cos 1^{\circ}$} 
		 {$y=\sqrt{x}$}
	\loigiai{ 
		Nhận xét, $y=\cos 1^{\circ}$ là  một hàm hằng nên suy ra $y'=0$.
	 }
\end{ex}

% nguồn: Ok_TracNghiemLop11_Lan1
\begin{ex}%[1D7N2-1]%[ma:51d1f0bb]
	Đạo hàm của hàm số $y=\dfrac{1}{x}$ tại  điểm $x=-2$ là
	 \choice
		 {\True $y'\left(-2\right)=-\dfrac{1}{4}$}
		 {$y'\left(-2\right)=\dfrac{1}{4}$}
		 {$y'\left(-2\right)=-\dfrac{1}{2}$}
		 {$y'\left(-2\right)=\dfrac{1}{2}$}
	\loigiai{Vì $  y'=\left(\dfrac{1}{x}\right)'=-\dfrac{1}{x^2} $ nên $ y'(-2)=-\dfrac{1}{4} $.}
\end{ex}

% nguồn: Ok_TracNghiemLop11_Lan1
\begin{ex}%[1D7N2-1]%[ma:4f79e9b4]
	Mệnh đề nào sau đây {\bf đúng}? 
	 \choice
		{$\left(\sin^2{2x}\right)^\prime  = \sin4x$}
		{$\left(\cos{2x}\right)^\prime  = 2\sin2x$}
		{\True  $\left(\cos^2{2x}\right)^\prime = -2\sin4x$}
		{$\left(\sin{2x}\right)^\prime  =-2 \cos2x$}
	\loigiai{ 
		Ta  có
		 $$\left(\sin^2{2x}\right)^\prime=2\sin 2x\cdot (\sin2x)^\prime=2\sin 2x\cdot 2\cos 2x=2\sin 4x.$$
		$$\left(\cos{2x}\right)^\prime=-\sin  2x\cdot (2x)^\prime=-2\sin 2x.$$
		$$\left(\cos^2{2x}\right)^\prime = 2\cos2x\cdot (\cos2x)^\prime =-2\cdot 2\cos2x  \cdot \sin2x=-2\sin4x.$$
		$$\left(\sin{2x}\right)^\prime=\cos 2x\cdot (2x)^\prime=2\cos 2x.$$ 
	} 
\end{ex}

% nguồn: Ok_TracNghiemLop11_Lan1
\begin{ex}%[1D7N2-1]%[ma:67d77ba7]
	Khẳng định nào  sau đây là \textbf{sai}?
	\choice 
		{$(\sin x)'=\cos  x$}
		 {$(\cot x)'=\dfrac{-1}{\sin^2x}$}
		{\True  $(\tan x)'=\dfrac{1}{\sin^2x}$}
		 {$(\cos x)'=-\sin x$}
	 \loigiai{
		Theo công thức tính đạo hàm của hàm số lượng giác,  khẳng định $(\tan x)'=\dfrac{1}{\sin^2x}$ là sai. Công thức đúng là $(\tan x)'=\dfrac{1}{\cos^2x}$.
	} 
\end{ex}

% nguồn: Ok_TracNghiemLop11_Lan1
\begin{ex}%[1D7N2-1]%[ma:8d42a675]
	Đạo hàm  của hàm số $y=\tan x$ là
	\choice 
		 {$\dfrac{1}{\sin^2 x}$}
		 {$-\dfrac{1}{\sin^2 x}$}
		{\True  $\dfrac{1}{\cos^2 x}$}
		{$-\dfrac{1}{\cos^2 x}$} 
	 \loigiai{
		Ta có $y'=\dfrac{1}{\cos^2 x}$. 
	} 
\end{ex}

% nguồn: Ok_TracNghiemLop11_Lan1
\begin{ex}%[1D7N2-1]%[ma:f298315f]
	Hàm số $y=5^{x^2+1}$ có đạo hàm  là?
	 \choice
		{$5^{x^2+1}$} 
		{$\left(x^2+1\right) 5^{x^2+1}\ln 5$} 
		 {$\left(x^2+1\right) 5^{x^2}$}
		{\True $2x  5^{x^2+1}\ln 5$}
	\loigiai{$y'=(x^2+1)'\cdot 5^{x^2+1}\cdot \ln5=2x\cdot 5^{x^2+1}\cdot  \ln5$.
	} 
\end{ex}

% nguồn: Ok_TracNghiemLop11_Lan1
\begin{ex}%[1D7N2-1]%[ma:d687b250]
	Đạo hàm của hàm số $y=2\sin^2 x-\cos 2x+x$ là 
	\choice 
		{\True $y'=4\sin  2x+1$}
		{$y'=4\sin  x+\sin 2x+1$}
		{$y'=4\sin x -2\sin 2x+1$} 
		{$y'=1$} 
	\loigiai{ 
		Ta có $y' = \left(2\sin^2 x-\cos 2x+x\right)' = 4\sin x \cdot  \cos x+ 2\sin 2x+1 = 2\sin 2x + 2\sin 2x +1 = 4\sin 2x+1$.
	 }
\end{ex}

% nguồn: Ok_TracNghiemLop11_Lan1
\begin{ex}%[1D7N2-1]%[ma:2946d285]
	Hàm số  $y=x^4$ có đạo hàm trên $\left(-\infty;+\infty \right)$ là
	\choice 
		{\True  $y'=4x^3$}
		{$y'=3x^3$} 
		 {$y'=4x^4$}
		 {$y'=3x^4$}
	\loigiai{Ta có  $ y'=(x^4)'=4\cdot x^3 $.}
\end{ex}

% nguồn: Ok_Toan11_TracNghiem_Lan2
\begin{ex}%[1D7H2-1]%[ma:c015d55b]
	Cho hàm số  $y=\dfrac{4}{x-1}$. Khi đó $y'(-1)$ bằng
	\choice 
		 {\True $-1$}
		{$1$} 
		 {$-2$}
		{$2$} 
	\loigiai{ 
		Ta có $y'=\dfrac{-4}{(x-1)^2}\Rightarrow  y'(-1)=-1$.
	} 
\end{ex}

% nguồn: Ok_Toan11_TracNghiem_Lan2
\begin{ex}%[1D7H2-1]%[ma:2be8e112]
	Đạo  hàm của hàm số $y=\left(x^3-2x^2\right)^2$ bằng
	\choice 
		 {$6x^5-20x^4-16x^3$}
		{$6x^5-20x^4+4x^3$} 
		 {$6x^3+16x^3$}
		{\True $6x^5-20x^4+16x^3$} 
	\loigiai{ 
		Ta có $y'=2\left(x^3-2x^2\right)\cdot  \left(x^3-2x^2\right)'=2\left(x^3-2x^2\right)\left(3x^2-4x\right)=6x^5-20x^4+16x^3$.
	 }
\end{ex}

% nguồn: Ok_Toan11_TracNghiem_Lan2
\begin{ex}%[1D7H2-1]%[ma:4c6aca35]
	Tính đạo hàm của hàm số $f(x)=\dfrac{2x}{x-1}$ tại điểm  $x=-1$.
	 \choice
		 {$f'(-1)=1$}
		 {\True $f'(-1)=-\dfrac{1}{2}$}
		 {$f'(-1)=-2$}
		{$f'(-1)=0$} 
	\loigiai{Ta có $ f'(x)=\dfrac{-2}{(x-1)^2}$.\\ 
		Vậy $f'(-1)=-\dfrac{1}{2}$. 
	} 
\end{ex}

% nguồn: Ok_Toan11_TracNghiem_Lan2
\begin{ex}%[1D7H2-1]%[ma:d6fd40a0]
	 Hàm số $f(x)=5^{2x^{2}-1}$ có đạo hàm là
	 \choice
		{$2x  \cdot 5^{2x^{2}-1} \cdot ln~5$}
		{$4x \cdot 5^{2x^{2}-1}$} 
		{\True $4x \cdot  5^{2x^{2}-1} \cdot \ln~5$}
		 {$5^{2x^{2}-1}$}
	 \loigiai{
		Ta có $(5^{2x^{2}-1})'=(2x^{2}-1)' \cdot  5^{2x^{2}-1} \cdot \ln~5=4x \cdot 5^{2x^{2}-1} \cdot \ln~5$.
	} 
\end{ex}

% nguồn: Ok_Toan11_TracNghiem_Lan2
\begin{ex}%[1D7H2-1]%[ma:0466d2e9]
	Tính đạo hàm của hàm số $y=\log_2  (2x+1)$.
	 \choice
		{\True  $y'=\dfrac{2}{(2x+1)\ln 2}$}
		{$y'=\dfrac{1}{(2x+1)\ln  2}$}
		{$y'=\dfrac{2}{2x+1}$} 
		{$y'=\dfrac{1}{2x+1}$} 
	 \loigiai{
		Ta  có $y'=\left(\log_2 (2x+1)\right)'=\dfrac{(2x+1)'}{(2x+1)\ln 2}=\dfrac{2}{(2x+1)\ln 2}$.
	 }
\end{ex}

% nguồn: Ok_Toan11_TracNghiem_Lan2
\begin{ex}%[1D7H2-1]%[ma:08771446]
	Đạo hàm của hàm  số $y=10^{x^{2}}$ là
	 \choice
		 {$10^{x^{2}} \cdot \ln~10$}
		 {$2x \cdot 10^{x^{2}}$}
		{\True $2x \cdot  10^{x^{2}} \cdot \ln~10$}
		{$\dfrac{1}{2}+\log_{2}a$} 
	\loigiai{ 
		Ta có  $y'=(x^{2})' \cdot 10^{x^{2}} \cdot \ln~10=2x \cdot 10^{x^{2}} \cdot \ln~10$.
	} 
\end{ex}

% nguồn: Ok_Toan11_TracNghiem_Lan2
\begin{ex}%[1D7H2-1]%[ma:4c2ef29c]
	Hàm số  $y=2^{2x}$ có đạo hàm là
	\choice 
		{$y'=2^{2x}\ln~2$} 
		{\True  $y'=2^{2x+1}\ln~2$}
		{$y'=2^{2x-1}$} 
		{$y'=2x \cdot  2^{2x-1}$}
	 \loigiai{
		Ta có $y'=(2x)' \cdot 2^{2x}\ln~2=2 \cdot 2^{2x}\ln~2  = 2^{2x+1}\ln~2$.
	 }
\end{ex}

% nguồn: Ok_Toan11_TracNghiem_Lan2
\begin{ex}%[1D7H2-1]%[ma:8ba7b574]
	 Cho hàm số $y=x^2-x+2$. Tính $y'(1)$.
	 \choice
		 {$y'(1)=-1$}
		{\True  $y'(1)=1$}
		{$y'(1)=2$} 
		 {$y'(1)=0$}
	\loigiai{ 
		 Ta có $y'=2x-1\Rightarrow y'(1)=2\cdot 1-1=1$.
	} 
\end{ex}

% nguồn: Ok_Toan11_TracNghiem_Lan2
\begin{ex}%[1D7H2-1]%[ma:e3c3d753]
	Đạo hàm của hàm số  $y=x^4-4mx^2-3m-1$ ($m$ là tham số) là
	\choice 
		 {\True $y'=4x^3-8mx$}
		{$y'=4x^3-8mx-3m-1$} 
		{$y'=4x^3-8mx-1$} 
		{$y'=4x^2-8mx$} 
	\loigiai{Ta  có $y'=4x^3-8mx$.
	} 
\end{ex}

% nguồn: Ok_Toan11_TracNghiem_Lan2
\begin{ex}%[1D7H2-1]%[ma:48476bc4]
	Đạo hàm của hàm số $y=\dfrac{x^4}{2}+\dfrac{5x^3}{3}-\sqrt{2x}+a^2$ ($a$ là  hằng số) bằng
	\choice 
		{$2x^3+5x^2-\dfrac{1}{\sqrt{2x}}+2a$} 
		{$2x^3+5x^2+\dfrac{1}{2\sqrt{2x}}$} 
		 {\True $2x^3+5x^2-\dfrac{1}{\sqrt{2x}}$}
		 {$2x^3+5x^2-\sqrt{2}$}
	\loigiai{Ta có  $y'=2x^3+5x^2-\dfrac{1}{\sqrt{2x}}$.
	 }
\end{ex}

% nguồn: Ok_Toan11_TracNghiem_Lan2
\begin{ex}%[1D7H2-1]%[ma:86ae8993]
	Cho hàm số $f(x)=ax^3+\dfrac{b}{x}$  có $f'(1)=1$, $f'(-2)=-2$. Khi đó, $f'\left(\sqrt{2}\right)$ bằng
	\choice 
		{$\dfrac{12}{5}$} 
		{\True $\dfrac{-2}{5}$} 
		 {$2$}
		{$-\dfrac{12}{5}$} 
	\loigiai{ 
		Ta có  $f'(x)=3ax^2-\dfrac{b}{x^2}$, $\forall x\ne0$.\\
		Suy ra $\heva{& f'(1)=3a-b\\ &f'(-2)=12a-\dfrac{b}{4}.}$\\ 
		Theo giả thiết, ta  có
		\begin{align*}
			\heva{& f'(1)=1\\ & f'(-2)=-2}\Leftrightarrow\heva{&3a-b=1\\ &  12a-\dfrac{b}{4}=-2}\Leftrightarrow\heva{&a=-\dfrac{1}{5}\\ &b=-\dfrac{8}{5}.}
		\end{align*}
		Do đó, $f'(x)=-\dfrac{3}{5}x^2+\dfrac{8}{5x^2}$.\\ 
		Vậy  $f'\left(\sqrt{2}\right)=-\dfrac{3}{5}\cdot 2+\dfrac{8}{5\cdot 2}=-\dfrac{2}{5}$.
	 }
\end{ex}

% nguồn: Ok_Toan11_TracNghiem_Lan2
\begin{ex}%[1D7H2-1]%[ma:82f19537]
	Tính đạo hàm của  hàm số $y=\dfrac{2x^2-3x+7}{x^2+2x+3}$.
	 \choice
		{$y'=\dfrac{-7x^2+2x+23}{(x^2+2x+3)^2}$} 
		 {\True $y'=\dfrac{7x^2-2x-23}{(x^2+2x+3)^2}$}
		{$y'=\dfrac{7x^2-2x-23}{x^2+2x+3}$} 
		{$y'=\dfrac{8x^3+3x^2+14x+5}{(x^2+2x+3)^2}$} 
	 \loigiai{
		 Ta có $y=\dfrac{2x^2-3x+7}{x^2+2x+3}$ nên
		\begin{align*}
			y'=\dfrac{(4x-3)(x^2+2x+3)-(2x+2)(2x^2-3x+7)}{(x^2+2x+3)^2}=\dfrac{7x^2-2x-23}{(x^2+2x+3)^2}. 
		\end{align*}
	} 
\end{ex}

% nguồn: Ok_Toan11_TracNghiem_Lan2
\begin{ex}%[1D7H2-1]%[ma:7a8410aa]
	Đạo hàm của hàm số  $y=3^{x^{3}+2}$ là
	 \choice
		{$y=3x^{2} \cdot 3^{x^{3}+2}$} 
		 {\True $y=x^{2} \cdot 3^{x^{3}+3} \cdot \ln~3$}
		{$y=3x^{2} \cdot  (x^{3}+2) \cdot 3^{x^{3}+1}$}
		{$y=3^{x^{3}+2}  \cdot \ln~3$}
	\loigiai{ 
		Ta có $y'=(x^{3}+2)'  \cdot 3^{x^{3}+2} \cdot \ln~3=3x^{2} \cdot 3^{x^{3}+2} \cdot \ln~3=x^{2} \cdot 3^{x^{3}+3} \cdot \ln~3$.
	} 
\end{ex}

% nguồn: Ok_Toan11_TracNghiem_Lan2
\begin{ex}%[1D7H2-1]%[ma:c9d32967]
	Đạo hàm của hàm số $y=\sqrt{1-2x^2}$ là kết quả nào sau đây? 
	\choice 
		 {$y'=\dfrac{1}{2\sqrt{1-2x^2}}$}
		 {$y'=\dfrac{-4x}{2\sqrt{1-2x^2}}$}
		{\True $y'=\dfrac{-2x}{\sqrt{1-2x^2}}$} 
		{$y'=\dfrac{2x}{\sqrt{1-2x^2}}$} 
	\loigiai{Ta có $y'=\dfrac{-4x}{2\sqrt{1-2x^2}}=\dfrac{-2x}{\sqrt{1-2x^2}}$. 
	} 
\end{ex}

% nguồn: Ok_Toan11_TracNghiem_Lan2
\begin{ex}%[1D7H2-1]%[ma:e713291a]
	Cho hàm số $y=x^3-3x+2025$. Bất phương trình  $y'<0$ có tập nghiệm là
	\choice 
		{\True $S=(-1;1)$} 
		 {$S=(-\infty;-1)\cup(1;+\infty)$}
		{$S=(1;+\infty)$} 
		{$S=(-\infty;-1)$} 
	 \loigiai{
		Ta có $y=x^3-3x+2015$ nên $y'=3x^2-3$.  Do đó,
		\begin{align*}
			y'<0\Leftrightarrow 3x^2-3<0\Leftrightarrow-1<x<1. 
		\end{align*}
		Vậy bất phương  trình $y'<0$ có tập nghiệm là $S=(-1;1)$.
	 }
\end{ex}

% nguồn: Ok_Toan11_TracNghiem_Lan2
\begin{ex}%[1D7H2-1]%[ma:b21341e2]
	Tính đạo hàm của hàm  số $y=\dfrac{2x}{x-1}$.
	\choice 
		{$y'=\dfrac{2}{(x-1)^2}$} 
		 {$y'=\dfrac{2}{x-1}$}
		{\True $y'=\dfrac{-2}{(x-1)^2}$} 
		 {$y'=\dfrac{-2}{x-1}$}
	 \loigiai{Ta có $y'=\dfrac{-2}{(x-1)^2}$.
	 }
\end{ex}

% nguồn: Ok_Toan11_TracNghiem_Lan2
\begin{ex}%[1D7H2-1]%[ma:e28fc0ea]
	Tính đạo hàm của hàm số  $f(x)=e^{2x-3}$.
	 \choice
		{\True  $f'(x)=2 \cdot e^{2x-3}$}
		 {$f'(x)=e^{2x-3}$}
		 {$f'(x)=-2 \cdot e^{2x-3}$}
		{$f'(x)=2 \cdot  e^{x-3}$}
	\loigiai{ 
		Ta có $f'(x)=(2x-3)'e^{2x-3}=2 \cdot  e^{2x-3}$.
	 }
\end{ex}

% nguồn: Ok_Toan11_TracNghiem_Lan2
\begin{ex}%[1D7H2-1]%[ma:2e89290e]
	Tính đạo hàm của  hàm số $y=2^{2x+3}$.
	\choice 
		{$y'=2^{2x+2}\ln~4$} 
		 {$y'=4^{x+2}\ln~4$}
		{\True $y'=2^{2x+2}\ln~16$} 
		 {$y'=2^{2x+3}\ln~2$}
	 \loigiai{
		Ta có $y'=(2x+3)' \cdot 2^{2x+3}\ln~2=2 \cdot 2^{2x+3} \ln 2 = 2^{2x+4} \ln 2 = 2^{2x+2} \cdot 2^2 \ln 2 =  2^{2x+2} \ln 16$.
	} 
\end{ex}

% nguồn: Ok_Toan11_TracNghiem_Lan2
\begin{ex}%[1D7H2-1]%[ma:95f44dbb]
	Đạo hàm của  hàm số $y=2\, 022^{2-x}$ là
	\choice 
		{$y'=\dfrac{2\, 022^{2-x}}{\ln~2\, 022}$} 
		{\True $y'=-2\, 022^{2-x}\ln~2\, 022$} 
		{$y'=2\, 022^{2-x}\ln~2\,  022$}
		{$y'=x  \cdot 2\, 022^{2-x}$}
	\loigiai{ 
		Ta có $y' = (2-x)' \cdot 2\, 022^{2-x} \cdot \ln 2\, 022 = -2\,  022^{2-x} \cdot \ln 2\, 022$.
	 }
\end{ex}

% nguồn: Ok_Toan11_TracNghiem_Lan2
\begin{ex}%[1D7H2-1]%[ma:94787474]
	Đạo hàm của hàm số  $y=\log_2 (2x+1)$ là
	\choice 
		{$y'=\dfrac{1}{(2x+1)\ln  2}$}
		 {$y'=\dfrac{1}{2x+1}$}
		{$\dfrac{2}{2x+1}$} 
		{\True $\dfrac{2}{(2x+1)\ln 2}$} 
	\loigiai{ 
		Ta có $\left(\log_2  (2x+1)\right)'=\dfrac{(2x+1)'}{(2x+1)\ln 2}=\dfrac{2}{(2x+1)\ln 2}$.
	} 
\end{ex}

% nguồn: Ok_Toan11_TracNghiem_Lan2
\begin{ex}%[1D7H2-1]%[ma:4f54547a]
	Tính đạo  hàm của hàm số $y=\dfrac{x+1}{1-3x}$.
	 \choice
		{\True  $y'=\dfrac{4}{(1-3x)^2}$}
		{$y'=\dfrac{4x}{(1-3x)^2}$} 
		{$y'=\dfrac{-4}{(1-3x)^2}$} 
		 {$y'=\dfrac{4}{1-3x}$}
	\loigiai{Ta  có $y'=\left(\dfrac{x+1}{1-3x}\right)'=\dfrac{(x+1)'\cdot (1-3x)-(x+1)\cdot (1-3x)'}{(1-3x)^2}=\dfrac{(1-3x)+3(x+1)}{(1-3x)^2}=\dfrac{4}{(1-3x)^2}$.
	} 
\end{ex}

% nguồn: Ok_Toan11_TracNghiem_Lan2
\begin{ex}%[1D7H2-1]%[ma:7e81144a]
	Cho hàm số $y=x^3+mx^2+3x-5$ với $m$ là tham số. Tìm tập hợp $M$  tất cả các giá trị của $m$ để $y'=0$ có hai nghiệm phân biệt.
	 \choice
		{$M=(-3;3)$} 
		{$M=(-\infty;-3]\cup[3;+\infty)$} 
		 {$M=\mathbb{R}$}
		{\True  $M=(-\infty;-3)\cup(3;+\infty)$}
	 \loigiai{
		Ta có $y=x^3+mx^2+3x-5$ nên $y'=3x^2+2mx+3$.  Do đó,\\
		 $y'=0$ có hai nghiệm phân biệt khi và chỉ khi
		\begin{align*}
			\Delta'>0\Leftrightarrow m^2-9>0\Leftrightarrow \hoac{&m<-3\\  &m>3.}
		\end{align*}
		Vậy  $M=(-\infty;-3)\cup(3;+\infty)$.
	 }
\end{ex}

% nguồn: Ok_Toan11_TracNghiem_Lan2
\begin{ex}%[1D7H2-1]%[ma:49e46e6f]
	Cho hàm số $f(x)=x^4+2x^2-3$. Tìm $x$  để $f'(x)>0$.
	 \choice
		{$-1<x<0$} 
		{$x<0$} 
		{\True  $x>0$}
		 {$x<-1$}
	 \loigiai{
		Ta có $f'(x)=4x^3+4x$.  Do đó,
		\begin{align*}
			f'(x)>0\Leftrightarrow 4x^3+4x>0\Leftrightarrow 4x(x^2+1)>0\Leftrightarrow x>0. 
		\end{align*}
	 }
\end{ex}

% nguồn: Ok_Toan11_TracNghiem_Lan2
\begin{ex}%[1D7H2-1]%[ma:a48c8823]
	Tính đạo hàm của hàm số  $y=f(x)=\left(x^3-2x^2\right)^2$.
	\choice 
		 {\True $f'(x)=6x^5-20x^4+16x^3$}
		{$f'(x)=6x^5+16x^3$} 
		{$f'(x)=6x^5-20x^4+4x^3$} 
		 {$f'(x)=6x^5-20x^4-16x^3$}
	\loigiai{Ta có \begin{align*}
		f'(x)&=2\left(x^3-2x^2\right)\cdot\left(x^3-2x^2\right)'\\ 
		&=2\left(x^3-2x^2\right)\cdot\left(3x^2-4x\right)\\ 
		 &=6x^5-8x^4-12x^4+16x^3\\
		&=6x^5-20x^4+16x^3. 
	\end{align*}
} 
\end{ex}

% nguồn: Ok_Toan11_TracNghiem_Lan2
\begin{ex}%[1D7H2-1]%[ma:badd73c9]
	Trên tập số thực, đạo hàm của hàm số $y=3^{x^{2}-x}$  là
	 \choice
		{$y'=(2x-1) \cdot 3^{x^{2}-x}$} 
		{\True $y'=(2x-1) \cdot 3^{x^{2}-x}  \cdot \ln~3$}
		{$y'=(x^{2}-x)  \cdot 3^{x^{2}+x+1}$}
		 {$y'=3^{x^{2}-x-1}$}
	 \loigiai{
		Ta có $y=3^{x^{2}-x} \Rightarrow y'=(x^{2}-x)' \cdot  3^{x^{2}-x} \cdot \ln~3=(2x-1) \cdot 3^{x^{2}-x} \cdot \ln~3$.
	 }
\end{ex}

% nguồn: Ok_Toan11_TracNghiem_Lan2
\begin{ex}%[1D7H2-1]%[ma:97cc062b]
	Cho  hàm số $y=\ln (\cos x+1+m^2)$. Với giá trị nào của $m$ thì $y'\left(\dfrac{\pi}{2} \right)=\dfrac{-1}{5}$.
	\choice 
		{\True $m=\pm  2$}
		 {$m=2$}
		 {$m=\dfrac{1}{2}$}
		{$m=\pm  1$}
	 \loigiai{
		Ta có $y'=\dfrac{-\sin x}{\cos  x+1+m^2} \Rightarrow y'\left(\dfrac{\pi}{2} \right)=\dfrac{-1}{1+m^2}$.\\
		Mà $y'(0)=\dfrac{-1}{5} \Rightarrow \dfrac{-1}{1+m^2}=\dfrac{-1}{5} \Rightarrow m=\pm 2$. 
	 }
\end{ex}

% nguồn: Ok_Toan11_TracNghiem_Lan2
\begin{ex}%[1D7H2-1]%[ma:0c84d8bc]
	Cho hàm số $f(x)=\sqrt{x-1}$. Đạo hàm của hàm số tại $x=1$  là
	\choice 
		 {$\dfrac{1}{2}$}
		 {$1$}
		 {$0$}
		{\True Không tồn tại} 
	\loigiai{Ta  có $f'(x)=\dfrac{1}{2\sqrt{x-1}}$ với $x>1$.\\
		Vậy không tồn tại đạo hàm của hàm  số tại $x=1$.
	 }
\end{ex}

% nguồn: Ok_Toan11_TracNghiem_Lan2
\begin{ex}%[1D7H2-1]%[ma:0e7aca56]
	Tính đạo hàm của  hàm số $y=\left(x^2-2\right)(2x-1)$.
	 \choice
		 {$y'=4x$}
		{$y'=3x^2-6x+2$} 
		 {\True $y'=6x^2-2x-4$}
		 {$y'=2x^2-2x+4$}
	\loigiai{Ta có  $y'=\left(x^2-2\right)'(2x-1)+\left(x^2-2\right)(2x-1)'=2x(2x-1)+2\left(x^2-2\right)=6x^2-2x-4$.
	} 
\end{ex}

% nguồn: Ok_Toan11_TracNghiem_Lan2
\begin{ex}%[1D7H2-1]%[ma:64afd215]
	Hàm số $y=3^{x^{2}-x}$ có đạo  hàm là
	 \choice
		{$(2x-1) \cdot  3^{x^{2}-x}$}
		{$(x^{2}-x) \cdot 3^{x^{2}-x-1}$} 
		{\True $(2x-1)  \cdot 3^{x^{2}-x} \cdot \ln~3$}
		{$3^{x^{2}-x}  \cdot \ln~3$}
	 \loigiai{
		Ta  có $(3^{x^{2}-x})'=(x^2-x)' \cdot 3^{x^{2}-x} \cdot \ln~3=(2x-1) \cdot 3^{x^{2}-x} \cdot \ln~3$.
	 }
\end{ex}

% nguồn: Ok_Toan11_TracNghiem_Lan2
\begin{ex}%[1D7H2-1]%[ma:dbbe86f4]
	 Tính đạo hàm của hàm số $y=\dfrac{x^2+2x-3}{x+2}$.
	\choice 
		 {\True $y'=1+\dfrac{3}{(x+2)^2}$}
		 {$y'=\dfrac{x^2+6x+7}{(x+2)^2}$}
		 {$y'=\dfrac{x^2+4x+5}{(x+2)^2}$}
		 {$y'=\dfrac{x^2+8x+1}{(x+2)^2}$}
	\loigiai{Ta  có $y=x-\dfrac{3}{x+2}\Rightarrow y'=1+\dfrac{3}{(x+2)^2}$.}
\end{ex}

% nguồn: Ok_Toan11_TracNghiem_Lan2
\begin{ex}%[1D7H2-1]%[ma:0f69dc1d]
	Cho hàm số $y=2^{x^{2}-mx+1}$. Với  giá trị nào của tham số $m$ thì $y'(0)=\ln~2$?
	\choice 
		{\True  $-\dfrac{1}{2}$}
		 {$\dfrac{1}{2}$}
		 {$2$}
		{$-2$} 
	\loigiai{ 
		\begin{itemize}
			\item Ta có $y'=(x^{2}-mx+1)' \cdot 2^{x^{2}-mx+1} \cdot \ln~2=(2x-m)  \cdot 2^{x^{2}-mx+1} \cdot \ln~2$.
			\item Khi đó $y'(0)=(-m) \cdot  2^{1} \cdot \ln~2 = -2m \ln~2$.
			\item Theo đề  bài $y'(0)=\ln~2 \Leftrightarrow -2m \ln~2 = \ln~2 \Leftrightarrow m=-\dfrac{1}{2}$.
		\end{itemize}
	} 
\end{ex}

% nguồn: Ok_Toan11_TracNghiem_Lan2
\begin{ex}%[1D7H2-1]%[ma:50ccdf34]
	Cho hàm số $f(x)=\sqrt{x^2+3}$. Tính giá trị của biểu thức $S=f(1)+4f'(1)$. 
	\choice 
		 {\True $S=4$}
		 {$S=2$}
		{$S=6$} 
		 {$S=8$}
	 \loigiai{Ta có $f'(x)=\dfrac{x}{\sqrt{x^2+3}}$.\\
		$f(1)=2$ và $f'(1)=\dfrac{1}{2}$.\\ 
		Vậy  $S=f(1)+4f'(1)=2+4\cdot\dfrac{1}{2}=4$.
	 }
\end{ex}

% nguồn: Ok_Toan11_TracNghiem_Lan2
\begin{ex}%[1D7H2-1]%[ma:11e2bd57]
	Đạo hàm của  hàm số $y=(1-2x)^3$ là
	 \choice
		 {$y'=3\cdot (1-2x)^2$}
		{$y'=6\cdot (1-2x)^2$} 
		{$y'=-3\cdot  (1-2x)^2$}
		{\True  $y'=-6\cdot (1-2x)^2$}
	\loigiai{Áp dụng công thức $\left(u^n\right)'=n\cdot u^{n-1}\cdot  u'$, với $u=1-2x$, $n=3$.\\
		Suy ra $y'=3\cdot (1-2x)^2\cdot  (1-2x)'=-6\cdot (1-2x)^2$.
	} 
\end{ex}

% nguồn: Ok_Toan11_TracNghiem_Lan2
\begin{ex}%[1D7H2-1]%[ma:8c3b937d]
	Hàm số $y=\ln \left(2x+1\right)$ có  đạo hàm là
	\choice 
		{$y'=\dfrac{2}{x\ln \left(2x+1\right)}$} 
		 {$y'=\dfrac{1}{2x+1}$}
		 {\True $y'=\dfrac{2}{2x+1}$}
		{$y'=\dfrac{1}{\left(2x+1\right)\ln 2}$} 
	\loigiai{ 
		Hàm số $y=\ln \left(2x+1\right)$ có đạo hàm  là $y'=\dfrac{2}{2x+1}$.
	} 
\end{ex}

% nguồn: Ok_Toan11_TracNghiem_Lan2
\begin{ex}%[1D7H2-1]%[ma:8bd87557]
	Cho  $f(x)=x^5+x^3-2x-3$. Tính $f'(1)+f'(-1)+4f'(0)$.
	\choice 
		{\True  $4$}
		{$7$} 
		{$6$} 
		 {$5$}
	 \loigiai{
		 Ta có $f'(x)=5x^4+3x^2-2$.\\
		 Suy ra $f'(1)=6$; $f'(-1)=6$; $f '(0)=-2$.\\
		Vậy $f'(1)+f'(-1)+4f'(0)=4$. 
	 }
\end{ex}

% nguồn: Ok_Toan11_TracNghiem_Lan2
\begin{ex}%[1D7H2-1]%[ma:85dbd2a3]
	Cho $\left(\dfrac{3-2x}{\sqrt{4x-1}}\right)’=\dfrac{ax-b}{(4x-1)\sqrt{4x-1}}$, $\forall x>\dfrac{1}{4}$. Tính  $\dfrac{a}{b}$.
	\choice 
		{$-16$} 
		{$-4$} 
		{\True $-1$} 
		 {$4$}
	 \loigiai{
		 Với mọi $x>\dfrac{1}{4}$, ta có\allowdisplaybreaks
		\begin{align*}
			\left(\dfrac{3-2x}{\sqrt{4x-1}}\right)’&=\dfrac{(3-2x)’\sqrt{4x-1}-(3-2x)\left(\sqrt{4x-1}\right)’}{(4x-1)}\\ 
			&=\dfrac{-2\sqrt{4x-1}-\dfrac{6-4x}{\sqrt{4x-1}}}{(4x-1)}\\ 
			&=\dfrac{-4x-4}{(4x-1)\sqrt{4x-1}}. 
		\end{align*}
		Do  đó, $a=-4$, $b=4$. Suy ra $\dfrac{a}{b}=-1$.
	} 
\end{ex}

% nguồn: Ok_Toan11_TracNghiem_Lan2
\begin{ex}%[1D7H2-1]%[ma:d2dbca4e]
	Tính đạo  hàm của hàm số $y=x^2-\dfrac{1}{x}$.
	 \choice
		 {$y'=2x-\dfrac{1}{x^2}$}
		{$y'=x-\dfrac{1}{x^2}$} 
		 {$y'=x+\dfrac{1}{x^2}$}
		 {\True $y'=2x+\dfrac{1}{x^2}$}
	\loigiai{Tập xác định $\mathscr{D}=\mathbb{R}\setminus\{0\}$.\\ 
		Ta có $y'=2x+\dfrac{1}{x^2}$. 
	 }
\end{ex}

% nguồn: Ok_Toan11_TracNghiem_Lan2
\begin{ex}%[1D7H2-1]%[ma:f3e57d48]
	Đạo hàm của hàm  số $y=(2x-1)\sqrt{x^2+x}$ là
	\choice 
		{\True  $y'=\dfrac{8x^2+4x-1}{2\sqrt{x^2+x}}$}
		 {$y'=\dfrac{8x^2+4x+1}{2\sqrt{x^2+x}}$}
		 {$y'=\dfrac{4x+1}{2\sqrt{x^2+x}}$}
		{$y'=\dfrac{6x^2+2x-1}{2\sqrt{x^2+x}}$} 
	\loigiai{ 
		Ta có  $y'=2\sqrt{x^2+x}+\dfrac{(2x-1)(2x+1)}{2\sqrt{x^2+x}}=\dfrac{4x^2+4x+4x^2-1}{2\sqrt{x^2+x}}=\dfrac{8x^2+4x-1}{2\sqrt{x^2+x}}$.
	} 
\end{ex}

% nguồn: Ok_Toan11_TracNghiem_Lan2
\begin{ex}%[1D7H2-1]%[ma:d4fa0ff9]
	Đạo hàm cùa hàm  số $y=\log_4 (2x+5)$ là
	\choice 
		 {$y'=\dfrac{1}{(2x+5)\ln 4}$}
		{\True $y'=\dfrac{1}{(2x+5)\ln 2}$} 
		{$y'=\dfrac{2\ln  4}{(2x+5)}$}
		{$y'=\dfrac{2}{\left(2x+5\right)\ln 5}$} 
	 \loigiai{
		Ta có $y'=\dfrac{2}{(2x+5)\ln  4}=\dfrac{2}{(2x+5)\cdot 2\ln 2}=\dfrac{1}{(2x+5)\ln 2}$.
	 }
\end{ex}

% nguồn: Ok_Toan11_TracNghiem_Lan2
\begin{ex}%[1D7H2-1]%[ma:763b6f66]
	Đạo hàm của hàm số  $y=\ln \dfrac{x}{x+1}$
	 \choice
		{$-\dfrac{1}{x\left(x+1\right)}$} 
		 {$\dfrac{x}{x+1}$}
		{$\dfrac{x+1}{x}$} 
		{\True $\dfrac{1}{x\left(x+1\right)}$} 
	\loigiai{ 
		\[y'=\left(\ln \dfrac{x}{x+1} \right)'=\dfrac{\dfrac{1}{\left(x+1\right)^2}}{\dfrac{x}{x+1}}=\dfrac{1}{x\left(x+1\right)}.\] 
	 }
\end{ex}

% nguồn: Ok_Toan11_TracNghiem_Lan2
\begin{ex}%[1D7H2-1]%[ma:7ef2f199]
	Tính đạo  hàm của hàm số $y=\log_2 (3x)$
	\choice 
		 {$y'=\dfrac{1}{x\ln 4}$}
		{$y'=\dfrac{3}{x\ln  2}$}
		{\True $y'=\dfrac{1}{x\ln 2}$} 
		{$y'=\dfrac{3}{x\ln 4}$} 
	 \loigiai{
		Ta có  $y'=\left(\log_2 (3x)\right)'=\dfrac{(3x)'}{3x\ln 2}=\dfrac{1}{x\ln 2}$.
	} 
\end{ex}

% nguồn: Ok_Toan11_TracNghiem_Lan2
\begin{ex}%[1D7H2-1]%[ma:b4fa4dbd]
	Đạo hàm của hàm số  $y=\left(x^2-\dfrac{2}{x}\right)^3$ bằng
	 \choice
		{\True  $y'=6\left(x+\dfrac{1}{x^2}\right)\left(x^2-\dfrac{2}{x}\right)^2$}
		 {$y'=3\left(x^2-\dfrac{2}{x}\right)^2$}
		{$y'=6\left(x-\dfrac{1}{x^2}\right)\left(x^2-\dfrac{2}{x}\right)^2$} 
		{$y'=6\left(x-\dfrac{1}{x}\right)\left(x^2-\dfrac{2}{x}\right)^2$} 
	\loigiai{ 
		 Ta có $y'=3\left(x^2-\dfrac{2}{x}\right)^2\left(2x+\dfrac{2}{x^2}\right)=6\left(x+\dfrac{1}{x^2}\right)\left(x^2-\dfrac{2}{x}\right)^2$.
	 }
\end{ex}

% nguồn: Ok_Toan11_TracNghiem_Lan2
\begin{ex}%[1D7H2-1]%[ma:d875e8be]
	Cho $y=\sqrt{x^2-2x+3}$, $y'=\dfrac{ax+b}{\sqrt{x^2-2x+3}}$. Khi đó,  giá trị $a\cdot b$ bằng
	 \choice
		 {$-4$}
		{\True $-1$} 
		 {$0$}
		 {$1$}
	 \loigiai{
		Ta có $y=\sqrt{x^2-2x+3}$  nên $y'=\dfrac{(x^2-2x+3)’}{2\sqrt{x^2-2x+3}}=\dfrac{2x-2}{2\sqrt{x^2-2x+3}}=\dfrac{x-1}{\sqrt{x^2-2x+3}}$.\\
		Do đó, $a=1$, $b=-1$.\\ 
		Vậy $a\cdot b=-1$. 
	 }
\end{ex}

% nguồn: Ok_Toan11_TracNghiem_Lan2
\begin{ex}%[1D7H2-1]%[ma:8d7eb3fa]
	Cho hàm  số $y=\dfrac{2x+1}{x-2}$. Tính $y'(3)$.
	 \choice
		 {$y'(3)=5$}
		{\True  $y'(3)=-5$}
		 {$y'(3)=0$}
		 {$y'(3)=7$}
	\loigiai{ 
		Ta  có $y'=\dfrac{(2x+1)'(x-2)-(2x+1)(x-2)'}{(x-2)^2}=\dfrac{2(x-2)-(2x+1)\cdot 1}{(x-2)^2}=\dfrac{-5}{(x-2)^2}$.\\
		 Vậy $y'(3)=-5$.
	} 
\end{ex}

% nguồn: Ok_Toan11_TracNghiem_Lan2
\begin{ex}%[1D7H2-1]%[ma:854905ac]
	Đạo hàm  của hàm số $y=e^{3x}$ là
	 \choice
		{$y'=e^{3x}$} 
		{$y'=e^{3x} \cdot \ln~3$} 
		{\True  $y'=3e^{3x}$}
		{$y'=\dfrac{e^{3x}}{3}$} 
	 \loigiai{
		 Ta có $y'=(3x)' \cdot e^{3x}=3e^{3x}$.
	 }
\end{ex}

% nguồn: Ok_Toan11_TracNghiem_Lan2
\begin{ex}%[1D7H2-1]%[ma:aed1212b]
	Tính  đạo hàm của hàm số $y=2\, 021^{3-5x}$.
	\choice 
		{$y'=2\, 021^{3-5x}  \cdot \ln~2\, 021$}
		{\True  $y'=-5 \cdot 2\, 021^{3-5x} \cdot \ln~2\, 021$}
		{$y'=2\,  021^{3-5x}$}
		 {$y'=-5 \cdot 2\, 021^{3-5x} \cdot \log~2\, 021$}
	\loigiai{ 
		Ta có $y'=(3-5x)' \cdot 2\,  021^{3-5x} \cdot \ln~2\, 021=-5 \cdot 2\, 021^{3-5x} \cdot \ln~2\, 021$.
	} 
\end{ex}

% nguồn: Ok_Toan11_TracNghiem_Lan2
\begin{ex}%[1D7H2-1]%[ma:b0ab3549]
	Cho hàm số $f(x)=\ln (\cos x)$. Giá trị của $f'\left(-\dfrac{\pi}{4} \right)$ là 
	 \choice
		 {$0$}
		{\True  $-1$}
		 {$1$}
		{$\sqrt{2}$} 
	\loigiai{ 
		Ta có $f'(x)=\left[\ln (\cos x)\right]'=\dfrac{(\cos x)'}{\cos x}=\dfrac{-\sin x}{\cos  x}=-\tan x\Rightarrow f'\left(\dfrac{\pi}{4} \right)=-1$.
	} 
\end{ex}

% nguồn: Ok_Toan11_TracNghiem_Lan2
\begin{ex}%[1D7H2-1]%[ma:dc4227bf]
	 Đạo hàm của hàm số $y=\dfrac{1}{4}x^4-3x^2-2021x+2022$ bằng biểu thức nào sau đây?
	 \choice
		{$y'=x^4-6x^2-2021$} 
		{\True $y'=x^3-6x-2021$} 
		 {$y'=\dfrac{3}{4}x^3-6x-2021$}
		{$y'=x^3-6x-2021x+2022$} 
	\loigiai{ 
		Ta  có
		$$y'=\left(\dfrac{1}{4}x^4-3x^2-2021x+2022\right)'=\left(\dfrac{1}{4}x^4\right)'-\left(3x^2\right)'-(2021x)'+(2022)'=x^3-6x-2021.$$ 
	} 
\end{ex}

% nguồn: Ok_Toan11_TracNghiem_Lan2
\begin{ex}%[1D7H2-1]%[ma:f9921652]
	Cho hàm số $f(x)=2x^3+1$. Giá trị  $f'(-1)$ bằng
	 \choice
		{\True $6$} 
		{$3$} 
		 {$-2$}
		{$-6$} 
	\loigiai{Ta có $f'(x)=6x^2\Rightarrow f'(-1)=6\cdot  (-1)^2=6$.
	} 
\end{ex}

% nguồn: Ok_Toan11_TracNghiem_Lan2
\begin{ex}%[1D7H2-1]%[ma:dc4874f9]
	Tính đạo hàm của hàm số $y=\sqrt{x}+x$  tại điểm $x_0=4$ là
	 \choice
		 {$y'(4)=\dfrac{9}{2}$}
		 {$y'(4)=6$}
		{$y'(4)=\dfrac{3}{2}$} 
		{\True $y'(4)=\dfrac{5}{4}$} 
	\loigiai{Ta  có $y'=\dfrac{1}{2\sqrt{x}}+1$.\\
		Vậy $y'(4)=\dfrac{1}{2\sqrt{4}}+1=\dfrac{5}{4}$. 
	} 
\end{ex}

% nguồn: Ok_Toan11_TracNghiem_Lan2
\begin{ex}%[1D7H2-1]%[ma:c75834ec]
	Cho  hàm số $f(x)$ xác định trên $\mathbb{R}$ bởi $f(x)=2x^2+1$. Giá trị $f'(-1)$ bằng
	\choice 
		{$2$} 
		{$6$} 
		 {\True $-4$}
		 {$3$}
	 \loigiai{
		Ta có  $f'(x)=4x$.\\
		Vậy $f'(-1)=-4$. 
	 }
\end{ex}

% nguồn: Ok_Toan11_TracNghiem_Lan2
\begin{ex}%[1D7H2-1]%[ma:2d25412e]
	Cho hàm số $y=\dfrac{1}{3}x^3-2x^2-5x$. Tập nghiệm của  bất phương trình $y'\geq 0$ là
	\choice 
		 {$[-1;5]$}
		{$\varnothing$} 
		{$(-\infty;-1)\cup(5;+\infty)$} 
		{\True $(-\infty;-1]\cup[5;+\infty)$} 
	\loigiai{ 
		Ta có $y=\dfrac{1}{3}x^3-2x^2-5x$ nên $y'=x^2-4x-5$. Do  đó,
		\begin{align*}
			y'\geq 0\Leftrightarrow x^2-4x-5\geq  0\Leftrightarrow x\in(-\infty;-1]\cup[5;+\infty).
		\end{align*}
	} 
\end{ex}

% nguồn: Ok_Toan11_TracNghiem_Lan2
\begin{ex}%[1D7H2-1]%[ma:f1fe75f8]
	Tính đạo hàm của  hàm số $y=\log_2 \left(2x+1\right)$.
	\choice 
		{$y'=\dfrac{1}{\left(2x+1\right)\ln  2}$}
		{\True $y'=\dfrac{2}{\left(2x+1\right)\ln 2}$} 
		 {$y'=\dfrac{2}{2x+1}$}
		{$y'=\dfrac{1}{2x+1}$} 
	 \loigiai{
		Áp  dụng công thức tính đạo hàm $y'=\left(\log_2 \left(2x+1\right)\right)'=\dfrac{\left(2x+1\right)'}{\left(2x+1\right)\ln 2}=\dfrac{2}{\left(2x+1\right)\ln 2}$.
	} 
\end{ex}

% nguồn: Ok_Toan11_TracNghiem_Lan2
\begin{ex}%[1D7H2-1]%[ma:f0dbae1b]
	Hàm số  $y=2^{x^{2}-3x}$ có đạo hàm là
	 \choice
		{\True  $(2x-3) \cdot 2^{x^{2}-3x}\ln~2$}
		{$2^{x^{2}-3x}\ln~2$} 
		{$(2x-3) \cdot  2^{x^{2}-3x}$}
		{$(x^{2}-3x) \cdot 2^{x^{2}-3x+1}$} 
	 \loigiai{
		Ta  có $y'=(x^2-3x)' \cdot 2^{x^{2}-3x}\ln~2=(2x-3) \cdot 2^{x^{2}-3x}\ln~2$.
	 }
\end{ex}

% nguồn: Ok_Toan11_TracNghiem_Lan2
\begin{ex}%[1D7H2-1]%[ma:98d47b64]
	Đạo hàm  của hàm số $y=2^{2x}$ bằng.
	 \choice
		{$y'=2^{2x}\ln~2$} 
		{$y'=\dfrac{2^{2x+1}}{\ln~2}$} 
		 {\True $y'=4^{x}\ln~4$}
		{$y'=\dfrac{2^{2x}}{\ln~2}$} 
	 \loigiai{
		Ta có $(2^{2x})'=(2x)' \cdot 2^{2x} \cdot \ln~2=2 \cdot 2^{2x} \cdot \ln~2=2 \cdot 4^{x} \cdot \ln~2=4^{x}  \cdot \ln~4$.
	 }
\end{ex}

% nguồn: Ok_Toan11_TracNghiem_Lan2
\begin{ex}%[1D7H2-1]%[ma:5d49c101]
	Hàm số  $y=3^{x^{2}-3x}$ có đạo hàm là
	 \choice
		{$(2x-3) \cdot 3^{x^{2}-3x}$} 
		{$3^{x^{2}-3x} \cdot \ln~3$} 
		 {$(x^{2}-3x) \cdot 3^{x^{2}-3x-1}$}
		{\True $(2x-3) \cdot 3^{x^{2}-3x}  \cdot \ln~3$}
	 \loigiai{
		Ta có $y'=(x^2-3x)' \cdot 3^{x^{2}-3x} \cdot  \ln~3=(2x-3) \cdot 3^{x^{2}-3x} \cdot \ln~3$.
	} 
\end{ex}

% nguồn: Ok_Toan11_TracNghiem_Lan2
\begin{ex}%[1D7H2-1]%[ma:a9a34af9]
	Đạo hàm của hàm số $y=\left(x^2-2x+2\right)\mathrm{e}^x$  là
	 \choice
		 {$y'=\left(2x-2\right)\mathrm{e}^x$}
		{\True $y'=x^2\mathrm{e}^x$} 
		 {$y'=\left(x^2+2\right)\mathrm{e}^x$}
		 {$y'=-2x\mathrm{e}^x$}
	 \loigiai{
		 Ta có $y'=(2x-2)\mathrm{e}^x+\left(x^2-2x+2\right)\mathrm{e}^x=x^2\mathrm{e}^x$.
	 }
\end{ex}

% nguồn: Ok_Toan11_TracNghiem_Lan2
\begin{ex}%[1D7H2-1]%[ma:d0e3ced8]
	Hàm  số $f(x)=2^{3x+4}$ có đạo hàm là
	 \choice
		{$f'(x)=\dfrac{3  \cdot 2^{3x+4}}{\ln~2}$}
		{\True $f'(x)=3 \cdot 2^{3x+4}  \cdot \ln~2$}
		{$f'(x)=2^{3x+4} \cdot \ln~2$} 
		 {$f'(x)=\dfrac{2^{3x+4}}{\ln~2}$}
	 \loigiai{
		Ta có $f'(x)=(3x+4)' \cdot 2^{3x+4} \cdot \ln~2 = 3 \cdot 2^{3x+4} \cdot  \ln~2$.
	} 
\end{ex}

% nguồn: Ok_Toan11_TracNghiem_Lan2
\begin{ex}%[1D7H2-1]%[ma:c2afeebe]
	Hàm số $y=2^{x^{2}-x}$ có  đạo hàm là
	\choice 
		 {$2^{x^{2}-x} \cdot \ln~2$}
		 {\True $(2x-1) \cdot 2^{x^{2}-x} \cdot \ln~2$}
		{$(x^{2}-x)  \cdot 2^{x^{2}-x-1}$}
		 {$(2x-1) \cdot 2^{x^{2}-x}$}
	\loigiai{ 
		Ta có $y'=(x^{2}-x)' \cdot 2^{x^{2}-x}  \cdot ln~2=(2x-1) \cdot 2^{x^{2}-x} \cdot ln~2$.
	} 
\end{ex}

% nguồn: Ok_Toan11_TracNghiem_Lan2
\begin{ex}%[1D7H2-1]%[ma:70730a54]
	Tính đạo  hàm của hàm số $f(x)=\dfrac{2x+7}{x+4}$ tại $x=2$ ta được
	 \choice
		 {\True $f'(2)=\dfrac{1}{36}$}
		{$f'(2)=\dfrac{11}{6}$} 
		{$f'(2)=\dfrac{3}{2}$} 
		{$f'(2)=\dfrac{5}{12}$} 
	 \loigiai{
		Ta có  $f'(x)=\dfrac{1}{(x+4)^2}$.\\
		Vậy $f'(2)=\dfrac{1}{36}$. 
	} 
\end{ex}

% nguồn: Ok_Toan11_TracNghiem_Lan2
\begin{ex}%[1D7H2-1]%[ma:65473bae]
	Hàm số $f(x)=\log_2 \left(x^2-2x\right)$ có đạo  hàm
	 \choice
		 {$f'(x)=\dfrac{\ln 2}{x^2-2x}$}
		{$f'(x)=\dfrac{1}{\left(x^2-2x\right)\ln  2}$}
		{$f'(x)=\dfrac{\left(2x-2\right)\ln 2}{x^2-2x}$} 
		{\True $f'(x)=\dfrac{2x-2}{\left(x^2-2x\right)\ln  2}$}
	 \loigiai{
		Ta có $f'(x)=\dfrac{\left(x^2-2x\right)'}{\left(x^2-2x\right)\ln 2}=\dfrac{2x-2}{\left(x^2-2x\right)\ln  2}$.
	} 
\end{ex}

% nguồn: Ok_Toan11_TracNghiem_Lan2
\begin{ex}%[1D7H2-1]%[ma:c4e48ab0]
	 Hàm số nào sau đây có đạo hàm bằng $\dfrac{1}{\sqrt{2x}}$?
	\choice 
		{$f(x)=2\sqrt{x}$} 
		 {$f(x)=\sqrt{x}$}
		 {\True $f(x)=\sqrt{2x}$}
		{$f(x)=-\dfrac{1}{\sqrt{2x}}$} 
	 \loigiai{Ta có $\left(\sqrt{2x}\right)'=\dfrac{1}{\sqrt{2x}}$.\\
		Vậy hàm số $f(x)=\sqrt{2x}$ có đạo hàm bằng  $\dfrac{1}{\sqrt{2x}}$.
	} 
\end{ex}

% nguồn: Ok_Toan11_TracNghiem_Lan2
\begin{ex}%[1D7H2-1]%[ma:eacc39a7]
	Hàm số nào sau đây là đạo hàm của hàm  số $y=\dfrac{2x+3}{3x-1}$?
	 \choice
		 {\True $y'=\dfrac{-11}{(3x-1)^2}$}
		{$y'=\dfrac{7}{(3x-1)^2}$} 
		 {$y'=\dfrac{2}{(3x-1)^2}$}
		{$y'=\dfrac{2}{3}$} 
	 \loigiai{Ta có
		 $$
		y'=\left(\dfrac{2x+3}{3x-1}\right)'=\dfrac{(2x+3)'\cdot (3x-1)-(2x+3)\cdot (3x-1)'}{(3x-1)^2}=\dfrac{2(3x-1)-(2x+3)\cdot 3}{(3x-1)^2}=\dfrac{-11}{(3x-1)^2}. 
		$$ 
	 }
\end{ex}

% nguồn: Ok_Toan11_TracNghiem_Lan2
\begin{ex}%[1D7H2-1]%[ma:6642d6bb]
	Cho hàm  số $f(x)=\sin^22x$. Tính $f'(x)$.
	\choice 
		 {$f'(x)=2\sin2x$}
		 {$f'(x)=2\cos^22x$}
		{\True  $f'(x)=2\sin4x$}
		 {$f'(x)=-2\sin4x$}
	 \loigiai{Ta có $f'(x)=2\cdot\sin2x\cdot (\sin2x)'=2\cdot\sin2x\cdot 2\cdot\cos2x=2\cdot\sin4x$.
	} 
\end{ex}

% nguồn: Ok_Toan11_TracNghiem_Lan2
\begin{ex}%[1D7H2-1]%[ma:3a417a5b]
	Đạo hàm của  hàm số $y=\dfrac{\ln x}{x}$ là
	\choice 
		 {$y=-4$}
		{$y'=\dfrac{1+\ln  x}{x^2}$}
		 {$y'=-\dfrac{1}{x^3}$}
		{\True  $y'=\dfrac{1-\ln x}{x^2}$}
	 \loigiai{
		 Ta có $y'=\left(\dfrac{\ln x}{x} \right)'=\dfrac{\left(\ln x\right)'\cdot x-x'\cdot\ln x}{x^2}=\dfrac{1-\ln x}{x^2}$.
	} 
\end{ex}

% nguồn: Ok_Toan11_TracNghiem_Lan2
\begin{ex}%[1D7H2-1]%[ma:f0c491bd]
	 Đạo hàm của hàm số $y=e^{2-x}$ là
	\choice 
		 {$y'=e^{2-x}$}
		 {$y'=2e^{2-x}$}
		{\True  $y'=-e^{2-x}$}
		{$y'=(2-x)e^{2-x}$} 
	\loigiai{ 
		Ta có $y'=(2-x)'e^{2-x}  = -e^{2-x}$.
	} 
\end{ex}

% nguồn: Ok_Toan11_TracNghiem_Lan2
\begin{ex}%[1D7H2-1]%[ma:bc5a604d]
	Cho hàm số $f(x)=\dfrac{2x+a}{x-b}$ $(a,b\in\mathbb{R};b\neq 1)$.  Ta có $f '(1)$ bằng
	 \choice
		 {$\dfrac{-a+2b}{(b-1)^2}$}
		 {$\dfrac{a-2b}{(b-1)^2}$}
		{$\dfrac{a+2b}{(b-1)^2}$} 
		{\True $\dfrac{-a-2b}{(b-1)^2}$} 
	 \loigiai{
		Ta có $f'(x)=\dfrac{2(x-b)-2x-a}{(x-b)^2}=\dfrac{-a-2b}{(x-b)^2}$.\\ 
		Do đó, $f'(1)=\dfrac{-a-2b}{(b-1)^2}$. 
	} 
\end{ex}

% nguồn: Ok_Toan11_TracNghiem_Lan2
\begin{ex}%[1D7H2-1]%[ma:78fa69e7]
	 Tính đạo hàm của hàm số $y=9^{2x+1}$.
	 \choice
		{\True $y'=2 \cdot 9^{2x+1} \cdot  \ln~9$}
		{$y'=(2x+1) \cdot  9^{2x+1}$}
		{$y'=9^{2x+1} \cdot \ln~9$} 
		{$y'=(2x+1) \cdot  9^{2x+1} \cdot \ln~9$}
	 \loigiai{
		 Ta có
		$(9^{2x+1})'=(2x+1)' \cdot  9^{2x+1}\ln~9=2 \cdot 9^{2x+1}\ln~9$.
	} 
\end{ex}

% nguồn: Ok_Toan11_TracNghiem_Lan2
\begin{ex}%[1D7H2-1]%[ma:5d6882f5]
	Cho  hàm số $y=\sqrt{2x^2+5x-4}$. Đạo hàm $y'$ của hàm số là
	 \choice
		{\True  $y'=\dfrac{4x+5}{2\sqrt{2x^2+5x-4}}$}
		{$y'=\dfrac{2x+5}{2\sqrt{2x^2+5x-4}}$} 
		 {$y'=\dfrac{2x+5}{\sqrt{2x^2+5x-4}}$}
		 {$y'=\dfrac{4x+5}{\sqrt{2x^2+5x-4}}$}
	\loigiai{Ta  có $y'=\left(\sqrt{2x^2+5x-4}\right)'=\dfrac{\left(2x^2+5x-4\right)'}{2\sqrt{2x^2+5x-4}}=\dfrac{4x+5}{2\sqrt{2x^2+5x-4}}$.
	 }
\end{ex}

% nguồn: Ok_Toan11_TracNghiem_Lan2
\begin{ex}%[1D7H2-1]%[ma:5ab9aa0a]
	Tính  đạo hàm của hàm số $y=x(x+1)(x+2)(x+3)$ tại điểm $x=0$ là
	\choice 
		{$y'(0)=5$} 
		{\True $y'(0)=6$} 
		{$y'(0)=0$} 
		{$y'(0)=-6$} 
	\loigiai{Ta  có $y'=(x+1)(x+2)(x+3)+x(x+2)(x+3)+x(x+1)(x+3)+x(x+1)(x+2)$.\\
		 Vậy $y'(0)=6$.
	 }
\end{ex}

% nguồn: Ok_Toan11_TracNghiem_Lan2
\begin{ex}%[1D7H2-1]%[ma:0a830ce4]
	Đạo hàm của hàm số $y=\log_4 \left(2x^2-3\right)$  là
	\choice 
		 {$y'=\dfrac{4x}{\left(2x^2-3\right)\ln 2}$}
		{$y'=\dfrac{4x}{2x^2-3}$} 
		{$y'=\dfrac{1}{\left(2x^2-3\right)\ln  4}$}
		{\True  $y'=\dfrac{2x}{\left(2x^2-3\right)\ln 2}$}
	 \loigiai{
		Ta  có $y'=\dfrac{4x}{\left(2x^2-3\right)\ln 4}=\dfrac{4x}{\left(2x^2-3\right)2\ln 2}=\dfrac{2x}{\left(2x^2-3\right)\ln 2}$.
	} 
\end{ex}

% nguồn: Ok_Toan11_TracNghiem_Lan2
\begin{ex}%[1D7H2-1]%[ma:834cdcdc]
	Cho hàm số $f(x)=\sqrt{x^2-2x+3}$, tính $f'(2)$. 
	\choice 
		 {$\dfrac{1}{3}$}
		{$\sqrt{3}$} 
		 {\True $\dfrac{\sqrt{3}}{3}$}
		 {$2\sqrt{3}$}
	\loigiai{Ta  có $f'(x)=\dfrac{x-1}{\sqrt{x^2-2x+3}}\Rightarrow f'(2)=\dfrac{\sqrt{3}}{3}$.
	} 
\end{ex}

% nguồn: Ok_Toan11_TracNghiem_Lan2
\begin{ex}%[1D7H2-1]%[ma:f3ecd8fc]
	Hàm số $y=\dfrac{1}{x^2+5}$  có đạo hàm bằng
	\choice 
		 {$y'=\dfrac{1}{\left(x^2+5\right)^2}$}
		{$y'=\dfrac{2x}{\left(x^2+5\right)^2}$} 
		 {$y'=\dfrac{-1}{\left(x^2+5\right)^2}$}
		{\True $y'=\dfrac{-2x}{\left(x^2+5\right)^2}$} 
	\loigiai{Ta  có $y'=-\dfrac{\left(x^2+5\right)'}{\left(x^2+5\right)^2}=\dfrac{-2x}{\left(x^2+5\right)^2}$.
	 }
\end{ex}

% nguồn: Ok_Toan11_TracNghiem_Lan2
\begin{ex}%[1D7H2-1]%[ma:1ed90f4a]
	Tính đạo hàm của  hàm số $ y=x^3+2x+1$.
	 \choice
		{$y'=3x^2+2x$} 
		 {\True $y'=3x^2+2$}
		 {$y'=3x^2+2x+1$}
		 {$y'=x^2+2$}
	 \loigiai{
		Ta có  $y'=3x^2+2$.
	} 
\end{ex}

% nguồn: Ok_Toan11_TracNghiem_Lan2
\begin{ex}%[1D7H2-1]%[ma:9cdb421f]
	 Tính đạo hàm của hàm số $y=\dfrac{x+1}{4^{x}}$.
	\choice 
		{\True $y'=\dfrac{1-2(x+1)\ln~2}{2^{2x}}$} 
		 {$y'=\dfrac{1+2(x+1)\ln~2}{2^{2x}}$}
		 {$y'=\dfrac{1-2(x+1)\ln~2}{2^{x^{2}}}$}
		{$y'=\dfrac{1+2(x+1)\ln~2}{2^{x^{2}}}$} 
	 \loigiai{
		\begin{itemize}
			\item Ta  có $y'=\dfrac{(x+1)' \cdot 4^{x}-(x+1) \cdot (4^{x})'}{(4^{x})^{2}}=\dfrac{4^{x}-(x+1) \cdot 4^{x} \cdot \ln~4}{(4^{x})^{2}}$.
			 \item $y' = \dfrac{1-(x+1)\ln 4}{4^x} = \dfrac{1-(x+1)2\ln 2}{2^{2x}} = \dfrac{1-2(x+1)\ln~2}{2^{2x}}$.
		\end{itemize}
	 }
\end{ex}

% nguồn: Ok_Toan11_TracNghiem_Lan2
\begin{ex}%[1D7H2-1]%[ma:a5c9b944]
	Cho hàm  số $y=\dfrac{x+2}{x-1}$. Tính $y'(3)$.
	 \choice
		{$\dfrac{5}{2}$} 
		 {$\dfrac{3}{4}$}
		{$\dfrac{3}{2}$} 
		{\True  $-\dfrac{3}{4}$}
	 \loigiai{
		 Ta có $y'=\dfrac{-3}{(x-1)^2}$.\\
		Vậy $y'(3)=\dfrac{-3}{(3-1)^2}=\dfrac{-3}{4}$. 
	 }
\end{ex}

% nguồn: Ok_Toan11_TracNghiem_Lan2
\begin{ex}%[1D7H2-1]%[ma:291d0b76]
	Trên khoảng $\left(\dfrac{1}{2};+\infty \right)$, đạo hàm của hàm số  $y=\log (2x-1)$ là
	 \choice
		{$y'=\dfrac{1}{\left(2x-1\right)\ln 10}$} 
		{\True $y'=\dfrac{2}{\left(2x-1\right)\ln 10}$} 
		 {$y'=\dfrac{2}{2x-1}$}
		 {$y'=\dfrac{1}{2x-1}$}
	\loigiai{ 
		Trên khoảng $\left(\dfrac{1}{2};+\infty  \right)$, ta có $y'=\dfrac{(2x-1)'}{(2x-1)\ln 10}=\dfrac{2}{(2x-1)\ln 10}$.
	 }
\end{ex}

% nguồn: Ok_Toan11_TracNghiem_Lan2
\begin{ex}%[1D7H2-1]%[ma:2545c582]
	Đạo hàm của hàm số $y=3^{2x-1}$  là
	\choice 
		{$2  \cdot 3^{2x-1}$}
		{$3^{2x-1} \cdot  \ln~3$}
		{$2 \cdot 3^{2x-1} \cdot  \ln~2$}
		{\True $2  \cdot 3^{2x-1} \cdot \ln~3$}
	 \loigiai{
		 Ta có $y'=(2x-1)' \cdot 3^{2x-1} \cdot \ln~3 = 2 \cdot 3^{2x-1} \cdot \ln~3$.
	} 
\end{ex}

% nguồn: Ok_Toan11_TracNghiem_Lan2
\begin{ex}%[1D7H2-1]%[ma:16fce109]
	 Đạo hàm của hàm số $y=\sqrt{x^2+x+1}$ là
	\choice 
		 {\True $y'=\dfrac{2x+1}{2\sqrt{x^2+x+1}}$}
		{$y'=\dfrac{x+1}{\sqrt{x^2+x+1}}$} 
		 {$y'=\dfrac{2(2x+1)}{\sqrt{x^2+x+1}}$}
		{$y'=\dfrac{2x+1}{\sqrt{x^2+x+1}}$} 
	\loigiai{ 
		 Ta có $y'=\dfrac{\left(x^2+x+1\right)'}{2\sqrt{x^2+x+1}}=\dfrac{2x+1}{2\sqrt{x^2+x+1}}$.
	} 
\end{ex}

% nguồn: Ok_Toan11_TracNghiem_Lan2
\begin{ex}%[1D7H2-1]%[ma:3600825e]
	Tính đạo hàm của hàm số  $y=\left(x^3-5\right)\cdot\sqrt{x}$.
	 \choice
		 {$y'=\dfrac{7}{2}\sqrt{x^3}-\dfrac{5}{2\sqrt{x}}$}
		 {\True $y'=\dfrac{7}{2}\sqrt{x^5}-\dfrac{5}{2\sqrt{x}}$}
		 {$y'=3x^2-\dfrac{5}{2\sqrt{x}}$}
		 {$y'=3x^2-\dfrac{1}{2\sqrt{x}}$}
	 \loigiai{Ta có $y'=3x^2\cdot\sqrt{x}+\left(x^3-5\right)\cdot\dfrac{1}{2\sqrt{x}}=3x^2\sqrt{x}+\dfrac{1}{2}x^2\sqrt{x}-\dfrac{5}{2\sqrt{x}}=\dfrac{7}{2}x^2\sqrt{x}-\dfrac{5}{2\sqrt{x}}=\dfrac{7}{2}\sqrt{x^5}-\dfrac{5}{2\sqrt{x}}$.
	} 
\end{ex}

% nguồn: Ok_Toan11_TracNghiem_Lan2
\begin{ex}%[1D7H2-1]%[ma:3b1c8d5e]
	Đạo  hàm của hàm số $y=\dfrac{1}{4}-\dfrac{1}{3}x+x^2-0{,}25x^4$ là
	\choice 
		 {$y'=\dfrac{1}{3}x-2x-2x^3$}
		{$y'=-\dfrac{1}{3}+x-2x^3$} 
		 {$y'=\dfrac{4}{3}+x-2x^3$}
		 {\True $y'=-\dfrac{1}{3}+2x-x^3$}
	 \loigiai{
		Ta có  $y'=-\dfrac{1}{3}+2x-x^3$.
	 }
\end{ex}

% nguồn: Ok_Toan11_TracNghiem_Lan2
\begin{ex}%[1D7H2-1]%[ma:aeb7551e]
	Đạo hàm  của hàm số $y=(2-\sqrt{3})^{x}$ trên tập $\mathbb{R}$ là
	 \choice
		 {\True $y'=(2+\sqrt{3})^{-x}\ln(2-\sqrt{3})$}
		{$y'=(2+\sqrt{3})^{-x}\ln(2+\sqrt{3})$} 
		{$y'=(2+\sqrt{3})^{x}\ln(2-\sqrt{3})$} 
		 {$y'=(2-\sqrt{3})^{x}\ln(2+\sqrt{3})$}
	 \loigiai{
		 Ta có $y'=(2-\sqrt{3})^{x}\ln(2-\sqrt{3})=\dfrac{1}{(2+\sqrt{3})^{x}}\ln(2-\sqrt{3})=(2+\sqrt{3})^{-x}\ln(2-\sqrt{3})$.
	} 
\end{ex}

% nguồn: Ok_Toan11_TracNghiem_Lan2
\begin{ex}%[1D7H2-1]%[ma:796bfa6a]
	Đạo hàm của  hàm số $f(x)=\sqrt{2-3x^2}$ bằng biểu thức nào sau đây?
	\choice 
		 {\True $\dfrac{-3x}{\sqrt{2-3x^2}}$}
		 {$\dfrac{1}{2\sqrt{2-3x^2}}$}
		{$\dfrac{-6x^2}{2\sqrt{2-3x^2}}$} 
		{$\dfrac{3x}{\sqrt{2-3x^2}}$} 
	\loigiai{ 
		Ta có $f'(x)=\left(\sqrt{2-3x^2}\right)’=\dfrac{(2-3x^2)’}{2\sqrt{2-3x^2}}=\dfrac{-6x}{2\sqrt{2-3x^2}}=\dfrac{-3x}{\sqrt{2-3x^2}}$. 
	 }
\end{ex}

% nguồn: Ok_Toan11_TracNghiem_Lan2
\begin{ex}%[1D7H2-1]%[ma:dbf87554]
	Cho $f(x)=\sqrt{1-4x}+\dfrac{1-x}{x-3}$, với  $x\in\left(-\infty;\dfrac{1}{4}\right)$. Tìm $f'(x)$.
	\choice 
		{$\dfrac{2}{\sqrt{1-4x}}+\dfrac{2}{x-3}$} 
		 {$\dfrac{2}{\sqrt{1-4x}}+\dfrac{2}{(x-3)^2}$}
		 {$\dfrac{1}{2\sqrt{1-4x}}+1$}
		{\True $\dfrac{-2}{\sqrt{1-4x}}+\dfrac{2}{(x-3)^2}$} 
	\loigiai{ 
		Ta có  $f'(x)=\dfrac{(1-4x)’}{2\sqrt{1-4x}}+\dfrac{(1-x)’(x-3)-(1-x)(x-3)’}{(x-3)^2}=\dfrac{-2}{\sqrt{1-4x}}+\dfrac{2}{(x-3)^2}$.
	} 
\end{ex}

% nguồn: Ok_Toan11_TracNghiem_Lan2
\begin{ex}%[1D7H2-1]%[ma:2d76b9be]
	 Đạo hàm của hàm số $y=x^4-4x^2-3$ là
	\choice 
		{$y'=-4x^3+8x$} 
		 {$y'=4x^2-8x$}
		{\True  $y'=4x^3-8x$}
		 {$y'=-4x^2+8x$}
	\loigiai{Ta  có $y'=\left(x^4-4x^3-3\right)'=4x^3-8x$.
	} 
\end{ex}

% nguồn: Ok_Toan11_TracNghiem_Lan2
\begin{ex}%[1D7H2-1]%[ma:75a30ba3]
	Đạo hàm  của hàm số $y=\ln \left(1-x^2\right)$ là
	\choice 
		{\True  $\dfrac{2x}{x^2-1}$}
		{$\dfrac{-2x}{x^2-1}$} 
		 {$\dfrac{1}{1-x^2}$}
		{$\dfrac{1}{x^2-1}$} 
	 \loigiai{
		Ta  có $y'=\dfrac{\left(1-x^2\right)'}{1-x^2}=\dfrac{-2x}{1-x^2}=\dfrac{2x}{x^2-1}$.
	} 
\end{ex}

% nguồn: Ok_Toan11_TracNghiem_Lan2
\begin{ex}%[1D7H2-1]%[ma:b57a22c8]
	Đạo hàm của  hàm số $y=(-x^2+3x+7)^7$ là
	\choice 
		 {\True $y'=7(-2x+3)(-x^2+3x+7)^6$}
		 {$y'=7(-x^2+3x+7)^6$}
		 {$y'=(-2x+3)(-x^2+3x+7)^6$}
		 {$y'=7(-2x+3)(-x^2+3x+7)^7$}
	 \loigiai{
		Ta  có $y'=7(-x^2+3x+7)^6(-x^2+3x+7)’=7(-2x+3)(-x^2+3x+7)^6$.
	} 
\end{ex}

% nguồn: Ok_Toan11_TracNghiem_Lan2
\begin{ex}%[1D7H2-1]%[ma:ed64e9ee]
	Cho hàm số  $y=\dfrac{3x+5}{-1+2x}$. Đạo hàm $y'$ của hàm số là
	\choice 
		 {$\dfrac{7}{(2x-1)^2}$}
		{$\dfrac{1}{(2x-1)^2}$} 
		 {\True $\dfrac{-13}{(2x-1)^2}$}
		 {$\dfrac{13}{(2x-1)^2}$}
	\loigiai{Ta có $y'=\dfrac{(3x+5)'\cdot (2x-1)-(3x+5)\cdot (2x-1)'}{(2x-1)^2}=\dfrac{3(2x-1)-2(3x+5)}{(2x-1)^2}=\dfrac{-13}{(2x-1)^2}$. 
	 }
\end{ex}

% nguồn: Ok_Toan11_TracNghiem_Lan2
\begin{ex}%[1D7H2-1]%[ma:d618d5ba]
	Đạo hàm của hàm số  $y=\pi^{x^{2}+2x}$ là
	 \choice
		{$y'=(2x+2) \cdot \pi^{x^{2}+2x}$} 
		 {$y'=(2x+2) \cdot \dfrac{\pi^{x^{2}+2x}}{\ln~\pi}$}
		{\True  $y'=(2x+2) \cdot \pi^{x^{2}+2x} \cdot \ln~\pi$}
		{$y'=\dfrac{(2x+2)}{\pi^{x^{2}+2x} \cdot  \ln~\pi}$}
	\loigiai{ 
		Ta có $y=\pi^{x^{2}+2x} \Rightarrow y'=(x^{2}+2x)' \cdot \pi^{x^{2}+2x} \cdot \ln~\pi=(2x+2) \cdot \pi^{x^{2}+2x} \cdot \ln~\pi$. 
	 }
\end{ex}

% nguồn: Ok_Toan11_TracNghiem_Lan2
\begin{ex}%[1D7H2-1]%[ma:a042b43e]
	Cho hàm số $y=\dfrac{x^2+x}{x-2}$. Đạo hàm của hàm số  tại $x=1$ là
	 \choice
		 {$y'(1)=-4$}
		{\True  $y'(1)=-5$}
		{$y'(1)=-3$} 
		{$y'(1)=-2$} 
	\loigiai{Ta có $y'=\dfrac{\left(x^2+x\right)'\cdot  (x-2)-\left(x^2+x\right)\cdot (x-2)'}{(x-2)^2}=\dfrac{(2x+1)(x-2)-\left(x^2+x\right)\cdot 1}{(x-2)^2}=\dfrac{x^2-4x-2}{(x-2)^2}$.\\
		 Vậy $y'(1)=-5$.
	 }
\end{ex}

% nguồn: Ok_TracNghiemLop11_Lan1
\begin{ex}%[1D7H2-1]%[ma:eb06ba45]
	Tính đạo hàm của  hàm số $f(x)=\mathrm{e}^{2x-3}$.
	 \choice
		{\True  $f'(x)=2\cdot\mathrm{e}^{2x-3}$}
		{$f'(x)=\mathrm{e}^{2x-3}$} 
		{$f'(x)=-2\cdot\mathrm{e}^{2x-3}$} 
		 {$f'(x)=2\cdot\mathrm{e}^{x-3}$}
	 \loigiai{
	 Ta có $f'(x)=(2x-3)'\mathrm{e}^{2x-3}=2\cdot\mathrm{e}^{2x-3}$.}
\end{ex}

% nguồn: Ok_TracNghiemLop11_Lan1
\begin{ex}%[1D7H2-1]%[ma:41b1732c]
	 Hàm số $y=\dfrac{1}{x^2+5}$ có đạo hàm bằng:
	\choice 
		{$y'=\dfrac{1}{\left(x^2+5\right)^2}$} 
		 {$y'=\dfrac{2x}{\left(x^2+5\right)^2}$}
		 {$y'=\dfrac{-1}{\left(x^2+5\right)^2}$}
		 {\True $y'=\dfrac{-2x}{\left(x^2+5\right)^2}$}
	\loigiai{ 
		$y'=\dfrac{-2x}{\left(x^2+5\right)^2}$ 
	 }
\end{ex}

% nguồn: Ok_TracNghiemLop11_Lan1
\begin{ex}%[1D7H2-1]%[ma:085fa0b9]
	Tính đạo hàm của hàm số $f(x)=\dfrac{2x+7}{x+4}$ tại $x=2$ ta  được
	 \choice
		{\True  $f'\left(2\right)=\dfrac{1}{36}$}
		{$f'\left(2\right)=\dfrac{11}{6}$} 
		 {$f'\left(2\right)=\dfrac{3}{2}$}
		 {$f'\left(2\right)=\dfrac{5}{12}$}
	\loigiai{ 
		Ta có  $f'(x)=\dfrac{1}{(x+4)^2}\Rightarrow f'\left(2\right)=\dfrac{1}{36}$.
	} 
\end{ex}

% nguồn: Ok_TracNghiemLop11_Lan1
\begin{ex}%[1D7H2-1]%[ma:7b50388e]
	Hàm số $y=3x^2+2\sqrt{x}-1$ với $x>0$ có đạo  hàm là
	 \choice
		{\True $y'=6x+\dfrac{1}{\sqrt{x}}$} 
		{$y'=3x^2+\dfrac{1}{\sqrt{x}}$} 
		{$y'=3x^2+\dfrac{2}{\sqrt{x}}$} 
		{$y'=6x+\dfrac{2}{\sqrt{x}}$} 
	 \loigiai{
		Ta có:  $y'=6x+\dfrac{1}{\sqrt{x}}$.
	} 
\end{ex}

% nguồn: Ok_TracNghiemLop11_Lan1
\begin{ex}%[1D7H2-1]%[ma:387506ff]
	Tính đạo hàm của hàm số $y=x\left(x+1\right)\left(x+2\right)\left(x+3\right)$ tại điểm  $x_0=0$ là
	 \choice
		 {$y'\left(0\right)=5$}
		{\True $y'\left(0\right)=6$} 
		{$y'\left(0\right)=0$} 
		{$y'\left(0\right)=-6$} 
	 \loigiai{
		Ta  có $y=x\left(x+1\right)\left(x+2\right)\left(x+3\right)=\left(x^2+x\right)\left(x^2+5x+6\right)$
		 $\\
		\Rightarrow y'=\left(2x+1\right)\left(x^2+5x+6\right)+\left(x^2+x\right)\left(2x+5\right)$ 
	 $\Rightarrow y'\left(0\right)=6$.}
\end{ex}

% nguồn: Ok_TracNghiemLop11_Lan1
\begin{ex}%[1D7H2-1]%[ma:056815c5]
	Tính đạo hàm  của hàm số $f(x)=\dfrac{m-nx}{3x-1}$ với $m, n$ là tham số.
	\choice 
		 {$\dfrac{3m-n}{(3x-1)^2}$}
		{$\dfrac{m-3n}{(3x-1)^2}$} 
		{$\dfrac{-m+3n}{(3x-1)^2}$} 
		{\True  $\dfrac{n-3m}{(3x-1)^2}$}
	 \loigiai{$f'(x)=\left(\dfrac{-nx+m}{3x-1}\right)'=\dfrac{(-n)(-1)-3m}{(3x-1)^2}=\dfrac{n-3m}{\left(3x-1\right)^2}.$}
\end{ex}

% nguồn: Ok_TracNghiemLop11_Lan1
\begin{ex}%[1D7H2-1]%[ma:72a35291]
	Cho hàm số  $f(x)=\dfrac{1}{3}x^3-2x^2-5x+1$. Tập nghiệm của phương trình $f'(x)=0$ là
	\choice 
		{\True  $\left\{-1;5\right\}$}
		 {$\varnothing$}
		{$\left\{1;-5\right\}$} 
		 {$\left\{2-\sqrt{5};2+\sqrt{5}\right\}$}
	 \loigiai{
		Ta có $f'(x)=x^2-4x-5$. \\ 
		Do  đó $f'(x)=0 \Leftrightarrow x^2-4x-5=0 \Leftrightarrow \hoac{
		 &x=-1 \\
		&x=5.} $ 
	 }
\end{ex}

% nguồn: Ok_TracNghiemLop11_Lan1
\begin{ex}%[1D7H2-1]%[ma:b77ddb77]
	Tính đạo hàm của  hàm số $y=x\left(x+1\right)\left(x+2\right)\left(x+3\right)$ tại điểm $x_0=0$ là
	\choice 
		 {$y'(0)=5$}
		 {\True $y'(0)=6$}
		{$y'(0)=0$} 
		 {$y'(0)=-6$}
	 \loigiai{
		Ta có $y=x\left(x+1\right)\left(x+2\right)\left(x+3\right)=\left(x^2+x\right)\left(x^2+5x+6\right)$\\ 
		 $\Rightarrow{y}'=\left(2x+1\right)\left(x^2+5x+6\right)+\left(x^2+x\right)\left(2x+5\right)$
	 $\Rightarrow{y}'(0)=6.$}
\end{ex}

% nguồn: Ok_TracNghiemLop11_Lan1
\begin{ex}%[1D7H2-1]%[ma:f2e4ec18]
	Hàm số $y=\ln(2x+1)$ có đạo  hàm là
	\choice 
		{$y'=\dfrac{2}{x\ln(2x+1)}$} 
		{$y'=\dfrac{1}{2x+1}$} 
		{\True  $y'=\dfrac{2}{2x+1}$}
		{$y'=\dfrac{1}{(2x+1)\ln 2}$} 
	\loigiai{ 
	Hàm số $y=\ln(2x+1)$  có đạo hàm là $y'=\dfrac{2}{2x+1}$.}
\end{ex}

% nguồn: Ok_TracNghiemLop11_Lan1
\begin{ex}%[1D7H2-1]%[ma:7986f66e]
	Tìm mệnh đề \textbf{sai} trong các  mệnh đề sau
	\choice 
		{Hàm số $y=x^n (n\in \mathbb{N},n\ge 2)$ có đạo hàm trên $\mathbb{R}$  và $y'=n{x}^{n-1}$}
		 {Hàm số hằng $y=c$ có đạo hàm trên $\mathbb{R}$ và $y'=0$}
		{\True Hàm số $y=\sqrt{x}$ có đạo hàm trên khoảng  $(0;+\infty)$ và $y'=\dfrac{1}{\sqrt{x}}$}
		{Hàm số $y=x$ có  đạo hàm trên $\mathbb{R}$ và $y'=1$}
	 \loigiai{Hàm số $y=\sqrt{x}$ có đạo hàm trên khoảng $(0;+\infty)$ và $y'=\dfrac{1}{2\sqrt{x}}$.}
\end{ex}

% nguồn: Ok_TracNghiemLop11_Lan1
\begin{ex}%[1D7H2-1]%[ma:0a30ce8d]
	Đạo hàm của  hàm số $y=\sqrt{2x+1}-\dfrac{1}{x}$ là
	\choice 
		{\True  $\dfrac{1}{\sqrt{2x-1}}+\dfrac{1}{x^2}$}
		{$\dfrac{2x}{\sqrt{2x-1}}-\dfrac{1}{x^2}$} 
		 {$\dfrac{2x}{\sqrt{2x-1}}+\dfrac{1}{x^2}$}
		 {$\dfrac{1}{\sqrt{2x-1}}-\dfrac{1}{x^2}$}
	\loigiai{ 
		Ta có $y'=\dfrac{(2x-1)'}{2\sqrt{2x  -1}}-\left(-\dfrac{1}{x^2}\right)=\dfrac{1}{\sqrt{2x -1}}+\dfrac{1}{x^2}$.
	} 
\end{ex}

% nguồn: Ok_TracNghiemLop11_Lan1
\begin{ex}%[1D7H2-1]%[ma:48b42a94]
	Hàm số có đạo hàm  bằng $2x+\dfrac{1}{x^2}$ là
	 \choice
		 {$y=\dfrac{x^3+1}{x}$}
		{\True $y=\dfrac{x^3+5x-1}{x}$} 
		 {$y=2-\dfrac{2}{x^3}$}
		{$y=2+\dfrac{1}{x^2}$} 
	\loigiai{Ta  có:
		\begin{itemize}
			 \item $y'=\left(\dfrac{x^3+1}{x}\right)'=\left(x^2+\dfrac{1}{x}\right)'=2x-\dfrac{1}{x^2}$.
			\item $y'= \left(\dfrac{x^3+5x-1}{x}\right)'=\left(x^2+5-\dfrac{1}{x}\right)'=2x+\dfrac{1}{x^2} $. 
			\item $y'=\left(2-\dfrac{2}{x^3}\right)'=\dfrac{6}{x^4} $. 
			\item $  y'=\left(2+\dfrac{1}{x^2}\right)'=-\dfrac{2}{x^3} $.
		\end{itemize}
	 }
\end{ex}

% nguồn: Ok_TracNghiemLop11_Lan1
\begin{ex}%[1D7H2-1]%[ma:4842ff9b]
	Đạo hàm của hàm số $y=-x^3+3mx^2+3\left(1-m^2\right)x+m^3-m^2$ (với $m$ là tham số)  bằng
	 \choice
		 {$3x^2-6mx-3+3m^2$}
		{$-x^2+3mx-1-3m$} 
		 {$-3x^2+6mx+1-m^2$}
		{\True $-3x^2+6mx+3-3m^2$} 
	 \loigiai{
		Ta có  $y'=-3x^2+6mx+3-3m^2$.
	} 
\end{ex}

% nguồn: Ok_TracNghiemLop11_Lan1
\begin{ex}%[1D7H2-1]%[ma:bc77c876]
	Cho hàm số $y=x^2-x+2$. Tính $y'(1)$ . 
	\choice 
		{$y'(1)=-1$} 
		{\True $y'(1)=1$} 
		{$y'(1)=2$} 
		 {$ y'(1)=0$}
	\loigiai{ 
	$y'=2x-1$  $\Rightarrow y'(1)=2\cdot 1-1=1$.}
\end{ex}

% nguồn: Ok_TracNghiemLop11_Lan1
\begin{ex}%[1D7H2-1]%[ma:f03d5b15]
	Cho hàm số $f(x)=x^4+2x^2-3$. Tìm $x$  để $f'(x)>0$?
	\choice 
		 {$-1< x <0$}
		 {$x < 0$}
		{\True $x >0$} 
		{$x  <-1$}
	\loigiai{ 
		 $f'(x) > 0\Leftrightarrow 4x^3+4x > 0\Leftrightarrow 4x\left(x^2+1\right) > 0\Leftrightarrow x > 0$.
	 }
\end{ex}

% nguồn: Ok_TracNghiemLop11_Lan1
\begin{ex}%[1D7H2-1]%[ma:f4d95908]
	Cho hàm số  $f(x) = \dfrac{x^3}{3} - \dfrac{x^2}{2} - 12x + \dfrac{1}{3}$. Tập nghiệm $\mathscr{T}$ của bất phương trình $f'(x) \leq 0$ là
	 \choice
		 {$\mathscr{T} \left(- \infty; - 3 \right] \cup \left[4; + \infty \right)$}
		{$\mathscr{T} =  \varnothing$}
		{\True $\mathscr{T} =\left[-3;4 \right]$} 
		{$\mathscr{T}  = \left(- \infty; + \infty \right)$}
	 \loigiai{
		Tập xác định: $\mathscr{D} =  \mathbb{R}$. \\
		Ta có: $f'(x) = x^2 - x - 12$. \\ 
		$f'(x) \leq 0 \Leftrightarrow -3  \leq x \leq 4$.
	} 
\end{ex}

% nguồn: Ok_TracNghiemLop11_Lan1
\begin{ex}%[1D7H2-1]%[ma:4d5231bb]
	Cho hai hàm số $f(x)=x+2$; $g(x)=\dfrac{1}{1-x}$.  Giá trị của $\dfrac{f'(1)}{g'(0)}$ bằng
	 \choice
		{\True $1$} 
		 {$-2$}
		{$0$} 
		 {$2$}
	\loigiai{ 
		Ta có:$f'(x)=1 \Rightarrow f'(1)=1$; $g'(x)=\dfrac{1}{(1-x)^2}\Rightarrow g'(0)=1$.  Do đó $\dfrac{f'(1)}{g'(0)}=\dfrac{1}{1}=1$.
	 }
\end{ex}

% nguồn: Ok_TracNghiemLop11_Lan1
\begin{ex}%[1D7H2-1]%[ma:26822236]
	Đạo hàm của hàm số $y=\dfrac{2x+1}{1-x}$  là
	\choice 
		{$y'=\dfrac{-4x+1}{(1-x)^2}$} 
		 {$y'=\dfrac{-3}{(1-x)^2}$}
		{\True  $y'=\dfrac{3}{(1-x)^2}$}
		 {$y'=\dfrac{4x-1}{(1-x)^2}$}
	\loigiai{ 
		Tập xác định $\mathscr{D}=\mathbb{R}\setminus  \{1\}$.\\
		 Ta có
		 \allowdisplaybreaks
		\begin{eqnarray*}
			y'&=&\dfrac{(2x+1)'(1-x)-(1-x)'(2x+1)}{(1-x)^2}\\ 
			 &=&\dfrac{2(1-x)-(-1)(2x+1)}{(1-x)^2}\\
			&=&\dfrac{2-2x+2x+1}{(1-x)^2}\\ 
			 &=&\dfrac{3}{(1-x)^2}.
		\end{eqnarray*}
		Vậy $y'=\dfrac{3}{(1-x)^2}$. 
	 }
\end{ex}

% nguồn: Ok_TracNghiemLop11_Lan1
\begin{ex}%[1D7H2-1]%[ma:25de12d0]
	 Tính đạo hàm của hàm số $f(x)=\dfrac{1}{x}$ tại điểm $x=\dfrac{1}{2}$.
	 \choice
		{$f'(\frac{1}{2})=-\dfrac{1}{4}$} 
		{$f'(\frac{1}{2})=\dfrac{1}{4}$} 
		{\True $f'(\frac{1}{2})=-4$} 
		{$f'(\frac{1}{2})=4$} 
	\loigiai{ 
		Ta có $f'(x)=-\dfrac{1}{x^2}$, thay $x=\dfrac{1}{2}$  ta được $f'(\frac{1}{2})=-4$.
	} 
\end{ex}

% nguồn: Ok_TracNghiemLop11_Lan1
\begin{ex}%[1D7H2-1]%[ma:56a988f6]
	Cho hàm số $y=(x^2-3\cos  x)^{10}$. Tính $y'(0)$.
	 \choice
		 {$y'(0)=-10\cdot 3^9$}
		 {$y'(0)=1$}
		{\True $y'(0)=0$} 
		{$y'(0)=10\cdot  3^9$}
	\loigiai{ 
		Có $y'=10(x^2-3\cos x)^9(2x+3  \sin x) \Rightarrow y'(0)=0$.
	} 
\end{ex}

% nguồn: Ok_TracNghiemLop11_Lan1
\begin{ex}%[1D7H2-1]%[ma:7db7924c]
	Tính đạo hàm của hàm số  $y=\dfrac{2x}{x-1}$.
	\choice 
		 {$y'=\dfrac{-1}{(x-1)^2}$}
		 {$y'=\dfrac{2}{(x-1)^2}$}
		{$y'=\dfrac{3}{(x-1)^2}$} 
		 {\True $y'=\dfrac{-2}{(x-1)^2}$}
	\loigiai{ 
		Ta  có $y'=\dfrac{-2}{(x-1)^2}$.
	} 
\end{ex}

% nguồn: Ok_TracNghiemLop11_Lan1
\begin{ex}%[1D7H2-1]%[ma:2496ffa4]
	 Cho $f(x)=x+\sqrt{4-x^2}$. Phương trình $f'(x)=0$ có bao nhiêu nghiệm?
	 \choice
		 {$0$}
		{$3$} 
		 {\True $1$}
		 {$2$}
	\loigiai{ 
		Ta  có $f'(x)=1+\dfrac{-x}{\sqrt{4-x^2}}=\dfrac{\sqrt{4-x^2}-x}{\sqrt{4-x^2}}$.\\
		$f'(x)=0 \Leftrightarrow \heva{&\sqrt{4-x^2}-x=0\\&4-x^2>0} \Leftrightarrow \heva{&0\leq x<2\\ &4-x^2=x^2} \Leftrightarrow \heva{&0\leq  x<2\\ & x=\pm\sqrt{2}} \Leftrightarrow x=\sqrt{2}$.\\
		 Vậy phương trình $f'(x)=0$ có 1 nghiệm.
	 }
\end{ex}

% nguồn: Ok_TracNghiemLop11_Lan1
\begin{ex}%[1D7H2-1]%[ma:183d293e]
	Cho hàm số $ y=\dfrac{4}{x-1}$. Khi  đó $y'\left(-1\right)$bằng
	 \choice
		{\True $-1$} 
		 {$-2$}
		 {$2$}
		 {$1$}
	 \loigiai{
	Ta có $y'=-\dfrac{4}{\left(x-1\right)^2}$  $\Rightarrow{y}'\left(-1\right)=-1$.}
\end{ex}

% nguồn: Ok_TracNghiemLop11_Lan1
\begin{ex}%[1D7H2-1]%[ma:dc32e2f8]
	 Đạo hàm của hàm số $y=\sqrt{3x^2-4x+5}$ là
	\choice 
		{$y'=\dfrac{6x-4}{\sqrt{3x^2-4x+5}}$} 
		 {$y'=\dfrac{3x-1}{\sqrt{3x^2-4x+5}}$}
		 {$y'=\dfrac{1}{2\sqrt{3x^2-4x+5}}$}
		{\True $y'=\dfrac{3x-2}{\sqrt{3x^2-4x+5}}$} 
	 \loigiai{
	Ta có $y'=\dfrac{(3x^2-4x+5)'}{2\sqrt{3x^2-4x+5}}=\dfrac{6x-4}{2\sqrt{3x^2-4x+5}}=\dfrac{3x-2}{\sqrt{3x^2-4x+5}}$.} 
\end{ex}

% nguồn: Ok_TracNghiemLop11_Lan1
\begin{ex}%[1D7H2-1]%[ma:477911d4]
	Cho hàm số $y=x^3-3x+2017$. Bất phương trình $y' < 0$  có tập nghiệm là
	\choice 
		{\True $S=\left(-1;1\right)$} 
		 {$S=\left(-\infty;-1\right)\cup \left(1;+\infty \right)$}
		{$S=\left(1;+\infty \right)$} 
		{$S=\left(-\infty;-1\right)$} 
	 \loigiai{
		Ta có 
		\begin{itemize}
			 \item $y=x^3-3x+2017\Rightarrow y'=3x^2-3$.
			\item $y' < 0\Leftrightarrow x^2-1 < 0\Leftrightarrow-1 < x <  1$.
		\end{itemize}
	} 
\end{ex}

% nguồn: Ok_TracNghiemLop11_Lan1
\begin{ex}%[1D7H2-1]%[ma:e63ebf8a]
	Hàm số $y=\sqrt{\cot 2x}$ có đạo  hàm là
	 \choice
		{$-\dfrac{1+\tan^2 2x}{2\sqrt{\cot 2x}}$} 
		{$-\dfrac{1+\tan^2 2x}{\sqrt{\cot  2x}}$}
		{$-\dfrac{1+\cot^2 2x}{2\sqrt{\cot 2x}}$} 
		 {\True$-\dfrac{1+\cot^2 2x}{\sqrt{\cot 2x}}$}
	 \loigiai{$y'=\dfrac{(\cot2x)'}{2\sqrt{\cot2x}}=
		\dfrac{-2(1+\cot ^22x)}{2\sqrt{\cot 2x}}= -\dfrac{1+\cot^2 2x}{\sqrt{\cot 2x}}.$ 
	} 
\end{ex}

% nguồn: Ok_TracNghiemLop11_Lan1
\begin{ex}%[1D7H2-1]%[ma:1e682c72]
	Hàm  số $y=x^3-2x^2-4x + 2018$ có đạo hàm là
	 \choice
		 {$y'=3x^2-4x+2018$}
		 {$y'=3x^2-2x-4$}
		 {\True$y'=3x^2-4x-4$}
		{$y'=x^2-4x-4$} 
	\loigiai{ Ta có  $y'=3x^2-4x-4$.
	} 
\end{ex}

% nguồn: Ok_TracNghiemLop11_Lan1
\begin{ex}%[1D7H2-1]%[ma:f4ba2b07]
	Hàm số $y=2^{x^2-x}$ có đạo hàm là 
	 \choice
		{$2^{x^2-x}\cdot\ln 2$} 
		 {\True$(2x-1)\cdot2^{x^2-x}.\ln 2$}
		 {$(x^2-x)\cdot2^{x^2-x-1}$}
		 {$(2x-1)\cdot2^{x^2-x}$}
	 \loigiai{
		Ta  có $y'=(x^2-x)'\cdot2^{x^2-x}\cdot\ln 2=(2x-1)\cdot2^{x^2-x}\cdot\ln 2$.
	 }
\end{ex}

% nguồn: Ok_TracNghiemLop11_Lan1
\begin{ex}%[1D7H2-1]%[ma:6b6e9627]
	Cho hàm  số $f(x)=\sin2x$. Khi đó giá trị $f'\left(\dfrac{\pi}{4}\right)$ bằng
	\choice 
		{$-1$} 
		{\True  $0$}
		{$1$} 
		 {$2$}
	\loigiai{Ta  có $f'(x)=2\cos2x\Rightarrow f'\left(\dfrac{\pi}{4}\right)=2\cos\left(2\cdot\dfrac{\pi}{4}\right)=2\cos\dfrac{\pi}{2}=0$.}
\end{ex}

% nguồn: Ok_TracNghiemLop11_Lan1
\begin{ex}%[1D7H2-1]%[ma:be0e31a0]
	Tìm  đạo hàm của hàm số $y = \sqrt{\cos(2x - 3)}$.
	 \choice
		{$-\dfrac{\sin  (2x - 3)}{2\sqrt{\cos( 2x - 3)}}$}
		{\True $-\dfrac{\sin (2x - 3)}{\sqrt{\cos( 2x - 3)}}$} 
		{$\dfrac{\sin  (2x - 3)}{2\sqrt{\cos( 2x - 3)}}$}
		{$\dfrac{\sin (2x - 3)}{\sqrt{\cos( 2x  - 3)}}$}
	 \loigiai{
		 Ta có $y' = \dfrac{\left( \cos(2x-3) \right)'}{2\sqrt{\cos(2x-3)}} = -\dfrac{ \sin(2x-3)}{\sqrt{\cos(2x-3)}} $.
	} 
\end{ex}

% nguồn: Ok_TracNghiemLop11_Lan1
\begin{ex}%[1D7H2-1]%[ma:0df7a1a0]
	Hàm số $y=(x+1)\sqrt{x^2+1}$  có đạo hàm là
	 \choice
		 {\True $\dfrac{2x^2+x+1}{\sqrt{x^2+1}}$}
		 {$\dfrac{x^2+x+1}{\sqrt{x^2+1}}$}
		{$\dfrac{x^2+2x+1}{\sqrt{x^2+1}}$} 
		{$\dfrac{x^2+x+2}{\sqrt{x^2+1}}$} 
	\loigiai{ 
	$y^\prime =\sqrt{x^2+1}+(x+1)\cdot  \dfrac{x}{\sqrt{x^2+1}}=\dfrac{x^2+1+x^2+x}{\sqrt{x^2+1}}=\dfrac{2x^2+x+1}{\sqrt{x^2+1}}$.}
\end{ex}

% nguồn: Ok_TracNghiemLop11_Lan1
\begin{ex}%[1D7H2-1]%[ma:41d5363a]
	Cho hàm  số $y=f(x)=x^3-3x^2+12$. Tìm $x$ để $f'(x)<0$.
	\choice 
		{$x\in  (-2;0)$}
		{$x\in (-\infty;-2)\cup(0;+\infty)$} 
		 {$x\in (-\infty;0)\cup(2;+\infty)$}
		{\True $x \in (0;2)$} 
	\loigiai{ 
		Ta có $f'(x)<0 \Leftrightarrow 3x^2-6x<0 \Leftrightarrow 0<x<2$. 
	 }
\end{ex}

% nguồn: Ok_TracNghiemLop11_Lan1
\begin{ex}%[1D7H2-1]%[ma:6bc12cc7]
	Tính đạo hàm của hàm số $y=x\tan 2x$. 
	 \choice
		 {$\tan 2x+\dfrac{2x}{\cos^2x}$}
		{$\tan 2x+\dfrac{x}{\cos^22x}$} 
		{\True $2x\tan^22x+\tan  2x+2x$}
		{$\dfrac{2x}{\cos^22x}$} 
	\loigiai{ 
		Ta có \begin{eqnarray*}
		 y'&=&(x)\tan 2x+x\left(\tan 2x\right)'\\
		 &=&\tan 2x+x\dfrac{(2x)'}{\cos^22x}\\
		&=&\tan 2x+\dfrac{2x}{\cos^22x}\\ 
		&=&2x\tan^22x+\tan  2x+2x.
	\end{eqnarray*}
	} 
\end{ex}

% nguồn: Ok_TracNghiemLop11_Lan1
\begin{ex}%[1D7H2-1]%[ma:280e569c]
	Đạo hàm của hàm số  $y=(x^3-3x^2)^{17}$ là
	\choice 
		{$y'=17(x^3-3x^2)^{16}$} 
		 {$y'=17(x^3-3x^2)^{16}(x^2-3x)$}
		{\True $y'=51(x^3-3x^2)^{16}(x^2-2x)$} 
		 {$y'=17(x^3-3x^2)(3x^2-6x)$}
	\loigiai{ 
	Ta có  $y'=17(x^3-3x^2)^{2016}(3x^2-6x)=51(x^3-3x^2)^{16}(x^2-2x)$.}
\end{ex}

% nguồn: Ok_TracNghiemLop11_Lan1
\begin{ex}%[1D7H2-1]%[ma:6748d967]
	Tính đạo  hàm của hàm số $y=\dfrac{x(1+2x)}{1-x}$.
	\choice 
		 {$y'=\dfrac{1-6x^2}{(1-x)^2}$}
		 {$y'=\dfrac{4x+1}{(1-x)^2}$}
		 {$y'=\dfrac{-6x^2+2x+1}{(1-x)^2}$}
		 {\True $y'=\dfrac{-2x^2+4x+1}{(1-x)^2}$}
	 \loigiai{
		 $y=\dfrac{x(1+2x)}{1-x}=\dfrac{x+2x^2}{1-x}$ \\
	$y'=\dfrac{(x+2x^2)'(1-x)-(x+2x^2)(1-x)'}{(1-x)^2}=\dfrac{(1+4x)(1-x)-(x+2x^2)(-1)}{(1-x)^2}=\dfrac{-2x^2+4x+1}{(1-x)^2}$.} 
\end{ex}

% nguồn: Ok_TracNghiemLop11_Lan1
\begin{ex}%[1D7H2-1]%[ma:c9aad5b9]
	Cho hàm số  $y=\dfrac{x+2}{x-1}$. Tính $y'\left(3\right)$
	\choice 
		{$\dfrac{5}{2}$} 
		{\True $-\dfrac{3}{4}$} 
		 {$-\dfrac{3}{2}$}
		{$\dfrac{3}{4}$} 
	\loigiai{ 
	Ta có  $y=\dfrac{x+2}{x-1}\Rightarrow y'=\dfrac{-3}{\left(x-1\right)^2}\Rightarrow y'\left(3\right)=\dfrac{-3}{\left(3-1\right)^2}=-\dfrac{3}{4}$.}
\end{ex}

% nguồn: Ok_TracNghiemLop11_Lan1
\begin{ex}%[1D7H2-1]%[ma:3cffe898]
	Tính  đạo hàm của hàm số $y=\sqrt{x}+x$ tại điểm $x_0=4$ là:
	\choice 
		 {$y'(4)=\dfrac{9}{2}$}
		 {$y'(4)=6$}
		 {$y'(4)=\dfrac{3}{2}$}
		{\True  $y'(4)=\dfrac{5}{4}$}
	 \loigiai{
	Ta có $y'=\dfrac{1}{2\sqrt{x}}+1$ $\Rightarrow{y}'(4)=\dfrac{1}{2\sqrt{4}}+1=\dfrac{5}{4}$.} 
\end{ex}

% nguồn: Ok_TracNghiemLop11_Lan1
\begin{ex}%[1D7H2-1]%[ma:78ec6056]
	Cho hàm số $f(x)=\dfrac{x^3}{3}-\dfrac{3}{2}x^2-4x + 6$. Tìm nghiệm của  phương trình $f'(x)=0$.
	 \choice
		 {\True $x=-1$, $x=4$}
		{$x=1$,  $x=4$}
		{$x=0$, $x=3$} 
		{$x=-1$} 
	\loigiai{ 
		Ta có $f'(x)=0\Leftrightarrow x^2-3x-4=0\Leftrightarrow  \hoac{&x=-1\\&x=4}$.
	 }
\end{ex}

% nguồn: Ok_TracNghiemLop11_Lan1
\begin{ex}%[1D7H2-1]%[ma:90081d10]
	Tính đạo hàm của hàm số  $ f(x)=\dfrac{2x}{x-1}$ tại điểm $ x=-1$.
	 \choice
		{$ f'\left(-1\right)=1$} 
		{\True $  f'\left(-1\right)=-\dfrac{1}{2}$}
		 {$ f'\left(-1\right)=-2$}
		 {$ f'\left(-1\right)=0$}
	\loigiai{ 
	Ta có  $ f'(x)=\dfrac{-2}{\left(x-1\right)^2}$. Vậy $ f'\left(-1\right)=-\dfrac{1}{2}$.}
\end{ex}

% nguồn: Ok_TracNghiemLop11_Lan1
\begin{ex}%[1D7H2-1]%[ma:ce1646cf]
	Cho hàm số $ f(x)$ xác định trên $\mathbb{R}$  bởi $ f(x)=2x^2+1$. Giá trị $f'\left(-1\right)$ bằng
	 \choice
		{$2$} 
		 {$6$}
		 {\True $-4$}
		{$3$} 
	 \loigiai{
	Ta có $  f'(x)=4x$ $\Rightarrow{f}'\left(-1\right)=-4$.}
\end{ex}

% nguồn: Ok_Toan11_TracNghiem_Lan2
\begin{ex}%[1D7V2-1]%[ma:9ca65d62]
	Tìm tất cả các giá trị thực của tham số $m$ để hàm  số $y=\log \left(x^2-2x-m+1\right)$ có tập xác định là $\mathbb{R}$.
	 \choice
		{$m\le  2$}
		{$m  > 2$}
		{$m\ge  0$}
		{\True $m  < 0$}
	\loigiai{ 
		Để hàm số có tâp xác định $\mathbb{R}$ khi và chỉ khi $x^2-2x-m+1> 0$, $\forall x\in \mathbb{R}$. 
		\[\Leftrightarrow \Delta ' <  0\Leftrightarrow \left(-1\right)^2-1\cdot\left(-m+1\right) < 0\Leftrightarrow m < 0.\]
	} 
\end{ex}

% nguồn: Ok_Toan11_TracNghiem_Lan2
\begin{ex}%[1D7V2-1]%[ma:f11724a7]
	Cho hàm số$ f(x)=\dfrac{3x+1}{\sqrt{x^2+4}}$. Tính giá trị biểu  thức $f'(0)$.
	\choice 
		{$-3$} 
		 {$-2$}
		{\True $\dfrac{3}{2}$} 
		{$3$} 
	\loigiai{ 
		Ta có \begin{align*}
		f'(x)&=\dfrac{3\sqrt{x^2+4}-(3x+1)\cdot\dfrac{x}{\sqrt{x^2+4}}}{\left(\sqrt{x^2+4}\right)^2}\\ 
		 &=\dfrac{\dfrac{3\left(x^2+4\right)-(3x+1)x}{\sqrt{x^2+4}}}{\left(\sqrt{x^2+4}\right)^2}\\
		 &=\dfrac{12-x}{\left(\sqrt{x^2+4}\right)^3}.
	\end{align*}
 Vậy $f'(0)=\dfrac{3}{2}$.}
\end{ex}

% nguồn: Ok_Toan11_TracNghiem_Lan2
\begin{ex}%[1D7V2-1]%[ma:d651bde6]
	Cho  hàm số $f(x)=\dfrac{x-1}{\sqrt{x^2+1}}$. Biết $f'(x)=\dfrac{ax+b}{\left(x^2+1\right)\sqrt{x^2+1}}$, $\forall x$. Tính $S=2a+b$.
	 \choice
		 {\True $S=3$}
		 {$S=1$}
		 {$S=-1$}
		 {$S=5$}
	\loigiai{Ta  có
		\begin{align*}f'(x)&=\left(\dfrac{x-1}{\sqrt{x^2+1}}\right)'=\dfrac{(x-1)'\sqrt{x^2+1}-(x-1)\left(\sqrt{x^2+1}\right)'}{x^2+1}\\
			&=\dfrac{\sqrt{x^2+1}-(x-1)\cdot\dfrac{x}{\sqrt{x^2+1}}}{x^2+1}\\ 
			 &=\dfrac{x^2+1-x^2+x}{\left(x^2+1\right)\sqrt{x^2+1}}\\
			 &=\dfrac{x+1}{\left(x^2+1\right)\sqrt{x^2+1}}.
		\end{align*}
		Suy ra $a=1$, $b=1$.\\ 
		 Vậy $S=2a+b=3$.
	 }
\end{ex}

% nguồn: Ok_Toan11_TracNghiem_Lan2
\begin{ex}%[1D7V2-1]%[ma:a3e24cab]
	Có bao nhiêu giá trị nguyên của $m$ để hàm số  $y=\dfrac{x+2}{x+5m}$ có đạo hàm dương trên khoảng $(-\infty;-10)$?
	 \choice
		 {$1$}
		 {\True $2$}
		{$3$} 
		 {vô số}
	\loigiai{ 
		Tập xác định: $\mathscr{D}=(-\infty;-5m)\cup(-5m;+\infty)$.\\ 
		Ta có $y'=\dfrac{5m-2}{(x+5m)^2}$, $\forall x\neq -5m$.\\ 
		Hàm số $y=\dfrac{x+2}{x+5m}$ có đạo hàm dương  trên khoảng $(-\infty;-10)$ khi và chỉ khi
		\begin{align*}
			\heva{&\dfrac{5m-2}{(x+5m)^2}>0\\ &-10\leq-5m}\Leftrightarrow\heva{& 5m-2>0\\ &m\leq 2}\Leftrightarrow\heva{&  m>\dfrac{2}{5}\\ &m\leq 2}\Leftrightarrow\dfrac{2}{5}<m\leq 2.
		\end{align*}
		 Vì $m\in\mathbb{Z}$ nên $m\in\left\{1;2\right\}$.\\
		Vậy có $2$ giá trị nguyên của $m$ thỏa mãn  yêu cầu bài toán.
	} 
\end{ex}

% nguồn: Ok_Toan11_TracNghiem_Lan2
\begin{ex}%[1D7V2-1]%[ma:13bbd48f]
	Cho hàm số $f(x)=\dfrac{mx^3}{3}-\dfrac{mx^2}{2}+(3-m)x-2$. Tìm $m$  để $f'(x)>0$, $\forall x\in\mathbb{R}$.
	 \choice
		{$0\leq m\leq\dfrac{12}{5}$} 
		 {$0<m<\dfrac{12}{5}$}
		{\True $0\leq m<\dfrac{12}{5}$} 
		{$0<m\leq\dfrac{12}{5}$} 
	 \loigiai{
		Ta  có $f'(x)=mx^2-mx+(3-m)$.
		\begin{itemize}
			\item Nếu $m=0$ thì $f'(x)=3>0$,  $\forall x\in\mathbb{R}$.
			\item Nếu $m\neq 0$ thì $f'(x)=mx^2-mx+(3-m)$ là tam thức bậc hai.  Do đó,
			\begin{align*}
				f'(x)>0,~\forall x\in\mathbb{R}\Leftrightarrow\heva{& m>0\\ &\Delta=m^2-4m(3-m)<0}\Leftrightarrow\heva{&  m>0\\ &5m^2-12m<0}\Leftrightarrow 0<m<\dfrac{12}{5}.
			\end{align*}
		\end{itemize}
		Vậy $0\leq m<\dfrac{12}{5}$. 
	} 
\end{ex}

% nguồn: Ok_Toan11_TracNghiem_Lan2
\begin{ex}%[1D7V2-1]%[ma:d7b0fdad]
	Cho hàm số $f(x)=\ln 2\,021+\ln \left(\dfrac{x}{x+1} \right)$. Tính giá trị biểu thức $S=f'(1)+f'(2)+\cdots+f'(2\,020)$, tổng gồm  $2\,020$ số hạng.
	\choice 
		 {$\dfrac{2\,021}{2\,020}$}
		 {\True $\dfrac{2\,020}{2\,021}$}
		{$\dfrac{2\,021}{2022}$} 
		{$\dfrac{2022}{2\,021}$} 
	 \loigiai{
		Ta có $f'(x)=\dfrac{\left(\dfrac{x}{x+1} \right)'}{\dfrac{x}{x+1}}$ $=\dfrac{\dfrac{1}{\left(x+1\right)^2}}{\dfrac{x}{x+1}}$  $=\dfrac{1}{x\cdot\left(x+1\right)}=\dfrac{1}{x}-\dfrac{1}{x+1}$
		 Suy ra $f'(1)=1-\dfrac{1}{2}$.
		\allowdisplaybreaks 
		\begin{itemize}
			\item $ f'(2)=\dfrac{1}{2}-\dfrac{1}{3}$. 
			\item  $ f'(3)=\dfrac{1}{3}-\dfrac{1}{4}$.
			 \item \dots \dots
			 \item $ f'(2\,020)=\dfrac{1}{2\,020}-\dfrac{1}{2\,021}$.
		\end{itemize}
		Vậy $S=f'(1)+f'(2)+\cdots+f'(2\,020)=1-\dfrac{1}{2\,021}=\dfrac{2\,020}{2\,021}$. 
	} 
\end{ex}

% nguồn: Ok_Toan11_TracNghiem_Lan2
\begin{ex}%[1D7V2-1]%[ma:80b675eb]
	Cho hàm số $f(x)=\sqrt{-5x^2+14x-9}$. Tập hợp các giá  trị của $x$ để $f'(x)<0$ là
	\choice 
		{$\left(\dfrac{7}{5};+\infty\right)$} 
		{$\left(-\infty;\dfrac{7}{5}\right)$} 
		 {\True $\left(\dfrac{7}{5};\dfrac{9}{5}\right)$}
		 {$\left(1;\dfrac{7}{5}\right)$}
	\loigiai{ 
		Tập xác định: $\mathscr{D}=\left[1;\dfrac{9}{5}\right]$.\\ 
		Ta có $f(x)=\sqrt{-5x^2+14x-9}$ nên $f'(x)=\dfrac{-5x+7}{\sqrt{-5x^2+14x-9}}$, $\forall x\in\left(1;\dfrac{9}{5}\right)$.  Do đó,
		\begin{align*}
			f'(x)<0\Leftrightarrow\dfrac{-5x+7}{\sqrt{-5x^2+14x-9}}<0\Leftrightarrow\heva{&-5x+7<0\\ &1<x<\dfrac{9}{5}}\Leftrightarrow\heva{&x>\dfrac{7}{5}\\ &1<x<\dfrac{9}{5}}\Leftrightarrow\dfrac{7}{5}<x<\dfrac{9}{5}. 
		\end{align*}
	} 
\end{ex}

% nguồn: Ok_Toan11_TracNghiem_Lan2
\begin{ex}%[1D7V2-1]%[ma:21b5af28]
	Cho hàm số  $f(x)=\dfrac{x}{(x-1)(x-2)\cdots(x-2022)}$. Tính $f'(0)$.
	 \choice
		{$\dfrac{1}{2023}$} 
		 {$\dfrac{1}{2022}$}
		{$\dfrac{1}{2023!}$} 
		 {\True $\dfrac{1}{2022!}$}
	\loigiai{ 
		Đặt $g(x)=(x-1)(x-2)\cdots(x-2022)$. Khi đó, $f(x)=\dfrac{x}{g(x)}$.  Do đó,
		\begin{align*}
			 f'(x)=\dfrac{x'\cdot g(x)-x\cdot g'(x)}{g^2(x)}=\dfrac{g(x)-x\cdot g'(x)}{g^2(x)}=\dfrac{1}{g(x)}-\dfrac{x\cdot g'(x)}{g^2(x)}.
		\end{align*}
		Vậy $f'(0)=\dfrac{1}{g(0)}-\dfrac{0\cdot g'(0)}{g^2(0)}=\dfrac{1}{g(0)}=\dfrac{1}{(-1)\cdot  (-2)\cdots(-2022)}=\dfrac{1}{2022!}$.
	 }
\end{ex}

% nguồn: Ok_Toan11_TracNghiem_Lan2
\begin{ex}%[1D7V2-1]%[ma:8b581d78]
	Cho hàm số $f(x)=\ln \left(1-\dfrac{1}{x^2} \right)$. Biết rằng $T=f'(2)+f'(3)+\cdots+f'(2\,019)+f'(2\,020)=\dfrac{a}{b}$ với $a$, $b$ là các  số nguyên dương nguyên tố cùng nhau. Giá trị của $2a-b$ bằng
	\choice 
		{$2$} 
		 {$4$}
		 {\True $-2$}
		 {$-4$}
	\loigiai{ 
		 Ta có $f'(x)=\dfrac{2}{x(x-1)(x+1)}=\dfrac{1}{x(x-1)}-\dfrac{1}{x(x+1)}$.
		Khi đó  \allowdisplaybreaks
		\begin{eqnarray*}
			 T&=&f'(2)+f'(3)+\cdots+f'(2\,019)+f'(2\,020)\\
			&=&\dfrac{1}{1\cdot 2}-\dfrac{1}{2\cdot 3}+\dfrac{1}{2\cdot 3}-\dfrac{1}{3\cdot 4}+\cdots+\dfrac{1}{2\,018\cdot 2\,019}-\dfrac{1}{2\,019\cdot 2\,020}+\dfrac{1}{2\,019\cdot  2\,020}-\dfrac{1}{2\,020\cdot 2\,021}\\
			&=&\dfrac{1}{2}-\dfrac{1}{2\,020\cdot  2\,021}\\
			&=&\dfrac{1\,010\cdot 2\,021-1}{2\,020\cdot 2\,021}. 
		\end{eqnarray*}
		Nên $a=1\,010\cdot 2\,021-1$, $b=2\,020\cdot 2\,021\Rightarrow 2a-b=-2$. 
	 }
\end{ex}

% nguồn: Ok_Toan11_TracNghiem_Lan2
\begin{ex}%[1D7V2-1]%[ma:bbe3f787]
	Cho hàm số $f(x)=\ln \left(\dfrac{x}{x+2} \right)$. Tổng $S=f' (1)+f' (3)+f' (5)+\cdots+f'  (2\,021)$ bằng
	\choice 
		{$\dfrac{4\,035}{2\,021}$} 
		 {$\dfrac{2\,021}{2\,022}$}
		{$2\,021$} 
		 {\True $\dfrac{2\,022}{2\,023}$}
	 \loigiai{
		Ta có $f(x)=\ln \left(\dfrac{x}{x+2} \right)\Rightarrow  f' (x)=\dfrac{2}{x\left(x+2\right)}=\dfrac{1}{x}-\dfrac{1}{x+2}$.\\
		 Vậy \allowdisplaybreaks
		\begin{eqnarray*}
			S&=&f' (1)+f' (3)+f'  (5)+\cdots+f' (2\,021)\\
			 &=&\dfrac{1}{1}-\dfrac{1}{3}+\dfrac{1}{3}-\dfrac{1}{5}+\cdots\ldots+\dfrac{1}{2\,021}-\dfrac{1}{2023}\\
			&=&1-\dfrac{1}{2023}\\ 
			 &=&\dfrac{2022}{2023}.
		\end{eqnarray*}
	} 
\end{ex}

% nguồn: Ok_Toan11_TracNghiem_Lan2
\begin{ex}%[1D7V2-1]%[ma:71e2cc33]
	 Cho hàm số $f(x)=\ln \dfrac{2\,018x}{x+1}$. Tính tổng $S=f'(1)+f'(2)+\cdots+f'(2\,018)$.
	\choice 
		 {$\ln 2\,018$}
		 {$1$}
		 {$2\,018$}
		{\True  $\dfrac{2\,018}{2\,019}$}
	\loigiai{ 
		Ta có $f'(x)=\left(\ln \dfrac{2\,018x}{x+1}  \right)'=\dfrac{1}{\dfrac{2\,018x}{x+1}}.\left(\dfrac{2\,018x}{x+1}.\right)'$ $=\dfrac{x+1}{2\,018x}.\dfrac{2\,018}{\left(x+1\right)^2}=\dfrac{1}{x.\left(x+1\right)}$.
		 Vậy \allowdisplaybreaks
		\begin{eqnarray*}
			S&=&f'(1)+f'(2)+\cdots+f'(2\,018)\\ 
			&=&\dfrac{1}{1.2}+\dfrac{1}{2.3}+\cdots+\dfrac{1}{2\,018.2\,019}\\ 
			 &=&\dfrac{1}{1}-\dfrac{1}{2}+\dfrac{1}{2}-\dfrac{1}{3}+\cdots+\dfrac{1}{2\,018}-\dfrac{1}{2\,019}\\
			 &=&1-\dfrac{1}{2\,019}\\
			&=&\dfrac{2\,018}{2\,019}. 
		\end{eqnarray*}
	} 
\end{ex}

% nguồn: Ok_Toan11_TracNghiem_Lan2
\begin{ex}%[1D7V2-1]%[ma:9d993fc5]
	Cho hàm số $y=(m+2)x^3+\dfrac{3}{2}(m+2)x^2+3x-1$, với $m$ là tham  số. Số các giá trị nguyên $m$ để $y'\geq0$, $\forall x\in\mathbb{R}$ là
	\choice 
		{\True  $5$}
		 {Vô số}
		{$3$} 
		 {$4$}
	 \loigiai{
		Ta có $y'=3(m+2)x^2+3(m+2)x+3$.  Do đó,
		\begin{align*}
			 y'\geq 0\Leftrightarrow3(m+2)x^2+3(m+2)x+3\geq 0\Leftrightarrow(m+2)x^2+(m+2)x+1\geq 0. \tag{1}
		\end{align*}
		\begin{enumerate}[{\bfseries TH1:}]
			\item $m+2=0\Leftrightarrow m=-2$. Khi đó, $y'=1\geq 0$, $\forall x\in\mathbb{R}$. 
			\item $m+2\neq 0\Leftrightarrow m\neq-2$.\\ 
			Phương trình $(1)$ luôn thỏa mãn  $\forall x\in\mathbb{R}$ khi và chỉ khi\allowdisplaybreaks
			\begin{align*}
				\heva{& m+2>0\\  &\Delta\leq 0}&\Leftrightarrow\heva{&m>-2\\ &(m+2)^2-4(m+2)\leq 0}\\
				&\Leftrightarrow\heva{&m>-2\\  &(m+2)(m-2)\leq 0}\\
				&\Leftrightarrow\heva{&m>-2\\  &-2\leq m\leq 2}\\
				&\Leftrightarrow-2<m\leq 2. 
			\end{align*}
		\end{enumerate}
		Kết hợp hai trường hợp, ta được  $m\in\left\{-2;-1;0;1;2\right\}$.
	} 
\end{ex}

% nguồn: Ok_Toan11_TracNghiem_Lan2
\begin{ex}%[1D7V2-1]%[ma:7d62b8c3]
	Cho hàm số $y=f(x)=(2m-1)\mathrm{e}^x+3$.  Giá trị của $m$ để $f'\left(-\ln 3\right)=\dfrac{5}{3}$ là
	\choice 
		 {$m=\dfrac{7}{9}$}
		{$m=\dfrac{2}{9}$} 
		{\True  $m=3$}
		 {$m=-\dfrac{3}{2}$}
	 \loigiai{
		 Ta có $f'(x)=(2m-1)\mathrm{e}^x$. Suy ra $f'(-\ln 3)=(2m-1)\mathrm{e}^{-\ln 3}=\dfrac{2m-1}{\mathrm{e}^{\ln 3}}=\dfrac{2m-1}{3}$.\\
		Mà $f'(-\ln  3)=\dfrac{5}{3}$ nên $\dfrac{2m-1}{3}=\dfrac{5}{3}\Leftrightarrow m=3$.
	 }
\end{ex}

% nguồn: Ok_Toan11_TracNghiem_Lan2
\begin{ex}%[1D7V2-1]%[ma:fb0b0593]
	Cho hàm số $f(x)=\sqrt{x^2-2x}$. Tập nghiệm $S$  của bất phương trình $f'(x)\geq f(x)$ có bao nhiêu giá trị nguyên?
	 \choice
		 {$1$}
		{$2$} 
		 {\True $0$}
		{$3$} 
	 \loigiai{
		Tập xác định của hàm số là $\mathscr{D}=(-\infty;0]\cup[2;+\infty)$.\\ 
		Ta có $f'(x)=\dfrac{(x^2-2x)’}{2\sqrt{x^2-2x}}=\dfrac{x-1}{\sqrt{x^2-2x}}$, $\forall  x\in(-\infty;0)\cup(2;+\infty)$.\\
		Do đó, với $x\in(-\infty;0)\cup(2;+\infty)$,  ta có\allowdisplaybreaks
		\begin{align*}
			f'(x)\geq  f(x)&\Leftrightarrow\dfrac{x-1}{\sqrt{x^2-2x}}\geq\sqrt{x^2-2x}\\
			&\Leftrightarrow  x-1\geq x^2-2x\\
			&\Leftrightarrow x^2-3x+1\leq0\\ 
			&\Leftrightarrow\dfrac{3-\sqrt{5}}{2}\leq  x\leq\dfrac{3+\sqrt{5}}{2}.
		\end{align*}
		Kết hợp với điều kiện $x\in(-\infty;0)\cup(2;+\infty)$, ta được  $S=\left(2;\dfrac{3+\sqrt{5}}{2}\right]$.\\
		Ta có $\dfrac{3+\sqrt{5}}{2}\approx 2{,}6$. Do đó, trên khoảng $\left(2;\dfrac{3+\sqrt{5}}{2}\right]$ không có số nguyên  nào.\\
		Vậy tập nghiệm $S$ của bất phương trình $f'(x)\geq f(x)$ không có  giá trị nguyên nào.
	 }
\end{ex}

% nguồn: Ok_Toan11_TracNghiem_Lan2
\begin{ex}%[1D7V2-1]%[ma:f46cd2bd]
	 Cho hàm số $f(x)=-x^3+3mx^2-12x+3$, với $m$ là tham số thực. Số giá trị nguyên của $m$ để $f'(x)\leq0$ với $\forall x\in\mathbb{R}$ là
	\choice 
		 {$1$}
		 {\True $5$}
		{$4$} 
		 {$3$}
	\loigiai{ 
		Ta có $f(x)=-x^3+3mx^2-12x+3$  nên $f'(x)=-3x^2+6mx-12$. Do đó,\\
		$f'(x)\leq0$, $\forall x\in\mathbb{R}$ khi và chỉ khi 
		\begin{align*}
			-3x^2+6mx-12\leq0,~\forall x\in\mathbb{R}\Leftrightarrow\heva{&a<0\\ &\Delta'\leq 0}\Leftrightarrow\heva{&-3<0\\  &9m^2-36\leq 0}\Leftrightarrow-2\leq m\leq 2.
		\end{align*}
		Vì $m\in\mathbb{Z}$ nên  $m\in\left\{-2;-1;0;1;2\right\}$.\\
		 Vậy có $5$ giá trị nguyên $m$ thỏa mãn.
	 }
\end{ex}

% nguồn: Ok_Toan11_TracNghiem_Lan2
\begin{ex}%[1D7V2-1]%[ma:f6806f61]
	Tìm tất cả các giá trị thực của tham số $m$ để hàm số $y=\ln \left(x^2-2x+m+1\right)$ có tập xác  định là $\mathbb{R}$.
	 \choice
		{$0< m <  3$}
		{$m  <-1$ hoặc $m > 0$}
		 {\True $m > 0$}
		 {$m=0$}
	\loigiai{ 
		Hàm số có tâp xác định $\mathbb{R}$ khi và chỉ khi  $x^2-2x+m+1> 0,\forall x\in \mathbb{R}\Leftrightarrow \heva{&a=1> 0\text{ luôn đúng} \\&\Delta '=1-(1+m)< 0}\Leftrightarrow m > 0$.
	} 
\end{ex}

% nguồn: Ok_Toan11_TracNghiem_Lan2
\begin{ex}%[1D7V2-1]%[ma:fce336c5]
	Cho hàm số $f(x)=\ln \dfrac{x+1}{x+4}$. Tính  giá trị của biểu thức $P=f'(0)+f'(3)+f'(6)+\cdots+f'(2\,019)$.
	\choice 
		{$\dfrac{1}{4}$} 
		 {$\dfrac{2\,024}{2\,023}$}
		 {\True $\dfrac{2\,022}{2\,023}$}
		 {$\dfrac{2\,020}{2\,023}$}
	\loigiai{ 
		Với $x\in [0;+\infty)$ ta có $x+1> 0$ và $x+4> 0$ nên  $f(x)=\ln \dfrac{x+1}{x+4}=\ln \left(x+1\right)-\ln \left(x+4\right)$.\\
		 Từ đó $f'(x)=\dfrac{1}{x+1}-\dfrac{1}{x+4}$.\\
		 Do đó \allowdisplaybreaks
		\begin{eqnarray*}
			 P&=&f'(0)+f'(3)+f'(6)+\cdots+f'(2\,019)\\
			&=&\left(1-\dfrac{1}{4} \right)+\left(\dfrac{1}{4}-\dfrac{1}{7} \right)+\left(\dfrac{1}{7}-\dfrac{1}{10} \right)+\cdots+\left(\dfrac{1}{2\,020}-\dfrac{1}{2023}  \right)\\
			&=&1-\dfrac{1}{2023}\\ 
			&=&\dfrac{2022}{2023}. 
		\end{eqnarray*}
	 }
\end{ex}

% nguồn: Ok_Toan11_TracNghiem_Lan2
\begin{ex}%[1D7V2-1]%[ma:bd657198]
	Cho hàm số $f(x)=\ln 2\,018+\ln  \left(\dfrac{x}{x+1} \right)$. Tính $S=f'(1)+f'(2)+f'(3)+\cdots+f'(2\,017)$.
	\choice 
		 {$S=\dfrac{4035}{2\,018}$}
		{\True  $S=\dfrac{2\,017}{2\,018}$}
		{$S=\dfrac{2016}{2\,017}$} 
		{$S=2\,017$} 
	 \loigiai{
		Ta  có $f(x)=\ln 2\,018+\ln \left(\dfrac{x}{x+1} \right)\Rightarrow f'(x)=\dfrac{1}{x\left(x+1\right)}=\dfrac{1}{x}-\dfrac{1}{x+1}$.\\
		Do  đó $S=\dfrac{1}{1}-\dfrac{1}{2}+\dfrac{1}{2}-\dfrac{1}{3}+\cdots+\dfrac{1}{2\,017}-\dfrac{1}{2\,018}=1-\dfrac{1}{2\,018}=\dfrac{2\,017}{2\,018}$.
	 }
\end{ex}

% nguồn: Ok_Toan11_TracNghiem_Lan2
\begin{ex}%[1D7V2-1]%[ma:650dd5a0]
	Đạo hàm của  hàm số $y=\dfrac{x+3}{\sqrt{x^2+1}}$ là
	\choice 
		{\True  $\dfrac{1-3x}{\left(x^2+1\right)\sqrt{x^2+1}}$}
		 {$\dfrac{1+3x}{\left(x^2+1\right)\sqrt{x^2+1}}$}
		{$\dfrac{1-3x}{x^2+1}$} 
		 {$\dfrac{2x^2-x-1}{\left(x^2+1\right)\sqrt{x^2+1}}$}
	\loigiai{Ta  có $y'=\dfrac{\sqrt{x^2+1}-\dfrac{(x+3)x}{\sqrt{x^2+1}}}{x^2+1}=\dfrac{1-3x}{\left(x^2+1\right)\sqrt{x^2+1}}$.
	 }
\end{ex}

% nguồn: Ok_TracNghiemLop11_Lan1
\begin{ex}%[1D7V2-1]%[ma:57849cde]
	Cho $\left[(2x-1)^2\cdot(2-3x)\right]'=ax^2+bx+c$. Tính $S=a+b+c$. 
	\choice 
		 {\True $S=-7$}
		 {$S=-87$}
		 {$S=-47$}
		{$S=17$} 
	 \loigiai{
		 Ta có:
		\begin{align*}
			&\left[(2x-1)^2\cdot(2-3x)\right]'=  \left[(2x-1)^2\right]'\cdot(2-3x)+(2x-1)^2\cdot \left(2-3x\right)'\\
			=& 4(2x-1)\cdot (2-3x)+(2x-1)^2\cdot (-3)=-36x^2+40x-11. 
		\end{align*}
		Do đó $a=-36$; $b=40$;  $c=-11$.\\
		Vậy $S=a+b+c=(-36)+40+(-11)=-7$. 
	 }
\end{ex}

% nguồn: Ok_TracNghiemLop11_Lan1
\begin{ex}%[1D7Y2-1]%[ma:b1d41e6b]
	Hàm số $y = \sin \left( \dfrac{\pi  }{2} - 2x \right)$ có đạo hàm bằng
	\choice 
		{$y'  = 2\sin 2x$}
		{\True $y' = - 2\sin 2x$} 
		{$y'  = \cos \left( \dfrac{\pi }{2} - 2x \right)$}
		{$y' = 2\cos \left( \dfrac{\pi }{2}  - 2x \right)$}
	 \loigiai{
		Ta có $y=\cos 2x$ nên $y'=-2\sin  2x$.
	 }
\end{ex}

% nguồn: Ok_TracNghiemLop11_Lan1
\begin{ex}%[1D7Y2-1]%[ma:dba6f517]
	Đạo hàm của hàm số $y = \dfrac{x^2 - 2x - 1}{x - 2}$ bằng 
	\choice 
		{\True $y' = \dfrac{x^2 - 4x + 5}{(x - 2)^2}$} 
		{$y' = \dfrac{x^2 - 6x +  5}{(x - 2)^2}$}
		{$y' = \dfrac{x^2 - 6x  + 4}{(x - 2)^2}$}
		{$y' = \dfrac{x^2 -  6x - 1}{(x - 2)^2}$}
	 \loigiai{
		Ta  có
		$$y'=\dfrac{(2x-2)(x-2)-(x^2-2x-1)}{(x-2)^2}=\dfrac{x^2-4x+5}{(x-2)^2}.$$ 
	} 
\end{ex}

% nguồn: Ok_TracNghiemLop11_Lan1
\begin{ex}%[1D7B2-1]%[ma:ba897310]
	Đạo hàm cấp hai  của hàm số $y = \tan x$ là
	\choice 
		{$y'' = 2\tan x \left(1 - \tan^2 x  \right)$}
		{\True $y'' = 2\tan x  \left(1 + \tan^2 x \right)$}
		{$y'' = - 2\tan x \left(1  - \tan^2 x \right)$}
		{$y'' = - 2\tan \left(1 + \tan^2 x \right)$} 
	\loigiai{ 
		$y'=\tan^2 x + 1$;\hspace{1cm} $y'' = 2\tan  x \cdot (\tan x)' = 2\tan x \cdot (\tan^2 x + 1)$.
	 }
\end{ex}

% nguồn: Ok_TracNghiemLop11_Lan1
\begin{ex}%[1D7Y2-1]%[ma:4b10dd80]
	Đạo hàm của hàm số $y = x\sin  x$ bằng
	 \choice
		{\True $y' = \sin x + x\cos  x$}
		{$y' = \sin x  - x\cos x$}
		 {$y' = x\cos x$}
		{$y' =  - x\cos x$}
	\loigiai{ 
		Ta  có $y'=\sin x+x\cos x.$
	 }
\end{ex}


%%%%%%%%%%%%%%%% ĐÚNG/SAI %%%%%%%%%%%%%%%%

% nguồn: Ok_dungsai_lop11_Lan1
\begin{ex}%[1D7N2-1]%[ma:411505f7]
 Cho bốn hàm số $f(x)=\ln x$, $g(x)=\mathrm{e}^x$,  $h(x)=2024^x$ và $k(x)=\log_3x$. Xét tính đúng sai của các mệnh đề sau
  \choiceTF
  {$f'(x)=-\dfrac{1}{x}, \forall  x>0$}
  {$g'(x)=x\cdot\mathrm{e}^{x-1} ,\forall x\in\mathbb{R}$} 
  {$h'(x)=2024^x, \forall x\in\mathbb{R}$} 
  {\True $k'(x)=\dfrac{1}{x\ln  3}, \forall x>0$}
 \loigiai{ 
  \begin{itemchoice}
   \itemch Sai vì $f'(x)=\dfrac{1}{x},  \forall x>0$.
   \itemch Sai vì $g'(x)=\mathrm{e}^x, \forall x\in\mathbb{R}$. 
   \itemch Sai  vì $h'(x)=2024^x\cdot\ln 2024, \forall x\in\mathbb{R}$.
   \itemch Đúng vì $k'(x)=\dfrac{1}{x\ln 3}, \forall  x>0$.
  \end{itemchoice}
  }
\end{ex}

% nguồn: Ok_dungsai_lop11_Lan1
\begin{ex}%[1D7N2-1]%[ma:08214571]
 Cho hàm số $f(x)=\dfrac{4}{5}x^5-6$. 
  \choiceTF
  {\True $f'(x)=4x^4$} 
   {\True $f'(2)=64$}
   {$f(0)=6$}
  {$f'(x)=4\Leftrightarrow x=1$} 
 \loigiai{ 
 \begin{itemchoice}
  \itemch Đúng\\ 
  Ta có $f'(x)=4x^4$. 
  \itemch Sai\\ 
  Ta có  $f'(x)=4x^4\Rightarrow f'(2)=4\cdot 2^2$.
   \itemch Sai\\
   Ta có $f(0)=\dfrac{4}{5}\cdot 0^5-6=-6$.
  \itemch  Sai\\
  Ta có $f'(x)=4\Leftrightarrow 4x^4=4\Leftrightarrow x=\pm  1$.
 \end{itemchoice}
 } 
\end{ex}

% nguồn: Ok_dungsai_lop11_Lan1
\begin{ex}%[1D7N2-1]%[ma:6b853c91]
 Cho $u=u(x)$, $v=v(x)$, $w=w(x)$ là các hàm số có đạo hàm tại điểm $x$ thuộc khoảng xác định. Xét tính đúng sai của  các mệnh đề sau
 \choiceTF 
  {\True  $(u+v+w)'=u'+v'+w'$}
  {\True  $(u+v-w)'=u'+v'-w'$}
   {$(uv)'=u'v'$}
  {$\left(\dfrac{u}{v}\right)'=\dfrac{u'}{v'}$  với $v=v(x)\neq 0, v'=v'(x)\neq 0$}
 \loigiai{ 
  \begin{itemchoice}
    \itemch Đúng.
   \itemch  Đúng.
   \itemch Sai vì $(uv)'=u'v+uv'$. 
    \itemch Sai vì $\left(\dfrac{u}{v}\right)'=\dfrac{u'v-uv'}{v^2}$.
  \end{itemchoice}
  }
\end{ex}

% nguồn: Ok_dungsai_lop11_Lan1
\begin{ex}%[1D7N2-1]%[ma:e2eb2689]
 Cho bốn hàm số $f(x)=\sin x$, $g(x)=\cos  x$, $h(x)=\tan x$ và $k(x)=\cot x$. Xét tính đúng sai của các mệnh đề sau
 \choiceTF 
  {\True $f'(x)=\cos  x, \forall x\in\mathbb{R}$}
  {$g'(x)=\sin  x, \forall x\in\mathbb{R}$}
  {\True $h'(x)=\dfrac{1}{\cos^2x}, \forall x\neq  \dfrac{\pi}{2}+k\pi\left(k\in\mathbb{Z}\right)$}
   {$k'(x)=\dfrac{1}{\sin^2x}, \forall x\neq k\pi, \left(k\in\mathbb{Z}\right)$}
 \loigiai{ 
  \begin{itemchoice}
   \itemch Đúng vì $f'(x)=\cos x, \forall  x\in\mathbb{R}$.
    \itemch Sai vì $g'(x)=-\sin x, \forall x\in\mathbb{R}$.
   \itemch  Đúng vì $h'(x)=\dfrac{1}{\cos^2x}, \forall x\neq\dfrac{\pi}{2}+k\pi, \left(k\in\mathbb{Z}\right)$.
   \itemch Sai vì $k'(x)=-\dfrac{1}{\sin^2x}, \forall x\neq  k\pi, \left(k\in\mathbb{Z}\right)$.
  \end{itemchoice}
 } 
\end{ex}

% nguồn: Ok_dungsai_lop11_Lan1
\begin{ex}%[1D7N2-1]%[ma:96f3caaa]
 Xét tính đúng, sai của các mệnh đề  sau
  \choiceTF
  {\True Hàm số $y=x^2+x+1$ có đạo hàm  trên $\mathbb{R}$ là $y'=2x+1$}
  {Đạo hàm của hàm số  $y=x^3-2x^2+\dfrac{5}{x}+7$ là $y'=x^2-2x-\dfrac{5}{x^2}$}
  {\True Một hàm số có đạo hàm $y'=3x^2-\dfrac{1}{x^2}$  là $y=\dfrac{x^4+1}{x}$}
  {\True Một hàm số có đạo hàm bằng $\dfrac{1}{2\sqrt{x}}$ là  $f(x)=\sqrt{x}$}
 \loigiai{ 
  \begin{itemchoice}
   \itemch  Ta có $y'=(x^2+x+1)’=2x+1$.
   \itemch Tập xác định  $\mathscr{D}=\mathbb{R}\setminus \{0\}$.\\
    Ta có $y'=3x^2-4x-\dfrac{5}{x^2}$.
   \itemch Ta  có $y=\dfrac{x^4+1}{x}=x^3+\dfrac{1}{x}\Rightarrow y'=3x^2-\dfrac{1}{x^2}$.
    \itemch Ta có $(\sqrt{x})’=\dfrac{1}{2\sqrt{x}}$.
  \end{itemchoice}
 } 
\end{ex}

% nguồn: Ok_dungsai_lop11_Lan1
\begin{ex}%[1D7H2-1]%[ma:9d9bcbaa]
 Cho hàm của hàm số $f(x)=\dfrac{2x-1}{x+1}$. Xét tính đúng sai của các mệnh đề  sau
  \choiceTF
  {\True $f'(x)=\dfrac{3}{(x+1)^2}, \forall x\neq -1$} 
   {\True $f'(0)=3$}
  {$f'(x)<0\Leftrightarrow  x<-1$}
  {Phương trình $f'(x)=\dfrac{1}{3}$  có nghiệm duy nhất}
 \loigiai{ 
  Ta có $f'(x)=\dfrac{3}{(x+1)^2},  \forall x\neq -1$.
  \begin{itemchoice}
   \itemch  Đúng.
   \itemch  Đúng vì $f'(0)=\dfrac{3}{(0+1)^2}=3$.
   \itemch Sai  vì $f'(x)=\dfrac{3}{(x+1)^2}>0, \forall x\neq -1$.
   \itemch  Sai vì $f'(x)=\dfrac{1}{3}\Leftrightarrow (x+1)^2=9\Leftrightarrow\hoac{&x+1=3\\&x+1=-3}\Leftrightarrow\hoac{&x=2\\&x=-4.}$
  \end{itemchoice}
 } 
\end{ex}

% nguồn: Ok_dungsai_lop11_Lan1
\begin{ex}%[1D7H2-1]%[ma:af6b9b43]
 Cho hai hàm số $f(x)=x^{2024}$ và  $g(x)=3x$. Xét tính đúng sai của các mệnh đề sau
 \choiceTF 
   {$f'(x)=x^{2023}$}
   {\True $g'(x)=3$}
  {\True $\left[f(x)+g(x)\right]'=2024x^{2023}+3$} 
  {$\left[g(x)-2\right]'=1$} 
 \loigiai{ 
  \begin{itemchoice}
   \itemch Sai vì $f'\left(x\right)=2024x^{2023}$. 
   \itemch  Đúng vì $g'(x)=3$.
    \itemch Đúng vì $\left[f(x)+g(x)\right]'=f'(x)+g'(x)=2024x^{2023}+3$.
   \itemch  Sai vì $\left[g(x)-2\right]'=g'(x)-2'=3$.
  \end{itemchoice}
  }
\end{ex}

% nguồn: Ok_dungsai_lop11_Lan1
\begin{ex}%[1D7H2-1]%[ma:67a8c16c]
 Cho hàm của hàm số $f(x)=\sin 2x$. Xét tính đúng sai của các  mệnh đề sau
 \choiceTF 
   {$f'(x)=\cos 2x, \forall x\in\mathbb{R}$}
   {\True $f'\left(\dfrac{\pi}{2}\right)=-2$}
   {$f'\left(\pi\right)=1$}
   {\True $-2\leq f'(x)\leq 2, \forall x\in\mathbb{R}$}
 \loigiai{ 
  Ta có $f'(x)=2\cos 2x,  \forall x\in\mathbb{R}$.
  \begin{itemchoice}
   \itemch Sai. 
   \itemch Đúng vì $f'\left(\dfrac{\pi}{2}\right)2\cos\left(2\cdot\dfrac{\pi}{2}\right)=-2$. 
   \itemch  Sai vì $f'\left(\pi\right)=2\cos\left(2\pi\right)=2$.
   \itemch Đúng vì $-1\leq\cos 2x\leq 1\Leftrightarrow -2\leq 2\cos 2x\leq 2, \forall  x\in\mathbb{R}$.
  \end{itemchoice}
 } 
\end{ex}

% nguồn: Ok_dungsai_lop11_Lan1
\begin{ex}%[1D7H2-1]%[ma:86f4fcbb]
 Cho hàm số $y=f(x)=\dfrac{2x+1}{x-1}$. Xét tính đúng sai của các  mệnh đề sau
 \choiceTF 
  {\True $f'(x)=\dfrac{-3}{(x-1)^2}, \forall  x\in\mathbb{R}\backslash\left\{1\right\}$}
   {$f'(x)=\dfrac{3}{(x-1)^2}, \forall x\in\mathbb{R}\backslash\left\{1\right\}$}
   {\True $f'(2)=-3$}
  {\True $f'(0)=-3$} 
  \loigiai{
  Ta có $f'(x)=\dfrac{2(x-1)-1(2x+1)}{(x-1)^2} =\dfrac{-3}{(x-1)^2},  \forall x\in\mathbb{R}\backslash \left\{1\right\}$.
  \begin{itemchoice}
   \itemch  Đúng.
    \itemch Sai.
    \itemch Đúng vì $f'(2)=\dfrac{-3}{(2-1)^2}=-3$.
   \itemch  Đúng vì $f'(0)=\dfrac{-3}{(0-1)^2}=-3$.
  \end{itemchoice}
 } 
\end{ex}

% nguồn: Ok_dungsai_lop11_Lan1
\begin{ex}%[1D7H2-1]%[ma:eae736f0]
 Cho hàm số $y=-4 x^3+\dfrac{x^2}{2}-2 x+3$, biết $y'=a x^2+b x+c$.  Khi đó
 \choiceTF 
   {$a+b+c=-10$}
  {Phương trình $y'=0$ có hai  nghiệm phân biệt}
  {\True Đồ thị hàm số $y'$ cắt trục tung tại điểm  $(0;-2)$}
  {Đồ thị hàm số $y'$  cắt đường thẳng $y=3$ tại hai điểm phân biệt}
  \loigiai{
 \begin{enumerate}
  \item  Ta có
  \begin{align*}
    y'=\left(-4 x^3+\dfrac{x^2}{2}-2 x+3\right)'=-12x^2+x-2,
  \end{align*}
  đồng nhất với $y'=ax^2+bx+c$ ta được $a=-12$, $b=1$,  $c=-2$.\\
  Do đó 
  \begin{align*}
   a+b+c=-12+1-2=-13. 
  \end{align*}
  \item 
   Ta có
  \begin{align*}
   y'=0\Leftrightarrow  -12x^2+x-2=0,
  \end{align*}
  phương trình này có biệt thức 
  \begin{align*}
   \Delta =  1^2-4.(-12)(-2)=-47<0,
  \end{align*}
  do dó phương trình  này vô nghiệm.
  \item 
  Với $x=0$, ta có 
  \begin{align*}
    y'(0)=-12.0^2+x-2=-2,
  \end{align*}
  do đó, đồ thị của hàm số $y'$ đi qua điểm $(0;-2)$ tức cắt trục tung tại  điểm $(0;-2)$.
   \item
  Với $y'=3$ ta xét phương trình 
  \begin{align*}
   -12x^2+x-2=3\Leftrightarrow -12x^2+x-5=0\Leftrightarrow x\in  \varnothing.
  \end{align*}
  Phương  trình này vô nghiệm. Do đó, đồ thị hàm số của $y'$ cắt đường thẳng $y=3$ tại hai điểm phân biệt là sai.
 \end{enumerate}
  }
\end{ex}

% nguồn: Ok_dungsai_lop11_Lan1
\begin{ex}%[1D7H2-1]%[ma:d0d5d996]
 Cho hàm số  $y=\cos^2 x$. Xét tính đúng sai của các khẳng định sau.
  \choiceTF
   {$y(0)=0$}
  {\True $y'=-\sin  2x$}
  {$y'\left(\dfrac{\pi}{2}\right)=-1$} 
  {\True Phương  trình $y'=-1$ có nghiệm là $x=\dfrac{\pi}{4}+k\pi,\ k\in\mathbb{Z}$}
  \loigiai{
  Ta có $y'= \left(\cos^2 x\right)'=2\cos x\cdot (-\sin x)=-\sin  2x$.
  \begin{itemchoice}
    \itemch Sai vì $y(0)=\cos^2 0=1$.
   \itemch Đúng vì $y'=2.\cos x.(-\sin  x)=-\sin 2x$.
   \itemch Sai  vì $y'\left(\dfrac{\pi}{2}\right)=-\sin\left(2.\dfrac{\pi}{2}\right)=0$.
   \itemch Đúng vì phương trình $y'=-1\Leftrightarrow -\sin  2x=-1\Leftrightarrow x=\dfrac{\pi}{4}+k\pi,\ k\in\mathbb{Z}.$
  \end{itemchoice}
 } 
\end{ex}

% nguồn: Ok_dungsai_lop11_Lan1
\begin{ex}%[1D7H2-1]%[ma:fd2c0fbe]
 Cho hàm số $y=f(x)=2x^2+x+1$.  Xét tính đúng sai của các mệnh đề sau
  \choiceTF
   {\True $f'(x)=4x+1, \forall x\in\mathbb{R}$}
  {$f'(x)=0\Leftrightarrow x=\dfrac{1}{4}$} 
  {\True $f'(1)=5$} 
   {$f'(x)<-3\Leftrightarrow x>-1$}
 \loigiai{ 
  \begin{itemchoice}
   \itemch Đúng vì $f'(x)=4x+1,  \forall x\in\mathbb{R}$.
   \itemch Sai vì $f'(x)=0\Leftrightarrow  4x+1=0\Leftrightarrow x=-\dfrac{1}{4}$.
   \itemch Đúng vì $f'(1)=4\cdot  1+1=5$.
   \itemch  Sai vì $f'(x)<-3\Leftrightarrow 4x+1<-3\Leftrightarrow x<-1$.
  \end{itemchoice}
 } 
\end{ex}

% nguồn: Ok_dungsai_lop11_Lan1
\begin{ex}%[1D7H2-1]%[ma:06e7996a]
 Tính được đạo hàm của các hàm  số sau. Khi đó
 \choiceTF{\True $y=\log_2(9x-5)$ có $y'=\dfrac{9}{(9x-5)\ln 2}$}{\True $y=2\mathrm{e}^{3x+1}$  có $y'=6\mathrm{e}^{3x+1}$}{\True $y=3^{x^3-1}$ có $y'=3\ln 3 \cdot x^2\cdot3^{x^3-1}$}{$y=\ln \sqrt{x}$ có $y'=-\dfrac{1}{2x}$}
 \loigiai{ 
  \begin{enumerate}
   \item $y'=\dfrac{(9x-5)'}{(9x-5)\ln 2}=\dfrac{9}{(9x-5) \ln  2}$
   \item $y'=2(3x+1)' \cdot  \mathrm{e}^{3x+1}=6\mathrm{e}^{3x+1}$.
   \item  $y'=\left(x^3-1\right)' \cdot 3^{x^3-1} \cdot \ln 3=3\ln 3 \cdot x^2 \cdot 3^{x^3-1}$.
    \item $y'=\dfrac{(\sqrt{x})'}{\sqrt{x}}=\dfrac{1}{2 \sqrt{x} \sqrt{x}}=\dfrac{1}{2x}$
  \end{enumerate}
 } 
\end{ex}

% nguồn: Ok_dungsai_lop11_Lan1
\begin{ex}%[1D7H2-1]%[ma:098d04a5]
 Cho hàm số $f(x)=\sqrt{x^2-4}$ và $g(x)=3x-2$. Xét tính  đúng sai của các mệnh đề sau
 \choiceTF 
  {Giá trị  của hàm số $g(x)$ tại $x=2$ bằng $3$}
  {\True Đạo hàm của hàm số  $f(x)$ tại $x=3$ bằng $\dfrac{3\sqrt{5}}{5}$}
  {\True Đạo hàm của hàm số $f\left[g(x)\right]$ tại $x=3$ bằng  $\dfrac{7\sqrt{5}}{5}$}
  {Đạo hàm của hàm số $g\left[f(x)\right]$ tại $x=3$ bằng $\dfrac{3\sqrt{5}}{5}$} 
  \loigiai{
  \begin{itemchoice}
   \itemch Sai.  $g(2)=3\cdot 2-2=4$.
   \itemch  Đúng. $f(x)'=\frac{x}{\sqrt{x^2-4}}\Rightarrow f(3)'=\dfrac{3}{\sqrt{3^2-4}}=\dfrac{3\sqrt{5}}{5}$.
   \itemch Đúng. Ta có $\left(f\left[g(x\right]\right)'=g'(x)\cdot f'\left[g(x)\right]=3\cdot\left(\dfrac{g(x)}{\sqrt{g^2(x)-4}}\right)$.\\ 
   mà $g(3)=7$. Suy ra đạo  hàm của hàm số $f\left[g(x)\right]$ tại $x=3$ bằng $\dfrac{7\sqrt{5}}{5}$.\\
   \itemch Sai. Ta có  $\left(g\left[f(x)\right]\right)'=f'(x)\cdot g'\left[f(x)\right]=\dfrac{x}{\sqrt{x^2-4}}\cdot 3$.\\
   Đạo hàm của hàm số  $g\left[f(x)\right]$ tại $x=3$ bằng $\dfrac{9\sqrt{5}}{5}$.
  \end{itemchoice}
 } 
\end{ex}

% nguồn: Ok_dungsai_lop11_Lan1
\begin{ex}%[1D7H2-1]%[ma:fecc0a6b]
 Tính được đạo hàm của các hàm số sau. Khi đó 
  \choiceTF
  {$y=2\sin x-3\cos x$ có $y^{\prime}=2\cos x-3\sin x$} 
  {\True  $y=3\cot x-\tan x$có $y^{\prime}=-\dfrac{2}{\sin^2x}-\dfrac{1}{\cos^2x}$}
  {$y=x\cos x$  có $y^{\prime}=\cos x+x\sin x$}
  {\True $y=2 x\sin  ^2 x$ có $y^{\prime}=2\sin^2x+2x\sin 2x$}
  \loigiai{
  \begin{itemchoice}
    \itemch $y^{\prime}=2(\sin x)^{\prime}-3(\cos x)^{\prime}=2\cos x+3\sin x$.
   \itemch $y^{\prime}=2(\cot x)^{\prime}-(\tan x)^{\prime}=-\dfrac{2}{\sin ^2  x}-\dfrac{1}{\cos ^2 x}$.\\
    \itemch $y^{\prime}=x^{\prime}\cdot\cos x+(\cos x)^{\prime}\cdot x=\cos x-x\sin x$.\\
   \itemch  Ta có\\
   $\begin{aligned}
   y^{\prime}&=(2 x)^{\prime}\cdot\sin ^2 x+\left(\sin ^2 x\right)^{\prime}\cdot 2 x=2\sin ^2  x+2\sin x(\sin x)^{\prime}\cdot 2 x\\
   &=2\sin ^2 x+4 x\sin x\cos x=2\sin ^2 x+2 x\sin 2 x. 
  \end{aligned}$
 \end{itemchoice}
} 
\end{ex}

% nguồn: Ok_dungsai_lop11_Lan1
\begin{ex}%[1D7H2-1]%[ma:a3d579ba]
 Tính được đạo hàm của các hàm số sau.  Khi đó
 \choiceTF 
  {\True $y=\log_2(9 x-5)$  có $y^{\prime}=\dfrac{9}{(9x-5)\ln 2}$}
  {\True $y=2 e^{3  x+1}$ có $y^{\prime}=6e^{3x+1}$}
  {\True $y=3^{x^3-1}$ có  $y^{\prime}=3\cdot\ln 3\cdot{x^2}\cdot{3^{x^3-1}}$}
  {$y=\ln\sqrt{x}$  có $y^{\prime}=-\dfrac{1}{2x}$}
 \loigiai{ 
  \begin{itemchoice}
   \itemch Đúng.  Vì $y^{\prime}=\dfrac{(9 x-5)^{\prime}}{(9 x-5)\ln 2}=\dfrac{9}{(9 x-5)\ln 2}$.
   \itemch Đúng. Vì $y^{\prime}=2(3 x+1)^{\prime}\cdot e^{3  x+1}=6 e^{3 x+1}$.
   \itemch Đúng. Vì $y^{\prime}=\left(x^3-1\right)^{\prime}\cdot 3^{x^3-1}\cdot\ln 3=3\cdot\ln 3\cdot  x^2\cdot 3^{x^3-1}$.
   \itemch  Sai. Vì $y^{\prime}=\dfrac{(\sqrt{x})^{\prime}}{\sqrt{x}}=\dfrac{1}{2\sqrt{x}\sqrt{x}}=\dfrac{1}{2 x}$.
 \end{itemchoice}}
\end{ex}

% nguồn: Ok_dungsai_lop11_Lan1
\begin{ex}%[1D7H2-1]%[ma:cf18e790]
  Tính được đạo hàm của các hàm số sau. Khi đó
 \choiceTF 
  {$y=2\sin x-3\cos x$ có $y'=2\cos x-3\sin  x$}
  {\True  $y=2\cot x-\tan x$ có $y'=-\dfrac{2}{\sin^2 x}-\dfrac{1}{\cos^2 x}$}
  {$y=x\cos  x$ có $y'=\cos x+x\sin x$}
  {\True $y=2x\sin^2 x$ có  $y'=2\sin^2 x+2x\sin x$}
 \loigiai{ 
  \begin{itemize}
   \item $y'=2(\sin x)'-3(\cos x)' =2\cos x+3\sin  x$.
   \item $y'=2(\cot  x)'-(\tan x)' =-\dfrac{2}{\sin^2 x}-\dfrac{1}{\cos^2 x}$.
   \item $y'=x'\cos  x+x (\cos x)'= \cos x-x\sin x.$
   \item 
   \begin{eqnarray*}
     y'&=&(2x)'\sin^2 x+2x(\sin^2 x)'\\
    &=&2\sin^2 x+2x \cdot 2\sin x.(\sin  x)'\\
    &=& 2\sin ^2 x+4x\sin x  \cos x\\
    &=& 2\sin^2  x+2x\sin 2x.
   \end{eqnarray*}
  \end{itemize}
 } 
\end{ex}

% nguồn: Ok_dungsai_lop11_Lan1
\begin{ex}%[1D7H2-1]%[ma:ba7ea7bb]
 Cho hàm số $y=-x^3+3x+3$. Xét tính đúng sai  của các mệnh đề sau
 \choiceTF 
  {\True $y'=0$ có hai nghiệm  phân biệt}
   { $y'<0, \forall x\in\mathbb{R}$}
  {\True $y'$ đạt giá trị lớn  nhất tại $x=0$}
  {$y'(0)=1$} 
 \loigiai{ 
   Ta có $y=-x^3+3x+3\Rightarrow y'=-3x^2+3$.
  \begin{itemchoice}
   \itemch ĐÚNG $y'=0\Leftrightarrow  x=\pm 1$.
   \itemch SAI  vì $y'\left(0\right)>0$.
   \itemch ĐÚNG $y'=-3x^{2} +3\le 3, \forall x\in\mathbb{R}$. Nên GTLN của  $y'$ bằng 3 khi $x=0$.
   \itemch  SAI vì $y'(0)=3$.
  \end{itemchoice}
  }
\end{ex}

% nguồn: Ok_dungsai_lop11_Lan1
\begin{ex}%[1D7H2-1]%[ma:923b9879]
 Cho  hàm số $f(x)=\sqrt{x+1}$, các mệnh đề sau đây đúng hay sai?
  \choiceTF
   {\True $f'(3)=\dfrac{1}{4}$}
   {\True $f'(0)=\dfrac{1}{2}$}
   {$f'(2)=\dfrac{1}{6}$}
   {$f'(8)=\dfrac{1}{18}$}
  \loigiai{
  Ta  có $f'(x)=\dfrac{1}{2\sqrt{x+1}}$.
  \begin{itemchoice}
   \itemch $f'(3)=\dfrac{1}{4}$. 
    \itemch $f'(0)=\dfrac{1}{2}$.
   \itemch $f'(2)=\dfrac{1}{2\sqrt{3}}$. 
   \itemch  $f'(8)=\dfrac{1}{6}$.
  \end{itemchoice}
 } 
\end{ex}

% nguồn: Ok_dungsai_lop11_Lan1
\begin{ex}%[1D7H2-1]%[ma:bfb26e8a]
 Cho hàm số $y=\log_2x$. Xét tính đúng sai  của các mệnh đề sau
  \choiceTF
  {$y'=\dfrac{\ln 2}{x}, \forall x>0$} 
  {\True $y'=\dfrac{1}{x\ln  2}, \forall x>0$}
  {\True  $y'(1)=\dfrac{1}{2}$}
   {$y'(4)=\dfrac{1}{\ln 8}$}
 \loigiai{ 
  \begin{itemchoice}
    \itemch Sai vì $y'=\left(\log_2x\right)'=\dfrac{1}{x\ln 2}, \forall x>0$.
   \itemch Đúng vì $y'=\dfrac{1}{x\ln  2}, \forall x>0$.
   \itemch Đúng vì $y'=\dfrac{1}{x\ln 2} \Rightarrow y'(1)=\dfrac{1}{\ln  2}$.
    \itemch Sai vì $y'=\dfrac{1}{x\ln 2}\Rightarrow $$y'(4)=\dfrac{1}{4\ln 2} =\dfrac{1}{\ln 16}$.
  \end{itemchoice}
 } 
\end{ex}

% nguồn: Ok_dungsai_lop11_Lan1
\begin{ex}%[1D7H2-1]%[ma:0205d18c]
 Cho hàm số $y=(2 \sqrt{x}+3)(4 \sqrt{x}-1)$. Khi  đó:
  \choiceTF
  {\True $y'\left( 1  \right)=13$}
  {\True $y'\left( 1 \right)+y'\left( 25 \right)=22$} 
  {Đồ thị của hàm  số $y'$ đi qua điểm $A\left( 25;-9 \right)$}
  {\True $y'\left( 1 \right)>y'\left(  2 \right)$}
  \loigiai{\phantom{a}\\
  ${\begin{aligned}
  y^{\prime} & =(2 \sqrt{x}+3)^{\prime}(4  \sqrt{x}-1)+(4 \sqrt{x}-1)^{\prime}(2 \sqrt{x}+3) \\
  & =2 \cdot  \dfrac{1}{2 \sqrt{x}}(4 \sqrt{x}-1)+4 \cdot \frac{1}{2 \sqrt{x}}(2 \sqrt{x}+3)=8+\frac{5}{\sqrt{x}}.
 \end{aligned}
 }$ 
} 
\end{ex}

% nguồn: Ok_dungsai_lop11_Lan1
\begin{ex}%[1D7H2-1]%[ma:49f0dd09]
 Cho hàm của hàm số $f(x)=2^x$. Xét  tính đúng sai của các mệnh đề sau
  \choiceTF
  {\True $f'(x)=2^x\ln  2, \forall x\in\mathbb{R}$}
   {$f'(x)\geq 2, \forall x\in\mathbb{R}$}
   {$f'(0)=0$}
  {\True $f'(1)=2\ln  2$}
  \loigiai{
  Ta  có $f'(x)=2^x\ln 2, \forall x\in\mathbb{R}$.
  \begin{itemchoice}
   \itemch  Đúng.
   \itemch  Sai vì $f'(-2)=2^{-2}\ln 2<2$.
   \itemch Sai vì $f'(0)=2^0\ln 2=\ln 2$. 
   \itemch Đúng vì $f'(1)=2^1\ln 2=2\ln 2$. 
  \end{itemchoice}
 } 
\end{ex}

% nguồn: Ok_dungsai_lop11_Lan1
\begin{ex}%[1D7H2-1]%[ma:e363fe7f]
 Cho  hàm của hàm số $f(x)=x^3+3x^2+x-5$. Xét tính đúng sai của các mệnh đề sau
 \choiceTF 
  {\True $f'(x)=3x^2+6x+1, \forall x\in\mathbb{R}$} 
   {$f'(-1)=2$}
  {\True $f'(x)<1\Leftrightarrow x\in  (-2;0)$}
  {\True $f'(x)=-2\Leftrightarrow x=-1$} 
 \loigiai{ 
  Ta  có $f'(x)=3x^2+6x+1, \forall x\in\mathbb{R}$.
  \begin{itemchoice}
   \itemch Đúng. 
   \itemch Sai vì  $f'(-1)=3\cdot (-1)^2+6\cdot (-1)+1=-2$.
   \itemch Đúng vì $f'(x)<1\Leftrightarrow 3x^2+6x+1<1\Leftrightarrow 3x^2+6x<0\Leftrightarrow x\in(-2;0)$. 
   \itemch Đúng vì $f'(x)=-2\Leftrightarrow 3x^2+6x+1=-2\Leftrightarrow 3x^2+6x+3=0\Leftrightarrow  x=-1$.
  \end{itemchoice}
 } 
\end{ex}

% nguồn: Ok_dungsai_lop11_Lan1
\begin{ex}%[1D7H2-1]%[ma:2a02a005]
 Tính được đạo hàm của các hàm số sau. Khi  đó
  \choiceTF
   {\True $y=(x^2+x)\mathrm{e}^x$ có $y'=(x^2+3x+1)\mathrm{e}^x$}
  {\True $y=\dfrac{x^3}{\ln x}$ có $y'=\dfrac{3x^2\ln  x-x^2}{\ln^2x}$}
  {\True $y=\dfrac{2x^2}{2^x}$ có  $y'=\dfrac{4x-2\ln 2\cdot x^2}{2^x}$}
   {\True $y=3x\cdot\log_3x$ có $y'=3\log_3x+\dfrac{3}{\ln 3}$}
 \loigiai{ 
  \begin{itemize}
   \item Ta có $y'=(2x+1)\mathrm{e}^x+(x^2+x)\mathrm{e}^x=(x^2+3x+1)\mathrm{e}^x$. 
   \item Ta có  $y'=\dfrac{3x^2\ln x-x^3\dfrac{1}{x}}{\ln^2x}=\dfrac{3x^2\ln x-x^2}{\ln^2x}$.
   \item Ta có $y=\dfrac{4x\cdot  2^x-2x^2\cdot 2^x\ln 2}{\left(2^x\right)^2}=\dfrac{4x-2x^2\ln 2}{2^x}$.
   \item Ta có $y'=3\log_3x+3x\cdot\dfrac{1}{x\cdot\ln  3}=3\log_3x+\dfrac{3}{\ln 3}$.
  \end{itemize}
 } 
\end{ex}

% nguồn: Ok_dungsai_lop11_Lan1
\begin{ex}%[1D7H2-1]%[ma:1079c9a8]
 Cho hàm số $y=x\cdot\mathrm{e}^{4x}$. Xét tính đúng sai của  các mệnh đề sau
 \choiceTF 
   {$y'=(4+4x)\cdot\mathrm{e}^{4x}$}
  {\True $y'(1)=5\cdot\mathrm{e}^4$} 
   {$y'(0)=-1$}
  {\True $y'=\mathrm{e}^{4x}+4y$} 
  \loigiai{
  Ta có $y'=\left(x\cdot\mathrm{e}^{4x}\right)'=x'\cdot \mathrm{e}^{4x}+x\cdot\left(\mathrm{e}^{4x}\right)'=\mathrm{e}^{4x}+4x\cdot\mathrm{e}^{4x}=(1+4x)\cdot\mathrm{e}^{4x}$. 
  \begin{itemchoice}
   \itemch Sai. 
   \itemch Đúng. 
   \itemch Sai vì  $y'(0)=1$.
    \itemch Đúng vì $y'=\mathrm{e}^{4x}+4x\cdot \mathrm{e}^{4x}=\mathrm{e}^{4x}+4y$.
  \end{itemchoice}
  }
\end{ex}

% nguồn: Ok_dungsai_lop11_Lan1
\begin{ex}%[1D7H2-1]%[ma:3c0fb340]
 Cho hàm số $y=f(x)=2\sin 3x+\cos 2x$. Xét tính đúng sai các mệnh  đề sau
 \choiceTF 
  {Hàm  số đã cho là hàm số chẵn}
  {\True Hàm số  tuần hoàn với chu kì $T=2\pi$}
   {$f(x)=f\left(\dfrac{\pi}{2}-x\right)$}
  {\True Đạo hàm của hàm số đã  cho là $y'=6\cos 3x-2\sin 2x$}
 \loigiai{ 
  \begin{itemchoice}
   \itemch Tập xác định $\mathscr{D}=\mathbb{R}$.\\ 
   Xét $f(-x)=2\sin  (-3x)+\cos (-2x)=-2\sin 3x+\cos 2x$.\\
    Suy ra $f(-x)\ne f(x)$. \\
   Vậy hàm số đã cho không phải hàm  số chẵn. Mệnh đề đã cho sai.
   \itemch Ta có $\sin 3x=\sin (3x+2\pi)$ và  $\cos 2x=\cos (2x+2\pi)$.\\
   Vậy hàm số đã  cho tuần hoàn với chu kì $T=2\pi$. Mệnh đề đã cho đúng.
    \itemch Ta có
   \begin{eqnarray*}
    f\left(\dfrac{\pi}{2}-x\right)&=&2\sin \left[3\left(\dfrac{\pi}{2}-x\right)\right]+\cos \left[2\left(\dfrac{\pi}{2}-x\right)\right]=2\sin  \left(\dfrac{3\pi}{2}-3x\right)+\cos (\pi-2x)\\
    &=&-2\sin \left(\dfrac{\pi}{2}-3x\right)-\cos  2x\\
     &=&-2\cos 3x-\cos 2x\ne f(x).
   \end{eqnarray*}
   Vậy  mệnh đề đã cho sai.
   \itemch Ta có $y'=2{(\sin  3x)}'+{(\cos 2x)}'=2\cos 3x\cdot(3x)'-\sin 2x\cdot(2x)'=6\cos 3x-2\sin 2x$.\\
   Mệnh đề đã  cho đúng.
  \end{itemchoice}
  }
\end{ex}

% nguồn: Ok_dungsai_lop11_Lan1
\begin{ex}%[1D7H2-1]%[ma:30851536]
  Cho hai hàm số $f(x)=\dfrac{3}{x+1}$ và $g(x)=\dfrac{x^2}{x+2}$. Xét tính đúng sai của các mệnh đề sau
  \choiceTF
  {\True $f'(x)=-\dfrac{3}{(x+1)^2},  \forall x\neq -1$}
  {\True $g'(x)=\dfrac{x^2+4x}{(x+2)^2}, \forall x\neq -2$} 
  {$\left[f(x)\cdot g(x)\right]'=f'(x)\cdot g'(x), \forall x\in\mathbb{R}\backslash  \left\{-1;-2\right\}$}
   {$g'(1)=-1$}
  \loigiai{
  \begin{itemchoice}
   \itemch  Đúng vì $f'(x)=-\dfrac{3}{(x+1)^2}, \forall x\ne -1$.
   \itemch  Đúng vì $g'(x)=\dfrac{2x\cdot (x+2)-x^2}{(x+2)^2}=\dfrac{x^2+4x}{(x+2)^2}$, $\forall x\neq -2$.
   \itemch  Sai vì $\left[f(x)\cdot g(x)\right]'=f'(x)\cdot g(x)+f(x)\cdot g'(x),\forall x\in\mathbb{R}\backslash \left\{-1;-2\right\}$.
   \itemch Sai vì  $g'(1)=\dfrac{5}{9}$.
  \end{itemchoice}
 } 
\end{ex}

% nguồn: Ok_dungsai_lop11_Lan1
\begin{ex}%[1D7H2-1]%[1D7H2-3]%[ma:10897302]
 Cho  hàm số $y=x^4-4x^2+3\sqrt{x}+2-\dfrac{1}{x}$. Khi đó
 \choiceTF 
  {\True $y^{\prime}(1)=-\dfrac{3}{2}$} 
  {Đồ thị của hàm số  $y^{\prime}$ đi qua điểm $A\left(1; \dfrac{3}{2}\right)$}
   {\True $y^{\prime}(4)=\dfrac{3597}{16}$}
  {\True Điểm $M$ thuộc đồ thị $(C)$ của hàm số $y=x^4-4x^2+3\sqrt{x}+2-\dfrac{1}{x}$ có hoành độ $x_0=1$. Khi đó, phương trình tiếp tuyến của $(C)$ tại $M$ vuông  góc với đường thẳng $y=\dfrac{2}{3} x$}
  \loigiai{
  $y^{\prime}=4x^3-8x+\dfrac{3}{2 \sqrt{x}}+0-\left(-\dfrac{1}{x^2}\right)=4 x^3-8  x+\dfrac{3}{2 \sqrt{x}}+\dfrac{1}{x^2}$.\\
  Khi đó 
  \begin{itemize}
   \item $y'(1)=-\dfrac{3}{2}$. 
   \item Do $y'(1)=-\dfrac{3}{2}$ nên đồ thị hàm số $y'$ không đi qua điểm $A\left(1;  \dfrac{3}{2}\right)$.
   \item $y'(4)=\dfrac{3597}{16}$. 
   \item  Gọi $\Delta\colon y=y'(x_0)\left(x-x_0\right)+y_0$ là phương trình tiếp tuyến của $(C)$ tại $M(x_0;y_0)$. \\
   Ta có $x_0=1$ nên $y_0=1$ và $y'(x_0)=-\dfrac{3}{2}$. Khi  đó
    \allowdisplaybreaks
   \begin{eqnarray*}
    \Delta \colon y=-\dfrac{3}{2}\left(x-1\right)+1 \Leftrightarrow \Delta \colon y=-\dfrac{3}{2}x+\dfrac{5}{2}. 
   \end{eqnarray*}
   Do đó, $\Delta$ vuông  góc với đường thẳng $y=\dfrac{2}{3}x$.
  \end{itemize}
 } 
\end{ex}

% nguồn: Ok_dungsai_lop11_Lan1
\begin{ex}%[1D7H2-1]%[ma:58483d69]
 Cho hàm số $f(x)=\sqrt{x^2-4}$. Xét  tính Đúng - Sai của các mệnh đề sau:
  \choiceTF
  {Hàm số đã cho có đạo  hàm tại $x=2$}
  {$f'(x)=\dfrac{1}{2  \sqrt{x^2-4}}$}
  {\True  $f'(3)=\dfrac{3\sqrt{5}}{10}$}
  {Phương trình $f'(x)=2$ có  hai nghiệm}
 \loigiai{ 
  \begin{itemchoice}
   \itemch Hàm số đã  cho không có đạo hàm tại $x=2$.
    \itemch $f'(x)=\dfrac{2x}{2 \sqrt{x^2-4}}=\dfrac{x}{\sqrt{x^2-4}}.$
   \itemch $f'(3)= \dfrac{3}{2  \sqrt{3^2-4}}=\dfrac{3}{2 \sqrt{5}}=\dfrac{3\sqrt{5}}{10}.$
   \itemch 
    $f'(x)=2
   \Leftrightarrow \dfrac{x}{\sqrt{x^2-4}}=2 
    \Leftrightarrow 2 \sqrt{x^2-4} =x
   \Leftrightarrow  \heva{&x > 2\\& 4x^2-16=x^2}
   \Leftrightarrow \heva{&x > 2\\& \hoac{&x=\tfrac{4 \sqrt{3}}{3}  \text{~(thỏa mãn)}\\&x=\tfrac{-4\sqrt{3}}{3} \text{~(loại)}.}}$\\
   Vậy  phương trình $f'(x)=2$ có một nghiệm.
  \end{itemchoice}
 } 
\end{ex}

% nguồn: Ok_dungsai_lop11_Lan1
\begin{ex}%[1D7H2-1]%[ma:7b7e5da8]
 Cho hàm số  $y=(2x-1)\cdot\log_8x$. Xét tính đúng sai của các mệnh đề sau
 \choiceTF 
  {$y'=2\cdot\log_8x+\dfrac{2x-1}{x\cdot\ln 8}$} 
   { $y'(2)=2+\dfrac{5}{8\ln 8}$}
  {$y'(8)=2+\dfrac{5}{8\cdot\ln 2}$} 
   {$y'(2)-y'(8)=\dfrac{1}{3}-\dfrac{2}{\ln 8}$}
 \loigiai{ 
  $y'=(2x-1)'\cdot\log_8x+(2x-1)\cdot\left(\log_8x\right)'=2\cdot\log_8x+\dfrac{2x-1}{x\cdot\ln 8}$. 
  \begin{itemchoice}
    \itemch Đúng.
   \itemch  Sai vì $y'(2)=2\log_{2^3}2+\dfrac{3}{2\ln 8}=\dfrac{2}{3}+\dfrac{3}{2\ln 8}$.
    \itemch Đúng vì $y'(8)=2+\dfrac{15}{8\cdot\ln 8}=2+\dfrac{5}{8\cdot\ln 2}$.
   \itemch Sai vì $y'(2)-y'(8)=\dfrac{2}{3}+\frac{3}{2\ln 8}-\left(2+\dfrac{5}{8\cdot\ln 2}\right)=-\dfrac{4}{3}+\dfrac{3}{2\ln 8}-\dfrac{15}{8\cdot\ln  8}=-\dfrac{4}{3}-\dfrac{3}{8\ln 8}$.
  \end{itemchoice}
 } 
\end{ex}

% nguồn: Ok_dungsai_lop11_Lan1
\begin{ex}%[1D7H2-1]%[ma:c7cf6910]
  Cho hàm số $y=\dfrac{x^2+1}{x-3}$. Các khẳng định sau là \textbf{đúng} hay \textbf{sai}?
 \choiceTF 
  {$y'(1)=\dfrac{3}{2}$} 
  {Tổng các nghiệm của phương trình $y'=0$ bằng  $-6$}
  {\True Đồ thị của  hàm số $y'$ đi qua điểm $A\left(1;-\dfrac{3}{2}\right)$}
   {$y'(1)<y'(2)$}
 \loigiai{Ta có $y'=\dfrac{2x(x-3)-(x^2+1)}{(x-3)^2}=\dfrac{x^2-6x-1}{(x-3)^2}$. 
 \begin{itemize}
  \item Ta có $y'(1)= -  \dfrac{3}{2}$.
  \item Ta có $y'=0\Leftrightarrow x^2-6x-1=0\Leftrightarrow\hoac{&x_1=3-\sqrt{10}\\&x_2=3+\sqrt{10}}$, suy  ra $x_1+x_2=6$.
  \item Vì $y'(1)= -  \dfrac{3}{2}$ nên đồ thị của hàm số $y'$ đi qua điểm $A\left(1;-\dfrac{3}{2}\right)$.
  \item Ta có $y'(1)=  - \dfrac{3}{2}$ và $y'(2)=-9$ nên $y'(1)>y'(2)$.
 \end{itemize}
  }
\end{ex}

% nguồn: Ok_dungsai_lop11_Lan1
\begin{ex}%[1D7H2-1]%[ma:5920e85d]
 Cho hàm số $y=\dfrac{\ln x}{2^x}$. Xét  tính đúng sai của các mệnh đề sau
 \choiceTF 
   {$y'(1)=1$}
  {$y'(2)=-\dfrac{1}{2}$} 
  {\True $y'\cdot 4^x=\dfrac{2^x}{x}-\left(\ln x\right)\cdot 2^x\cdot\ln 2$} 
  {\True  $y'(1)+y(1)=\dfrac{1}{2}$}
  \loigiai{
  Ta  có $y'=\dfrac{\left(\ln x\right)'\cdot 2^x-\ln x\cdot\left(2^x\right)'}{2^{2x}}=\dfrac{\frac{2^x}{x}-\left(\ln x\right)\cdot 2^x\cdot \ln 2}{2^{2x}}$.
  \begin{itemchoice}
   \itemch  Sai vì $y'(1)=\dfrac{1}{2}$.
   \itemch Sai vì $y'(2)=\dfrac{2-4\cdot\left(\ln 2\right)^2}{16}=\dfrac{1-2\cdot\left(\ln 2\right)^2}{8}$. 
   \itemch Đúng vì $y'\cdot 4^x=\dfrac{\frac{2^x}{x}-\left(\ln x\right)\cdot  2^x\cdot \ln 2}{2^{2x}}\cdot 4^{x}=\dfrac{2^{x}}{x}-\left(\ln x\right)\cdot 2^x\cdot \ln 2$.
   \itemch Đúng. 
  \end{itemchoice}
  }
\end{ex}

% nguồn: Ok_dungsai_lop11_Lan1
\begin{ex}%[1D7H2-1]%[ma:6ef0b788]
 Cho hàm số $y=(3-4x)^{11}$. Xét tính đúng sai của các mệnh đề  sau
  \choiceTF
   {\True $y'=-44\cdot (3-4x)^{10}$}
   {$y''=-440\cdot (3-4x)^{9}$}
   {\True $y'(1)=-44$}
   {$y'(1)+y'\left(\dfrac{1}{2}\right)=0$}
 \loigiai{ 
  \begin{itemchoice}
   \itemch Đúng  vì $y'=11\cdot (3-4x)^{10}\cdot (3-4x)'=-44\cdot (3-4x)^{10}$.
   \itemch Sai vì $y''=-44\cdot 10 (3-4x)^9\cdot (-4)=1760(3-4x)^9$. 
   \itemch  Đúng.
   \itemch Sai vì  $y'(1)=-44$; $y'\left(\dfrac{1}{2}\right)=-44$;
    $y'(1)+y'\left(\dfrac{1}{2}\right)=-88$.
  \end{itemchoice}
  }
\end{ex}

% nguồn: Ok_dungsai_lop11_Lan1
\begin{ex}%[1D7V2-1]%[ma:67415610]
 Cho hàm  số $y=\dfrac{x^2+1}{x-3}$. Khi đó
  \choiceTF
   {$ y'(1)=\dfrac{3}{2}$}
  {Tổng các  nghiệm của phương trình $ y'=0$ bằng $-6$}
  {\True Đồ thị của  hàm số $ y'$ đi qua điểm $ A\left(1;-\dfrac{3}{2}\right)$}
  {$  y'(1)<y'(2)$}
 \loigiai{ 
  $y^{\prime}=\dfrac{\left(x^2+1\right)^{\prime}(x-3)-(x-3)^{\prime}\left(x^2+1\right)}{(x-3)^2}=\dfrac{2 x(x-3)-1\left(x^2+1\right)}{(x-3)^2}=\dfrac{x^2-6  x-1}{(x-3)^2}$.
  \begin{itemchoice}
   \itemch  Sai. Vì $ y'(1)=-\dfrac{3}{2}$.
   \itemch  Sai. Tổng các nghiệm của phương trình $ y'=0$ bằng $6$.
   \itemch  Đúng. Vì $ y'(1)=-\dfrac{3}{2}$.
   \itemch Sai. $y'(1) = -1 > y'(2) = -9$. 
  \end{itemchoice}
 } 
\end{ex}

% nguồn: Ok_dungsai_lop11_Lan1
\begin{ex}%[1D7V2-1]%[ma:e36ae0b1]
 Cho hàm số $ y=5\cot x-3 $. 
 \choiceTF 
  {\True $  y'=-5(\cot^2 x+1) $}
  {\True $  y'=\dfrac{-5}{\sin^2{x}}$}
  {$ y'\left(\dfrac{\pi}{2}\right)=5$} 
  {\True $ y\left(\dfrac{\pi}{3}\right)=-\dfrac{9+5\sqrt{3}}{3}$} 
 \loigiai{ 
  \begin{itemchoice}
    \itemch Đúng\\
   Ta có $ y'=\dfrac{-5}{\sin^2{x}}=-5(\cot^2  x+1)$.
   \itemch Đúng\\ 
   Vì  $ y'=\dfrac{-5}{\sin^2{x}}$.
    \itemch Sai\\
    Ta có $ y'=\dfrac{-5}{\sin^2{x}}$.\\
   Suy ra $y'\left(\dfrac{\pi}{2}\right)=-5 $. 
   \itemch  Đúng\\
   Vì  $ y\left(\dfrac{\pi}{3}\right)=5\cot\left(\dfrac{\pi}{3}\right)-3=-\dfrac{9+5\sqrt{3}}{3}$.
  \end{itemchoice}
 } 
\end{ex}

% nguồn: Ok_dungsai_lop11_Lan1
\begin{ex}%[1D7V2-1]%[ma:1992ac8f]
 Cho hàm số $y=-4 x^3+\dfrac{x^2}{2}-2 x+3$, biết $ y'=a{x^2}+bx+c$. Khi đó 
  \choiceTF
  {$ a+b+c=-10$} 
  {Phương trình $ y'=0$ có hai nghiệm  phân biệt}
  {\True  Đồ thị hàm số $ y'$ cắt trục tung tại điểm $\left(0;-2\right)$}
  {Đồ thị hàm số $ y'$ cắt đường thẳng  $ y=3$ tại hai điểm phân biệt}
  \loigiai{
  $y^{\prime}=-4\cdot 3\cdot x^2+\dfrac{1}{2}\cdot 2\cdot x-2+0=-12  x^2+x-2$
  \begin{itemchoice}
   \itemch Sai. Vì $a+b+c=-12$. 
   \itemch Sai. Phương  trình $ y'=0$ vô nghiệm.
   \itemch Đúng. $y'(0) = -2$, nên đồ thị hàm số $  y'$ cắt trục tung tại điểm $\left(0;-2\right)$.
   \itemch Sai. Vì ta có $y'=3 \Leftrightarrow -12x^2+x-5=0$ vô  nghiệm.
  \end{itemchoice}
 .} 
\end{ex}

% nguồn: Ok_dungsai_lop11_Lan2
\begin{ex}%[1D7V2-1]%[ma:c9a92e98]
	Cho hàm số $f(x)=-x^3+3mx^2-12x+3$, với $m$  là tham số thực.
	 \choiceTF
		{Với  $m=1$ thì $f'(x)=-3x^2+5x-12$}
		{\True  Với $m=0$ thì $f'(1)=f'\left(-1\right)$}
		{Với $m=0$ thì $f'(x) > 0, \forall  x\in \mathbb{R}$}
		{$f'(x)\le 0, \forall  x\in \mathbb{R}$ khi và chỉ khi $-2< m < 2$}
	 \loigiai{
		Ta  có $f'(x)=-\left(x^3\right)'+3m\left(x^2\right)'-12(x)'+3'=-3x^2+6mx-12$.
		\begin{itemchoice}
			\itemch  Với $m=1$ thì $f'(x)=-3x^2+6x-12$.
			\itemch Với $m=0$ thì  $f'(x)=-3x^2-12$.\\
			 Khi đó $f'(1)=-3\cdot 1^2-12=-3\cdot (-1)^2-12=f'(-1)$.
			\itemch Với $m=0$ thì $-3x^2-12< 0, \forall x\in \mathbb{R}$ hay $f'(x) < 0, \forall x\in \mathbb{R}$ với  $m=0$.
			 \itemch \allowdisplaybreaks
			\begin{eqnarray*}
				 f'(x)\le 0, \forall x\in \mathbb{R}&\Leftrightarrow& \heva{&-3< 0\\&\Delta '\le 0}\\
				 &\Leftrightarrow& (3m)^2-(-3)(-12)\le 0\\
				&\Leftrightarrow&  9m^2-36\le 0\\
				 &\Leftrightarrow&-2\le m\le 2.
			\end{eqnarray*}
		\end{itemchoice}
	} 
\end{ex}

% nguồn: Ok_dungsai_lop11_Lan2
\begin{ex}%[1D7V2-1]%[ma:0a529c08]
	Cho hàm  số $f(x)=\ln x-\ln\left(x+1\right)$.
	\choiceTF 
		{\True Tập xác định của hàm  số là $\mathscr{D}=\left(0;+\infty\right)$}
		 {\True $f'(x)=\dfrac{1}{x}-\dfrac{1}{x+1}$}
		{$\lim\limits_{x\to+\infty}f(x)=1$} 
		{\True Giá trị của biểu thức $P=f'(1)+f'(2)+f'(3)+\cdots+f'(2\,023)+f'(2\,024)$ bằng $\dfrac{2\,024}{2\,025}$} 
	\loigiai{ 
		\begin{itemchoice}
			\itemch Điều kiện xác  định $\heva{&x > 0\\&x+1 > 0}\Leftrightarrow x > 0$.\\
			Nên tập xác định là  $\mathscr{D}=\left(0;+\infty\right)$.
			\itemch $f(x)=\ln x-\ln\left(x+1\right)\Rightarrow  f'(x)=\dfrac{1}{x}-\dfrac{1}{x+1}$.
			\itemch $\lim\limits_{x\to+\infty}f(x)=\lim\limits_{x\to+\infty}\left[\ln x-\ln\left(x+1\right)\right]=\lim\limits_{x\to+\infty}\left(\ln\dfrac{x}{x+1}\right)=\lim\limits_{x\to+\infty}\left(\ln\dfrac{1}{1+\frac{1}{x}}\right)=0$. 
			 \itemch $f'(x)=\dfrac{1}{x}-\dfrac{1}{x+1}$ nên $f'(1)=1-\dfrac{1}{2}$; $f'(2)=\dfrac{1}{2}-\dfrac{1}{3}$; $\ldots$; $f'(2\,024)=\dfrac{1}{2\,024}-\dfrac{1}{2\,025}$.\\
			Do đó 
			\begin{eqnarray*}
				&P &=f'(1)+f'(2)+f'(3)+\cdots+f'(2\,023)+f'(2\,024)\\ 
				& &=\left(1-\dfrac{1}{2}\right)+\left(\dfrac{1}{2}-\dfrac{1}{3}\right)+\cdots+\left(\dfrac{1}{2\,024}-\dfrac{1}{2\,025}\right)\\ 
				& &=1-\dfrac{1}{2\,025}=\dfrac{2\,024}{2\,025}. 
			\end{eqnarray*}
		\end{itemchoice}
	} 
\end{ex}

% nguồn: Ok_dungsai_lop11_Lan2
\begin{ex}%[1D7V2-1]%[ma:735bd0dd]
	Cho hàm số $y=x^4-4x^2-2\sqrt  x+2\,024$.
	\choiceTF 
		{\True  $y'=4x^3-8x-\dfrac{1}{\sqrt x}$}
		{Đồ thị của hàm số $y'$  đi qua điểm $A\left(1;2019\right)$}
		{\True $y'(1) <  y'(2)$}
		 {\True $y'(1)+y'(4)=\dfrac{437}{2}$}
	\loigiai{ 
		\begin{itemchoice}
			\itemch  Ta có $y'=4x^3-8x-\dfrac{1}{\sqrt x}$.
			\itemch Vì $y'(1)=-5$ nên đồ thị của hàm số $y'$  không đi qua điểm $A\left(1;2019\right)$.
			\itemch Vì $y'(1)=-5$; $y'(2)=16-\dfrac{1}{\sqrt 2}$. 
			\itemch Vì $y'(1)=-5$; $y'(4)=\dfrac{447}{2}$ nên  $y'(1)+y'(4)=\dfrac{437}{2}$.
	\end{itemchoice}}
\end{ex}

% nguồn: Ok_dungsai_lop11_Lan2
\begin{ex}%[1D7V2-1]%[ma:7e6ce619]
	Cho hàm số $y=(x-2) \sqrt{x^2+m}$ với $m$ là tham  số.
	 \choiceTF
		{Hàm số $f(x)$ có đạo hàm là $f'(x)=\dfrac{x^2-2  x+m}{\sqrt{x^2+m}}$}
		{\True $f'(0)=\sqrt{m}$} 
		{\True Hệ số góc tiếp  tuyến của đồ thị hàm số tại $x_0=1$ là $\dfrac{m}{\sqrt{m+1}}$}
		{\True Để phương trình  $f(x)=0$ có hai nghiệm phân biệt thì $m<\dfrac{1}{2}$}
	\loigiai{ 
		\begin{itemchoice}
			\itemch 
			Ta có $y'=(x-2)' \cdot \sqrt{x^2+m}+(x-2) \cdot\left(\sqrt{x^2+m}\right)'=\sqrt{x^2+m}+\dfrac{x(x-2)}{\sqrt{x^2+m}}=\dfrac{2 x^2-2 x+m}{\sqrt{x^2+m}}$. 
			 \itemch
			Ta có $f'(0)=\dfrac{2 \cdot  0^2-2 \cdot 0+m}{\sqrt{0^2+m}}=\sqrt{m}$.
			\itemch 
			Hệ số góc tiếp tuyến của đồ thị hàm  số tại $x_0=1$ là $k=f'(1)=\dfrac{2 \cdot 1^2-2 \cdot 1+m}{\sqrt{1^2+m}}=\dfrac{m}{\sqrt{m+1}}$.
			 \itemch
			Ta có $y'=0 \Leftrightarrow \dfrac{2 x^2-2 x+m}{\sqrt{x^2+m}}=0  \Leftrightarrow 2 x^2-2 x+m=0. \quad (1)$\\
			Phương trình $  f'(x)=0 $ có hai nghiệm phân biệt thì phương trình $ (1) $ có hai nghiệm phân biệt khi và chỉ khi
			$$\Delta'>0  \Leftrightarrow (-1)^2-2 \cdot m >0 \Leftrightarrow m <\dfrac{1}{2}.$$
		\end{itemchoice}
	} 
\end{ex}

% nguồn: Ok_dungsai_lop11_Lan2
\begin{ex}%[1D7V2-1]%[ma:ce00f793]
	Cho hàm số  $y=f(x)=\ln\dfrac{x}{x+1}-2\,025$.
	\choiceTF 
		{\True Tập xác định của hàm số là  $\left(-\infty;-1\right)\cup\left(0;+\infty\right)$}
		 {Đạo hàm của hàm số là $y'=-\dfrac{1}{x^2+x}$}
		{Giá trị  $y'(3)$ bằng $\dfrac{13}{12}$}
		{\True Tổng  $T=f'(1)+f'(2)+\cdots+f'\left(2\,025\right)$ bằng $\dfrac{2\,025}{2\,026}$}
	\loigiai{ 
		\begin{itemchoice}
			\itemch Điều kiện  xác định $\dfrac{x}{x+1}> 0\Leftrightarrow\hoac{&x <-1\\&x > 0.}$\\
			Tập xác định của hàm số là  $\mathscr{D}=\left(-\infty-1\right)\cup\left(0;+\infty\right)$.
			\itemch Ta  có $y'=\dfrac{\left(\dfrac{x}{x+1}\right)'}{\dfrac{x}{x+1}}=\dfrac{1}{\left(x+1\right)^2}\cdot \dfrac{x+1}{x}=\dfrac{1}{x\left(x+1\right)}=\dfrac{1}{x^2+x}$.
			\itemch Do  $y'=\dfrac{1}{x^2+x}\Rightarrow y'(3)=\dfrac{1}{12}$.
			\itemch Ta  có $f'(x)=\dfrac{1}{x\left(x+1\right)}=\dfrac{1}{x}-\dfrac{1}{x+1}$.\\
			 Do đó $f'(1)=1-\dfrac{1}{2}$; $f'(2)=\dfrac{1}{2}-\dfrac{1}{3}$; $\ldots$; $f'\left(2\,025\right)=\dfrac{1}{2\,025}-\dfrac{1}{2\,026}$.\\
			 Vậy $f'(1)+f'(2)+\cdots+f'\left(2\,025\right)=1-\dfrac{1}{2\,026}=\dfrac{2\,025}{2\,026}$.
	\end{itemchoice}}
\end{ex}

% nguồn: Ok_dungsai_lop11_Lan1
\begin{ex}%[1D7B2-1]%[ma:931a626a]
 Cho hai hàm số $f(x)= x + 2$ và $g(x) = x^2 -  2x + 3$.
  \choiceTF
   {\True $f'(x)=1$}
  {$g'(x)=x-2$} 
   {$[f(g(1))]'=2$}
  {\True $[g(f(1))]'=  4$}
 \loigiai{ 
   Ta có
  \begin{itemchoice}
   \itemch  $f'(x)=1$.
    \itemch $g'(x)=2x-2$.
   \itemch $[f(g(x))]'= f'(g(x))  \cdot g'(x) = 1\cdot (2x-2) =2x-2$, suy ra $[f(g(0))]'=-2$
   \itemch $[g(f(x))]'=g'\left(f(x)\right) \cdot f'(x) = \left[2(x+2)-2\right]\cdot 1= 2x+2$, suy ra $[g(f(1))]' = 4$ 
  \end{itemchoice}
  }
\end{ex}


%%%%%%%%%%%%%%%% TRẢ LỜI NGẮN %%%%%%%%%%%%%%%%

% nguồn: Ok_TraLoiNgan_Lop11_Lan1
\begin{ex}%[1D7N2-1]%[ma:c2e16767]
	Cho $f(x)=\sqrt{x^3+1}-3x^2+5ax-2$. Tìm $a$ để  $f'(2)=-5$.
	 \shortans{a=1.}
	 \loigiai{Ta có $f'(x)=\dfrac{3x^2}{2\sqrt{x^3+1}}-6x+5a$.\\
		 Do $f'(2)=-5\Leftrightarrow \dfrac{3\cdot2^2}{2\sqrt{2^1}}-6\cdot 2+5a=-5\Leftrightarrow a=-1$.
	} 
\end{ex}

% nguồn: Ok_TraLoiNgan_Lop11_Lan1
\begin{ex}%[1D7H2-1]%[ma:6c08d669]
	Cho hàm số $f(x)=\left(x^{2}-3 x\right) \sqrt{x-1}$. Nghiệm của phương trình $f'(x)=0$ là 
	\shortans{2} 
	 \loigiai{
		 Tập xác định $\mathscr{D} = [1;+\infty)$. \\
		Khi $x>1$, ta  có
		\[f'(x) = (2x-3)\sqrt{x-1} + \dfrac{x^2-3x}{2\sqrt{x-1}}  = \dfrac{2(2x-3)(x-1)+x^2-3x}{2\sqrt{x-1}} = \dfrac{5x^2-13x+6}{2\sqrt{x-1}}.\]
		Khi đó 
		\begin{align*}
			f'(x) =0 \Leftrightarrow \heva{&x>1\\&5x^2-13x+6 =0} \Leftrightarrow  \heva{&x>1\\&\hoac{&x=2\\&x= \dfrac{3}{5} \text{ (loại).}}}
		\end{align*}
		Vậy phương trình  $f'(x)=0$ có tập nghiệm là $S= \{2\}$.
	 }
\end{ex}

% nguồn: Ok_TraLoiNgan_Lop11_Lan1
\begin{ex}%[1D7H2-1]%[ma:df151e88]
	Cho  $\left(\dfrac{3-2x}{\sqrt{4x-1}}\right)'=\dfrac{ax+b}{(4x-1)\sqrt{4x-1}}$.
	Tính  $E=a-b$.
	 \shortans{0}
	 \loigiai{
		- Ta  có $\left(\dfrac{3-2x}{\sqrt{4x-1}}\right)'=\dfrac{-2\sqrt{4x-1}-\dfrac{2(3-2x)}{\sqrt{4x-1}}}{4x-1}=\dfrac{-2(4x-1)-2(3-2x)}{(4x-1)\sqrt{4x-1}}=\dfrac{-4x-4}{(4x-1)\sqrt{4x-1}}$.\\
		- Do đó $a=b=-4$.  \\
		Suy ra  $E=a-b=0$.
	 }
\end{ex}

% nguồn: Ok_TraLoiNgan_Lop11_Lan1
\begin{ex}%[1D7H2-1]%[ma:deec6993]
	Đạo hàm của hàm số $y=\dfrac{\sin ^2 x-\cos ^2 x}{\sin x\cdot\cos x}$  tại điểm $x=\dfrac{\pi}{6}$ (làm tròn đến phần trăm).
	 \shortans{5,33}
	\loigiai{ 
		$y=\dfrac{\sin ^2 x-\cos ^2 x}{\sin x\cdot\cos x}=\dfrac{-\cos 2 x}{\dfrac{1}{2}\sin 2  x}=-2\cot 2 x$\\
	$\Rightarrow  y^{\prime}=-2\dfrac{-2}{\sin ^2 2 x}=\dfrac{4}{\sin ^2 2 x}\Rightarrow y^{\prime}\left(\dfrac{\pi}{6}\right)=\dfrac{16}{3}\approx 5{,}33$.}
\end{ex}

% nguồn: Ok_TraLoiNgan_Lop11_Lan1
\begin{ex}%[1D7H2-1]%[ma:bd0a9384]
	Chứng minh hàm số $f(x)=\sin ^6 x+\cos ^6 x+3\sin  ^2 x\cdot\cos ^2 x-2 x$ có đạo hàm không phụ thuộc $x$.\\
	 \shortans{-2}
	\loigiai{ 
		$\begin{aligned}
		f^{\prime}(x)& =\left(\sin ^6 x+\cos ^6 x+3\sin ^2 x\cdot\cos  ^2 x-2 x\right)^{\prime}\\
		&=6\sin ^5 x\cdot\cos x-6\cos  ^5 x\sin x+6\sin x\cdot\cos ^3 x-6\sin ^3 x\cdot\cos x-2\\
		&=6\sin x\cos x\left(\sin ^4 x-\cos  ^4 x+\cos ^2 x-\sin ^2 x\right)-2\\
		 &=6\sin x\cos x\left(\sin ^2 x\left(\sin ^2 x-1\right)+\cos ^2 x\left(1-\cos ^2 x\right)\right)-2\\
		&=6\sin x\cos x\left(-\cos ^2 x\sin ^2  x+\cos ^2 x\sin ^2 x\right)-2=-2 .
\end{aligned}$}
\end{ex}

% nguồn: Ok_TraLoiNgan_Lop11_Lan2
\begin{ex}%[1D7H2-1]%[ma:912a931d]
	Tính  đạo hàm của hàm số sau $y=-\dfrac{2x+1}{(x+1)^2}$ tại $x_0=-2$.
	\shortans{-4} 
	 \loigiai{
		\begin{eqnarray*}
			&y'&=-\dfrac{(2x+1)'\cdot{(x+1)^2}-\left[(x+1)^2\right]'\cdot (2x+1)}{(x+1)^4}\\ 
			& &=-\dfrac{2(x+1)^2-2(x+1)(2x+1)}{(x+1)^4}\\ 
			 & &=\dfrac{2x^2+2x}{(x+1)^4}=\dfrac{2x}{(x+1)^3}.
		\end{eqnarray*}
	 Suy ra $y'\left(-2\right)=4$.}
\end{ex}

% nguồn: Ok_TraLoiNgan_Lop11_Lan2
\begin{ex}%[1D7H2-1]%[ma:c2842660]
	Tính đạo hàm  hàm số sau $y=\dfrac{3^x-3^{-x}}{3^x+3^{-x}}$ tại $x_0=1$ (kết quả làm tròn đến hàng phần chục).
	\shortans{0,4} 
	\loigiai{ 
		 Ta có
		\begin{eqnarray}
			&y'&=\dfrac{\left(3^x\ln 3+3^{-x}\ln 3\right)\left(3^x+3^{-x}\right)-\left(3^x-3^{-x}\right)\left(3^x\ln 3-3^{-x}\ln  3\right)}{\left(3^x+3^{-x}\right)^2}\\
			&  &=\dfrac{\left(3^x+3^{-x}\right)^2-\left(3^x-3^{-x}\right)^2}{\left(3^x+3^{-x}\right)^2}\ln 3=\dfrac{4\ln 3}{\left(3^x+3^{-x}\right)^2}.
		\end{eqnarray}
		Do đó $ y'(x_0)=y'(1)=\dfrac{4 \ln 3}{\left(3^1+3^{-1}\right)}\approx 0{,}4  $.
	} 
\end{ex}

% nguồn: Ok_TraLoiNgan_Lop11_Lan2
\begin{ex}%[1D7H2-1]%[ma:a9e60225]
	 Cho hàm số $f(x)=\dfrac{3x+1}{\sqrt{x^2+4}}$. Giá trị biểu thức $f'(0)$ là
	 \shortans{1,5}
	 \loigiai{
		Tập  xác định $ \mathscr{D}=\mathbb{R} $.\\
	$f'(x)=\dfrac{3\sqrt{x^2+4}-\left(3x+1\right)\cdot  \dfrac{x}{\sqrt{x^2+4}}}{\left(\sqrt{x^2+4}\right)^2}=\dfrac{12-x}{\sqrt{\left(x^2+4\right)^3}} \Rightarrow f'(0)=\dfrac{3}{2}=1{,}5$.}
\end{ex}

% nguồn: Ok_TraLoiNgan_Lop11_Lan2
\begin{ex}%[1D7H2-1]%[ma:f2d2c4e3]
	Tính đạo  hàm của hàm số sau $y=\dfrac{x^2-3x+3}{x-1}$ tại $x_0=2$.
	 \shortans{0}
	\loigiai{ 
		\begin{eqnarray*}
			 &y'&=\dfrac{\left(x^2-3x+3\right)'(x-1)-(x-1)'\left(x^2-3x+3\right)}{(x-1)^2}\\
			& &=\dfrac{(2x-3)(x-1)-\left(x^2-3x+3\right)}{(x-1)^2}\\ 
			&  &= \dfrac{x^2-2x}{(x-1)^2.}
		\end{eqnarray*}
	Vậy $  y'(2)=0 $.}
\end{ex}

% nguồn: Ok_TraLoiNgan_Lop11_Lan2
\begin{ex}%[1D7H2-1]%[ma:5511bad8]
	Cho hàm số $f(x)=a{x^3}+\dfrac{b}{x}$ có $f'(1)=1$,  $f'\left(-2\right)=-2$. Tính $f'\left(\sqrt 2\right)$.
	 \shortans{-0,4}
	 \loigiai{
		Ta  có $f'(x)=3a{x^2}-\dfrac{b}{x^2}$ nên ta có $$\heva{&f'(1)=3a-b\\& f'\left(-2\right)=12a-\dfrac{b}{4}} \Leftrightarrow\heva{&3a-b=1\\&12a-\dfrac{b}{4}=-2}\Leftrightarrow\heva{&a=-\dfrac{1}{5}\\&b=-\dfrac{8}{5}.}$$
	Suy ra $f'\left(\sqrt  2\right)=6a-\dfrac{b}{2}=-\dfrac{2}{5}=-0{,}4$.}
\end{ex}

% nguồn: Ok_TraLoiNgan_Lop11_Lan2
\begin{ex}%[1D7H2-1]%[ma:424fec7c]
	Tính đạo hàm của hàm  số sau $y=3\ln\left(x^2+1\right)$ tại $x_0=1$.
	 \shortans{3}
	\loigiai{ 
		$y'=\left(3\ln\left(x^2+1\right)\right)'=3\cdot \dfrac{\left(x^2+1\right)'}{x^2+1}=\dfrac{6x}{x^2+1}$.\\ 
	 $y'(1)=\dfrac{6\cdot 1}{1^2+1}=3$.}
\end{ex}

% nguồn: Ok_TraLoiNgan_Lop11_Lan2
\begin{ex}%[1D7H2-1]%[ma:563a2b2b]
	Đạo hàm của hàm số $y=\dfrac{\sin^2x-\cos^2x}{\sin x\cdot\cos x}$ tại điểm $x=\dfrac{\pi}{6}$ (kết quả làm  tròn đến hàng phần trăm).
	 \shortans{5,33}
	 \loigiai{
		 $y=\dfrac{\sin^2x-\cos^2x}{\sin x\cdot\cos x}=\dfrac{-\cos 2x}{\dfrac{1}{2}\sin 2x}=-2\cot 2x$.\\
	Suy ra $ y'=-2\dfrac{-2}{\sin^22x}=\dfrac{4}{\sin^22x}\Rightarrow y'\left(\dfrac{\pi}{6}\right)=\dfrac{16}{3}\approx  5{,}33$.}
\end{ex}

% nguồn: Ok_TraLoiNgan_Lop11_Lan2
\begin{ex}%[1D7H2-1]%[ma:6321ec38]
	Tính đạo hàm của hàm số sau $y=2\mathrm{e}^{2x-1}$  tại $x_0=1$.
	 \shortans{10,9}
	 \loigiai{
		Ta có  $y'=\left(2\mathrm{e}^{2x-1}\right)'=2\cdot (2x-1)'{\mathrm{e}^{2x-1}}=4\mathrm{e}^{2x-1}$.\\
	Suy ra $y'(1)=4\mathrm{e}^{2\cdot  1-1}=4\mathrm{e}\approx 10{,}9$.}
\end{ex}

% nguồn: Ok_TraLoiNgan_Lop11_Lan2
\begin{ex}%[1D7H2-1]%[ma:43ecc24d]
	Tính đạo hàm của hàm số sau $y=\left(x^2+2x-2\right)\cdot{5^x}$ tại $x_0=1$ (kết quả  làm tròn đến hàng đơn vị).
	 \shortans{28}
	 \loigiai{
		Ta có 
		 $$y'=\left(x^2+2x-2\right)'\cdot{5^x}+\left(5^x\right)'\cdot\left(x^2+2x-2\right)=(2x+2)\cdot{5^x}+\left(x^2+2x-2\right)\cdot{5^x}\ln 5.$$
	 Khi đó $y(1)\approx 28$.}
\end{ex}

% nguồn: Ok_TraLoiNgan_Lop11_Lan2
\begin{ex}%[1D7H2-1]%[ma:28813bdf]
	Tính đạo hàm của  hàm số $y=\dfrac{x^2+4x-1}{2x+3}$ tại $x_0=1$.
	 \shortans{0,88}
	 \loigiai{
		\begin{eqnarray*}
			&y'=&\dfrac{\left(x^2+4x-1\right)'\cdot  (2x+3)-(2x+3)'\cdot\left(x^2+4x-1\right)}{(2x+3)^2}\\
			& &=\dfrac{(2x+4)(2x+3)-2\left(x^2+4x-1\right)}{(2x+3)^2}\\ 
			 & &=\dfrac{4x^2+14x+12-2x^2-8x+2}{(2x+3)^2}=\dfrac{2x^2+6x+14}{(2x+3)^2}.
		\end{eqnarray*}
	Vậy $ y'(0)=\dfrac{2\cdot 0^2+6\cdot 0+14}{(2\cdot  0 +3)^2} $.}
\end{ex}

% nguồn: Ok_TraLoiNgan_Lop11_Lan2
\begin{ex}%[1D7H2-1]%[ma:69d4e03f]
	 Tính đạo của hàm số $f(x)=\sin ^6x+\cos ^6x+3\sin ^2x\cdot{\cos ^2}x-2x$.
	\shortans{-2} 
	 \loigiai{
		Ta có 
		\begin{eqnarray*}
			&f'(x)&=\left(\sin^6x+\cos^6x+3\sin^2x\cdot{\cos^2}x-2x\right)'\\ 
			& &=6\sin ^5x\cdot\cos x-6\cos ^5x\sin x+6\sin x\cdot{\cos  ^3}x-6\sin ^3x\cdot\cos x-2\\
			&  &=6\sin x\cos x\left(\sin^4x-\cos^4x+\cos^2x-\sin^2x\right)-2\\
			& &=6\sin x\cos x\left(\sin^2x\left(\sin^2x-1\right)+\cos^2x\left(1-\cos^2x\right)\right)-2\\ 
			& &=6\sin  x\cos x\left(-\cos^2x\cdot {\sin^2}x+\cos^2x\cdot{\sin^2}x\right)-2=-2.
	\end{eqnarray*}}
\end{ex}

% nguồn: Ok_TraLoiNgan_Lop11_Lan2
\begin{ex}%[1D7H2-1]%[ma:56f14d2d]
	 Tính đạo hàm của hàm số $y=5^x+\log_2x+\ln (2x)+\mathrm{e}^{-x}+\sqrt 3 $ tại $x_0=1$ (kết quả làm tròn đến hàng phần chục).
	 \shortans{10,1}
	\loigiai{ 
		 $$y=5^x+\log_2x+\ln (2x)+\mathrm{e}^{-x}+\sqrt 3\Rightarrow y'=5^x\cdot\ln 5+\dfrac{1}{x\ln 2}+\dfrac{1}{x}-\mathrm{e}^{-x}$$
	Suy ra $y'(1)\approx 10{,}1$.} 
\end{ex}

% nguồn: Ok_TraLoiNgan_Lop11_Lan2
\begin{ex}%[1D7H2-1]%[ma:f319f85e]
	Tính đạo  hàm của hàm số sau $y=\dfrac{1+x-x^2}{1-x+x^2}$ tại $x_0=1$.
	 \shortans{-2}
	\loigiai{ 
		\begin{eqnarray*}
			&y'=&\dfrac{\left(1+x-x^2\right)'\left(1-x+x^2\right)-\left(1-x+x^2\right)'\left(1+x-x^2\right)}{\left(1-x+x^2\right)^2}\\ 
			& &=\dfrac{(1-2x)\left(1-x+x^2\right)-(-1+2x)\left(1+x-x^2\right)}{\left(1-x+x^2\right)^2}=\dfrac{2-4x}{\left(1-x+x^2\right)^2}. 
		\end{eqnarray*}
	Vậy $ y'(1)= \dfrac{2-4\cdot 1}{(1-1+1^2)^2}=-2$.} 
\end{ex}

% nguồn: Ok_TraLoiNgan_Lop11_Lan1
\begin{ex}%[1D7V2-1]%[ma:38acdfd6]
	Cho hàm số $f(x)=\dfrac{\sin3x}{3}+\cos x-\sqrt{3}\left(\sin x+\dfrac{\cos3x}{3}\right)$. Tổng các nghiệm thuộc đoạn $[0;\pi]$ của phương trình $f'(x)=0$ bằng bao  nhiêu (kết quả làm tròn đến hàng phần mười)?
	 \shortans{2,6}
	\loigiai{Ta có\\ 
		\begin{eqnarray*}
			f'(x)&=&\dfrac{3\cos3x}{3}-\sin x-\sqrt{3}\left(\cos x-\dfrac{3\sin3x}{3}\right)=\cos3x-\sin x-\sqrt{3}(\cos x-\sin3x)\\ 
			&=&(\cos3x+\sqrt{3}\sin3x)-(\sqrt{3}\cos3x+\sin x)=2\sin\left(3x+\dfrac{\pi}{6}\right)-2\sin\left(x+\dfrac{\pi}{3}\right)\\ 
			&=&4\cos\left(2x+\dfrac{\pi}{4}\right)\cdot\sin\left(x-\dfrac{\pi}{12}\right). 
		\end{eqnarray*}
		Do đó $f'(x)=0\Leftrightarrow\hoac{&\cos\left(2x+\dfrac{\pi}{4}\right)=0\\ 
		 &\sin\left(x-\dfrac{\pi}{12}\right)=0}\Leftrightarrow\hoac{&2x+\dfrac{\pi}{4}=\dfrac{\pi}{2}+k\pi\\
		&x-\dfrac{\pi}{12}=k\pi}\Leftrightarrow\hoac{&x=\dfrac{\pi}{8}+\dfrac{k\pi}{2}\\&x=\dfrac{\pi}{12}+k\pi} (k  \in \mathbb{Z})$.\\
		Do $x\in[0;\pi]$ nên suy ra  $x\in\left\{\dfrac{\pi}{8};\dfrac{\pi}{12};\dfrac{5\pi}{8}\right\}$.\\
	 Tổng các nghiệm là $\dfrac{\pi}{8}+\dfrac{\pi}{12}+\dfrac{5\pi}{8} \approx 2{,}6$.}
\end{ex}

% nguồn: Ok_TraLoiNgan_Lop11_Lan1
\begin{ex}%[1D7V2-1]%[ma:685247e5]
	Có bao nhiêu giá trị nguyên của tham số $m$ để bất phương trình $f'(x)\leq 0$ có tập nghiệm là $\mathbb{R}$, biết $f(x) =  -\dfrac{mx^3}{3}+\dfrac{mx^2}{2}-(3-m)x+2$?
	 \shortans{3}
	\loigiai{ 
	Ta có $f'(x) = -mx^2 + mx + m  - 3$.\\
	Với $f'(x) \leq 0$ ta  được bất phương trình $-mx^2 + mx + m - 3 \leq 0$. \hfill $(*)$
	\begin{itemize}
		\item Khi $m = 0$, bất phương trình $(*)$ nghiệm đúng với mọi $x$ thuộc $\mathbb{R}$. Vậy $m  = 0$ thỏa mãn yêu cầu bài toán.
		\item Khi $m \neq 0$ thì bất phương trình $(*)$ là một bất phương trình trình bậc hai ẩn $x$. Do đó $(*)$ có tập nghiệm là  $\mathbb{R}$ khi và chỉ khi
		\begin{align*}
			\heva{& -m<0 \\& m^2+4m(m-3)\leq 0}\Leftrightarrow \heva{& m>0 \\& 5m^2 - 12m \leq 0} \Leftrightarrow \heva{& m>0  \\& 0 \leq m \leq \dfrac{12}{5}} \Leftrightarrow 0<m \leq \dfrac{12}{5}.
		\end{align*}
	\end{itemize}
	Kết hợp hai trường hợp ta được $m \in \left[0;\dfrac{12}{5}\right]$. Vậy có $3$ giá trị nguyên của $m$ thỏa mãn đề  bài là $
	\{0; 1;  2\}$.
	 }
\end{ex}

% nguồn: Ok_TraLoiNgan_Lop11_Lan1
\begin{ex}%[1D7V2-1]%[2D1V3-6]%[2D1N3-6]%[ma:c994c0ca]
	Cho hàm số $y=\dfrac{x^2+4x-1}{2x+3}$ có đạo hàm $y'=\dfrac{ax^2+bx+c}{(2x+3)^2}$. Khi  đó giá trị của biểu thức $P=a^2+b^2+c^2$ bằng
	 \shortans{236}
	\loigiai{\allowdisplaybreaks\begin{eqnarray*}
		y'&=&\dfrac{\left(x^2+4x-1\right)'\cdot (2x+3)-(2x+3)'\cdot\left(x^2+4x-1\right)}{(2x+3)^2}\\ 
		&=&\dfrac{(2x+4)(2x+3)-2\left(x^2+4x-1\right)}{(2x+3)^2}\\ 
		 &=&\dfrac{4x^2+14x+12-2x^2-8x+2}{(2x+3)^2}\\
		 &=&\dfrac{2x^2+6x+14}{(2x+3)^2}.
	\end{eqnarray*}
	$  \Rightarrow a=2$; $b=6$; $c=14$.\\
Vậy  $P=a^2+b^2+c^2=236$.}
\end{ex}

% nguồn: Ok_TraLoiNgan_Lop11_Lan1
\begin{ex}%[1D7V2-1]%[ma:4d24a44b]
	Cho hàm số $f(x)=2\sin2x+\cos2x$. Giá trị lớn nhất của hàm số $g(x)=f'(x)+2\sin2x-3\cos2x$ bằng bao nhiêu (kết quả làm tròn đến hàng  phần trăm)?
	 \shortans{2,47}
	\loigiai{Ta có  $f'(x)=(2\sin2x+\cos2x)'=4\cos2x-2\sin2x$.\\
		Do đó  $g(x)=4\sin2x+2\cos2x+2\sin2x-3\cos2x=6\sin2x-\cos2x$.\\
		Áp  dụng bất đẳng thức BCS, ta có\\
		$$|g(x)|\leq\sqrt{(6^2+(-1)^2)(\sin^22x+\cos^22x)}\Leftrightarrow|g(x)|\leq\sqrt{37}.$$ 
		$\max g(x)=\sqrt{37}\Leftrightarrow\dfrac{\sin 2x}{\cos2x}=-6\Leftrightarrow\tan2x=-6$.\\ 
	 Vậy $\max g(x) \approx 2{,}47$.}
\end{ex}

% nguồn: Ok_TraLoiNgan_Lop11_Lan1
\begin{ex}%[1D7V2-1]%[ma:ff2c09d7]
	Cho hàm số $y=f(x)=\sqrt{\tan  x+\cot x}$. Tính $f'\left(\dfrac{\pi}{4}\right)$.
	\shortans{0} 
	 \loigiai{
		\textbf{Cách 1:}\allowdisplaybreaks\begin{eqnarray*} f'(x)&=&\dfrac{1+\tan^2x-\left(1+\cot^2x\right)}{2\sqrt{\tan x+\cot x}}\\
		&=&\dfrac{\tan^2x-\cot^2x}{2\sqrt{\tan x+\cot  x}}.
	\end{eqnarray*}
	Suy ra $f'(0)=\dfrac{\tan^2\dfrac{\pi}{4}-\cot^2\dfrac{\pi}{4}}{2\sqrt{\tan\dfrac{\pi}{4}+\cot\dfrac{\pi}{4}}}=0$.\\ 
\textbf{Cách 2:  }Sử dụng máy tính Casio.}
\end{ex}

% nguồn: Ok_TraLoiNgan_Lop11_Lan2
\begin{ex}%[1D7V2-1]%[ma:588f3bf8]
	Cho hàm số $f(x)=\dfrac{2}{\sqrt{x+2}-\sqrt x}-\dfrac{1}{\sqrt{x+2}-\sqrt{x+1}}$. Tính tổng $T=f'(1)+f'(2)+f'(3)+\cdots+f'(2\,023)$ (kết quả  làm tròn đến hàng phần trăm).
	\shortans{0,49} 
	\loigiai{ 
		Tập xác  định $\mathscr{D}=[0;+\infty)$.\\
		Ta có $f(x)=\dfrac{2}{\sqrt{x+2}-\sqrt x}-\dfrac{1}{\sqrt{x+2}-\sqrt{x+1}}=(\sqrt{x+2}+\sqrt x)-(\sqrt{x+2}+\sqrt{x+1})$  $=\sqrt x-\sqrt{x+1}$.\\
		Suy ra  ${f'}(x)=\dfrac{1}{2}\left(\dfrac{1}{\sqrt x}-\dfrac{1}{\sqrt{x+1}}\right).$\\
		 Vậy
		\begin{eqnarray*}
			&T&=f'(1)+f'(2)+f'(3)+\cdots+f'(2\,023)\\ 
			 & &=\dfrac{1}{2}\left(\dfrac{1}{\sqrt 1}-\dfrac{1}{\sqrt 2}\right)+\dfrac{1}{2}\left(\dfrac{1}{\sqrt 2}-\dfrac{1}{\sqrt 3}\right)+\cdots+\dfrac{1}{2}\left(\dfrac{1}{\sqrt{2\,023}}-\dfrac{1}{\sqrt{2\,024}}\right)\\
			& &=\dfrac{1}{2}\left(\dfrac{1}{\sqrt  1}-\dfrac{1}{\sqrt{2\,024}}\right)\\
			& &=\dfrac{\sqrt{2\,024}-1}{2\sqrt{2\,024}}\approx  0{,}49 .
		\end{eqnarray*}
	} 
\end{ex}

% nguồn: Ok_TraLoiNgan_Lop11_Lan1
\begin{ex}%[1D7B2-1]%[ma:36e22482]
	Cho hàm  số $f(x) = - \dfrac{1}{3}x^3 + 4x^2 - 7x - 11$. Trong đoạn $[0;2025]$, bất phương trình $f'(x) \geqslant 0$ có tất cả bao nhiêu nghiệm nguyên?
	 \shortans{7}
	\loigiai{ 
		Ta có $f'(x)=-x^2+8x-7$, do đó $f'(x) \geq 0 \Leftrightarrow -x^2+8x-7 \geq 0  \Leftrightarrow 1\leq x\leq 7$. \\
		Vậy bất phương trình đã  cho có $7$ nghiệm nguyên thuộc đoạn $[0;2025]$.
	} 
\end{ex}


%%%%%%%%%%%%%%%% TỰ LUẬN %%%%%%%%%%%%%%%%

% nguồn: Ok_TuanLuan_Toan11_Lan1
\begin{bt}%[1D7N2-1]%[1D7H2-1]%[ma:3e000fd2]
	Tính  đạo hàm của các hàm số sau
	\begin{listEX}[2]
		 \item $y=\left(\mathrm{e}^{x}+3 x\right)(\ln \mathrm{x}+1)$;
		 \item $y=x^{5}-3 x^{3}+7 x+2024$.
	\end{listEX}
	\loigiai{ 
		\begin{listEX}[1]
			\item \begin{align*}
			y'= & \left(\mathrm{e}^{x}+3 x\right)'(\ln \mathrm{x}+1) +  \left(\mathrm{e}^{x}+3 x\right)(\ln \mathrm{x}+1)'\\
			= & \left(\mathrm{e}^{x}+3 \right)(\ln \mathrm{x}+1) + \left(\mathrm{e}^{x}+3 x\right)\dfrac{1}{x}\\ 
			= & \mathrm{e}^{x}\ln\mathrm{x} +\mathrm{e}^{x} +3\ln\mathrm{x} +3 +  \dfrac{\mathrm{e}^{x}}{x} +3\\
			 = & \mathrm{e}^{x}\ln\mathrm{x} +\mathrm{e}^{x}+ \dfrac{\mathrm{e}^{x}}{x} +3\ln\mathrm{x} +6.
		\end{align*}
		\item Đạo hàm  của các hàm số $y=x^{5}-3 x^{3}+7 x+2024$ là $y'=5x^{4}-9 x^{2}+7$.
	\end{listEX}
 }
\end{bt}

% nguồn: Ok_TuanLuan_Toan11_Lan2
\begin{ex}%[1D7N2-1]%[1D7H2-1]%[ma:ebf252cd]
	Tính đạo hàm  của các hàm số sau
	\begin{listEX}[2]
		 \item $y=\left(\mathrm{e}^{x}+3 x\right)(\ln \mathrm{x}+1)$;
		 \item $y=x^{5}-3 x^{3}+7 x+2024$.
	\end{listEX}
	 \loigiai{
		\begin{listEX}[1]
			\item \begin{align*}
			y'= & \left(\mathrm{e}^{x}+3 x\right)'(\ln \mathrm{x}+1) + \left(\mathrm{e}^{x}+3  x\right)(\ln \mathrm{x}+1)'\\
			= &  \left(\mathrm{e}^{x}+3 \right)(\ln \mathrm{x}+1) + \left(\mathrm{e}^{x}+3 x\right)\dfrac{1}{x}\\
			=  & \mathrm{e}^{x}\ln\mathrm{x} +\mathrm{e}^{x} +3\ln\mathrm{x} +3 + \dfrac{\mathrm{e}^{x}}{x} +3\\
			= &  \mathrm{e}^{x}\ln\mathrm{x} +\mathrm{e}^{x}+ \dfrac{\mathrm{e}^{x}}{x} +3\ln\mathrm{x} +6.
		\end{align*}
		\item Đạo hàm của các hàm số $y=x^{5}-3 x^{3}+7 x+2024$ là $y'=5x^{4}-9  x^{2}+7$.
	\end{listEX}
} 
\end{ex}

% nguồn: Ok_TuanLuan_Toan11_Lan1
\begin{bt}%[1D7H2-1]%[ma:774a4359]
	Tính đạo hàm của hàm số $y=5^x+\log_2x+\ln(2x)+\mathrm{e}^{-x}+\sqrt{3}$. 
	\loigiai{ 
		Ta  có $y'=5^x\cdot\ln5+\dfrac{1}{x\ln2}+\dfrac{1}{x}-\mathrm{e}^{-x}$.
	} 
\end{bt}

% nguồn: Ok_TuanLuan_Toan11_Lan1
\begin{bt}%[1D7H2-1]%[ma:cece1e39]
	Cho hàm số $f(x)=\sqrt{x^2+3}$.  Tính giá trị của biểu thức $S=f(1)+4 f^{\prime}(1)$.
	 \loigiai{
		Ta có $f^{\prime}(x)=\dfrac{x}{\sqrt{x^2+3}}\Rightarrow f^{\prime}(1)=\dfrac{1}{2}$, $f(1)=2\Rightarrow S=f(1)+4 f^{\prime}(1)=2+4\cdot\dfrac{1}{2}=4$. 
	} 
\end{bt}

% nguồn: Ok_TuanLuan_Toan11_Lan2
\begin{ex}%[1D7H2-1]%[ma:e692f39f]
	Tính đạo hàm của hàm số  $y=5^x+\log_2x+\ln(2x)+\mathrm{e}^{-x}+\sqrt{3}$.
	 \loigiai{
		Ta có $y'=5^x\cdot\ln5+\dfrac{1}{x\ln2}+\dfrac{1}{x}-\mathrm{e}^{-x}$. 
	} 
\end{ex}

% nguồn: Ok_TuanLuan_Toan11_Lan2
\begin{ex}%[1D7H2-1]%[ma:c19d3d1f]
	Cho hàm số $f(x)=\sqrt{x^2+3}$.  Tính giá trị của biểu thức $S=f(1)+4 f^{\prime}(1)$.
	\loigiai{ 
		Ta  có $f^{\prime}(x)=\dfrac{x}{\sqrt{x^2+3}}\Rightarrow f^{\prime}(1)=\dfrac{1}{2}$, $f(1)=2\Rightarrow S=f(1)+4 f^{\prime}(1)=2+4\cdot\dfrac{1}{2}=4$.
	 }
\end{ex}

% nguồn: Ok_TuanLuan_Toan11_Lan2
\begin{ex}%[1D7H2-1]%[ma:d1bfab33]
	Tính đạo hàm của các  hàm số sau
	\begin{enumEX}{2}
		\item $ y=4x^2-\sqrt  x+\dfrac{1}{x}$;
		\item $ y=x-\dfrac{1}{\sqrt  x}-\dfrac{1}{\sqrt[3]{x}}+1$;
		 \item $ y=\left(x+\dfrac{1}{\sqrt x}\right)^2+\left(x+\dfrac{1}{x}\right)^2$;
		\item  $ y=x+\sqrt x+\sqrt[3]{x}+\sqrt[4]{x}+\sqrt[5]{x}$.
	\end{enumEX}
	 \loigiai{
		\begin{enumerate}
			\item Ta có  $ y'=8x-\dfrac{1}{2\sqrt x}-\dfrac{1}{x^2}$.
			\item  Ta có $ y'=1+\dfrac{1}{2\sqrt{x^3}}+\dfrac{1}{3\sqrt[3]{x^4}}$.
			\item Ta có $ y=\left(x+\dfrac{1}{\sqrt  x}\right)^2+\left(x+\dfrac{1}{x}\right)^2=x+2+\dfrac{1}{x}+x^2+2+\dfrac{1}{x^2}\Rightarrow y'=1-\dfrac{1}{x^2}+2x-\dfrac{2}{x^3}$.
			\item  Ta có $ y'=1+\dfrac{1}{2\sqrt x}+\dfrac{1}{3\sqrt[3]{x^2}}+\dfrac{1}{4\sqrt[4]{x^3}}+\dfrac{1}{5\sqrt[5]{x^4}}$.
	\end{enumerate}}
\end{ex}

% nguồn: Ok_TuanLuan_Toan11_Lan2
\begin{ex}%[1D7H2-1]%[ma:1b6775d3]
	Tính đạo hàm  của hàm số $y=\dfrac{3^x-3^{-x}}{3^x+3^{-x}}$.
	 \loigiai{Ta có
	\begin{eqnarray*}
		&y'&=\dfrac{\left(3^x \ln 3+3^{-x} \ln 3\right)\left(3^x+3^{-x}\right)-\left(3^x-3^{-x}\right)\left(3^x \ln 3-3^{-x} \ln 3\right)}{\left(3^x+3^{-x}\right)^2}\\ & &=\dfrac{\left(3^x+3^{-x}\right)^2-\left(3^x-3^{-x}\right)^2}{\left(3^x+3^{-x}\right)^2}\cdot \ln 3\\& &=\dfrac{4 \ln 3}{\left(3^x+3^{-x}\right)^2}. 
	\end{eqnarray*}
	} 
\end{ex}

% nguồn: Ok_TuanLuan_Toan11_Lan2
\begin{ex}%[1D7H2-1]%[ma:923ea6e8]
	Tìm đạo  hàm của các hàm số sau
	\begin{enumEX}{2}
		 \item $y=5\sin x-3\cos x$;
		 \item $y=\sin\left(x^2-3x+2\right)$;
		 \item $ y=\sqrt{1+2\tan x}$;
		\item $y=\tan 3x-\cot  3x$.
	\end{enumEX}
	 \loigiai{
		\begin{enumerate}
			\item Ta có $y'=5\cos x+3\sin  x$.
			\item  Ta có $y'=\left(x^2-3x+2\right)'\cdot \cos\left(x^2-3x+2\right)$ $=\left(2x-3\right)\cdot \cos \left(x^2-3x+2\right)$.
			\item Ta  có $y'=\dfrac{\left(2\tan x\right)'}{2\sqrt{1+2\tan x}}$ $=\dfrac{\dfrac{2}{\cos^2x}}{2\sqrt{1+2\tan x}}$ $=\dfrac{1}{\cos^2x\cdot \sqrt{1+2\tan x}}$.
			\item Ta có thể làm bằng  hai cách sau
			\begin{itemize}
				\item \textbf{Cách 1.} Ta có  ngay
				 $$y'=\dfrac{3}{\cos^23x}+\dfrac{3}{\sin^23x}=\dfrac{3}{\sin^23x\cdot \cos^23x} =\dfrac{3}{\tfrac{1}{4}{\sin^2}6x}=\dfrac{12}{\sin^26x}.$$
				\item \textbf{Cách 2.} Ta biến  đổi
				$$y=\dfrac{\sin 3x}{\cos 3x}-\dfrac{\cos  3x}{\sin 3x}=\dfrac{\sin^23x-\cos^23x}{\cos 3x\cdot \sin 3x}=-\dfrac{2\cos 6x}{\sin 6x}=-2\cot 6x.$$
				Vậy  $y'=\dfrac{12}{\sin^26x}$.
			\end{itemize}
		\end{enumerate}
	 }
\end{ex}

% nguồn: Ok_TuanLuan_Toan11_Lan2
\begin{ex}%[1D7H2-1]%[ma:caa159f6]
	Cho  hàm số $f(x)=\mathrm{e}^{\sqrt{x^2+1}}$. Tính $f'(1)$.
	\loigiai{Ta có 
	 $$
	f(x)=\mathrm{e}^{\sqrt{x^2+1}} \Rightarrow f'(x)=\dfrac{x  \cdot \mathrm{e}^{\sqrt{x^2+1}}}{\sqrt{x^2+1}}.
	$$ 
	Vậy  $f'(1)=\dfrac{\mathrm{e}^{\sqrt{2}}}{\sqrt{2}}$.
	} 
\end{ex}

% nguồn: Ok_TuanLuan_Toan11_Lan2
\begin{ex}%[1D7H2-1]%[ma:ada3c31f]
	Tính đạo hàm của hàm số  $y=x \mathrm{e}^{-x}$.
	 \loigiai{Ta có
	$$ 
	y=x \mathrm{e}^{-x} \Rightarrow y'=\mathrm{e}^{-x}-x \cdot \mathrm{e}^{-x}=(1-x) \cdot \mathrm{e}^{-x}. 
	$$ 
	 }
\end{ex}

% nguồn: Ok_TuanLuan_Toan11_Lan2
\begin{ex}%[1D7H2-1]%[ma:4576a869]
	Tính đạo  hàm của hàm số $y=2^{x^2+2}$.
	 \loigiai{
		Ta có 
		 $$
		y=2^{x^2+2} \Rightarrow y'=\left(x^2+2\right)' \cdot 2^{x^2+2} \cdot \ln 2=2 x \cdot 2^{x^2+2} \cdot \ln  2=x \cdot 2^{x^2+3} \cdot \ln 2.
	$$} 
\end{ex}

% nguồn: Ok_TuanLuan_Toan11_Lan2
\begin{ex}%[1D7H2-1]%[ma:63c7f4b9]
	Tính  đạo hàm của các hàm số sau
	\begin{enumEX}{2}
		\item $  y=\left(1-x\right)\left(1-2x\right)\left(1-3x\right)$;
		\item  $ y=\left(x+\sqrt x\right)\left(x^2+x+1\right)$;
		\item $ y=x^2\left(x+4\right)^3$; 
		\item $ y=\left(1-\dfrac{1}{x}\right)\left(x-\dfrac{1}{x^2}\right)$; 
		\item $ y=\left(x^3+3x\right)\left(2-x\right)$. 
	\end{enumEX}
	\loigiai{ 
		\begin{enumerate}
			\item  Ta có
			\begin{eqnarray*}
				 &y&=\left(1-x\right)\left(1-2x\right)\left(1-3x\right)\\
				& &=\left(1-3x+2x^2\right)\left(1-3x\right)\\ &  &=1-3x-3x+9x^2+2x^2-6x^3\\
				 & &=1-6x+11x^2-6x^3.
			\end{eqnarray*}
			Suy  ra $y'=-6+22x-18x^2$.
			\item Ta có $y=\left(x+\sqrt x\right)\left(x^2+x+1\right)=x^3+x^2+x-x^2\sqrt x-x\sqrt x-\sqrt x $.\\ 
			Suy ra $ y'=3x^2+2x+1-\dfrac{5}{2}\sqrt{x^3}-\dfrac{3}{2}\sqrt x-\dfrac{1}{2\sqrt x}$. 
			\item Ta có  $ y=x^2\left(x+4\right)^3=x^2\left(x^3+12x^2+48x+64\right)=x^5+12x^4+48x^3+64x^2$.\\
			 Suy ra $y'=5x^4+48x^3+144x^2+128x$.
			\item Ta  có $ y=\left(1-\dfrac{1}{x}\right)\left(x-\dfrac{1}{x^2}\right)=x-\dfrac{1}{x^2}-1+\dfrac{1}{x^3}\Rightarrow y'=1+\dfrac{2}{x^3}-\dfrac{3}{x^4}$.
			\item Ta có  $ y=\left(x^3+3x\right)\left(2-x\right)=2x^3-x^4+6x-3x^2\Rightarrow y'=6x^2-4x^3+6-6x$.
		\end{enumerate}
	} 
\end{ex}

% nguồn: Ok_TuanLuan_Toan11_Lan2
\begin{ex}%[1D7H2-1]%[ma:ba0677cd]
	 Tính đạo hàm của các hàm số sau
	\begin{enumEX}{3}
		 \item $ y=\dfrac{2}{x}$;
		\item $  y=\dfrac{x+2}{2x-1}$;
		\item $  y=\dfrac{x^2+4x-1}{2x+3}$;
		 \item $ y=\dfrac{\sqrt x}{x+1}$;
		\item $ y=\dfrac{1+x-x^2}{1-x+x^2}$; 
		\item $ y=\dfrac{3}{2x-1}$. 
	\end{enumEX}
	\loigiai{ 
		\begin{enumerate}
			\item Ta  có $ y'=\dfrac{-2}{x^2}$.
			\item Ta có $ y'=\dfrac{2x-1-2\left(x+2\right)}{\left(2x-1\right)^2}=\dfrac{5}{\left(2x-1\right)^2}$. 
			\item Ta có $y'=\dfrac{\left(2x+3\right)\left(2x+4\right)-2\left(x^2+4x-1\right)}{\left(2x+3\right)^2}=\dfrac{4x^2+14x+12-2x^2-8x+2}{\left(2x+3\right)^2}=\dfrac{2x^2+6x+14}{\left(2x+3\right)^2}$. 
			\item  Ta có $y'=\dfrac{\dfrac{1}{2\sqrt x}\left(x+1\right)-\sqrt x}{\left(x+1\right)^2}=\dfrac{1-x}{2\sqrt x{\left(x+1\right)^2}}$.
			 \item Ta có $ y=\dfrac{2}{1-x+x^2}-1\Rightarrow y'=-\dfrac{2\left(-1+2x\right)}{\left(1-x+x^2\right)^2}=\dfrac{2\left(1-2x\right)}{\left(1-x+x^2\right)^2}$.
			\item Ta  có $ y'=-\dfrac{6}{\left(2x-1\right)^2}$.
	\end{enumerate}}
\end{ex}

% nguồn: Ok_TuanLuan_Toan11_Lan2
\begin{ex}%[1D7H2-1]%[ma:9c66f512]
	Tính đạo hàm cùa  hàm số $y=\mathrm{e}^{x^2-2} \cos x$.
	\loigiai{Ta có 
	 $$
	y=\mathrm{e}^{x^2-2} \cos x \Rightarrow y'=2 x \cdot \mathrm{e}^{x^2-2} \cos x-\mathrm{e}^{x^2-2} \sin x=(2 x \cos x-\sin x)\cdot  \mathrm{e}^{x^2-2}.
	 $$
	 }
\end{ex}

% nguồn: Ok_TuanLuan_Toan11_Lan2
\begin{ex}%[1D7H2-1]%[ma:2f22aead]
	Cho hàm số $y=\ln (\cos x)$. Với điều kiện hàm  số đã cho xác định, tìm đạo hàm của hàm số đó.
	 \loigiai{
	 Ta có $y'=\dfrac{(\cos x)'}{\cos x}=\dfrac{-\sin x}{\cos x}=-\tan x$.
	} 
\end{ex}

% nguồn: Ok_TuanLuan_Toan11_Lan2
\begin{ex}%[1D7H2-1]%[ma:a3cef07e]
	Tính đạo  hàm của hàm số $y=\left(x^2+2 x\right) \mathrm{e}^x$.
	\loigiai{Ta có 
	$$ 
	y=\left(x^2+2 x\right) \mathrm{e}^x \Rightarrow y'=(2 x+2) \cdot  \mathrm{e}^x+\left(x^2+2 x\right) \cdot \mathrm{e}^x=\left(x^2+4 x+2\right) \cdot \mathrm{e}^x.
	$$ 
	 }
\end{ex}

% nguồn: Ok_TuanLuan_Toan11_Lan2
\begin{ex}%[1D7V2-1]%[ma:540b649c]
	Tính đạo hàm  của các hàm số sau
	\begin{enumEX}{2}
		\item $ y=\left(x+\sqrt  x+\sqrt[3]{x}+\sqrt[4]{x}+\sqrt[5]{x}\right)^5$;
		\item  $ y=\left(\sqrt{x^2-\sqrt x}+\sqrt[3]{1-\dfrac{x+1}{x^2}}\right)^7$;
		\item $  y=\left(x^3+2\right)\sqrt[3]{1-x^2}+\sqrt{x^3-1}$;
		\item $ y=\dfrac{2x-\dfrac{1}{x^2}+1}{x+\sqrt x+\sqrt{x^3-1}}$; 
		 \item $ y=\dfrac{x}{\sqrt{x^2+x+1}}$;
		\item  $ y=\left(\sqrt[3]{x}+2\right)\left(1+\sqrt[3]{x^2}+3x\right)$.
	\end{enumEX}
	 \loigiai{
		\begin{enumerate}
			 \item Ta có $ y'=5\left(1+\dfrac{1}{2\sqrt x}+\dfrac{1}{3\sqrt[3]{x^2}}+\dfrac{1}{4\sqrt[4]{x^3}}+\dfrac{1}{5\sqrt[5]{x^4}}\right){\left(x+\sqrt x+\sqrt[3]{x}+\sqrt[4]{x}+\sqrt[5]{x}\right)^4}$.
			\item Ta có 
			\begin{eqnarray*}
				 &y'&=7\left(\sqrt{x^2-\sqrt x}+\sqrt[3]{1-\dfrac{x+1}{x^2}}\right)'{\left(\sqrt{x^2-\sqrt x}+\sqrt[3]{1-\dfrac{x+1}{x^2}}\right)^6}\\
				& &=7\left(\dfrac{\left(x^2+\sqrt  x\right)'}{2\sqrt{x^2-\sqrt x}}+\dfrac{\left(1-\dfrac{x+1}{x^2}\right)'}{3\sqrt[3]{\left(1-\dfrac{x+1}{x^2}\right)^2}}\right){\left(\sqrt{x^2-\sqrt x}+\sqrt[3]{1-\dfrac{x+1}{x^2}}\right)^6}\\
				& &=7\left(\dfrac{2x-\dfrac{1}{2\sqrt x}}{2\sqrt{x^2-\sqrt  x}}+\dfrac{-\left(\dfrac{x^2-2x\left(x+1\right)}{x^4}\right)}{3\sqrt[3]{\left(1-\dfrac{x+1}{x^2}\right)^2}}\right){\left(\sqrt{x^2-\sqrt x}+\sqrt[3]{1-\dfrac{x+1}{x^2}}\right)^6}\\
				&  &=7\left(\dfrac{2x-\dfrac{1}{2\sqrt x}}{2\sqrt{x^2-\sqrt x}}+\dfrac{\dfrac{-x-2}{x^3}}{3\sqrt[3]{\left(1-\dfrac{x+1}{x^2}\right)^2}}\right){\left(\sqrt{x^2-\sqrt x}+\sqrt[3]{1-\dfrac{x+1}{x^2}}\right)^6}.
			\end{eqnarray*}
			\item Ta có \begin{eqnarray*} &y'&=\left(x^3+2\right)'\sqrt[3]{1-x^2}+\left(x^3+2\right){\left(\sqrt[3]{1-x^2}\right)'}+\left(\sqrt{x^3-1}\right)'\\
			&  &=3x^2\sqrt[3]{1-x^2}+\left(x^3+2\right)\cdot \dfrac{-2x}{3\cdot \sqrt[3]{\left(1-x^2\right)^2}}+\dfrac{3x^2}{2\sqrt{x^3-1}}.
		\end{eqnarray*}
		 \item Ta có
		\begin{eqnarray*}
			&y'&=\dfrac{\left(2x-\dfrac{1}{x^2}+1\right)'\left(x+\sqrt x+\sqrt{x^3-1}\right)-\left(2x-\dfrac{1}{x^2}+1\right){\left(x+\sqrt x+\sqrt{x^3-1}\right)'}}{\left(x+\sqrt  x+\sqrt{x^3-1}\right)^2}\\
			& &=\dfrac{\left(2+\dfrac{2}{x^3}\right)\left(x+\sqrt x+\sqrt{x^3-1}\right)-\left(2x-\dfrac{1}{x^2}+1\right)\left(1+\dfrac{1}{2\sqrt x}+\dfrac{3x^2}{2\sqrt{x^3-1}}\right)}{\left(x+\sqrt  x+\sqrt{x^3-1}\right)^2}.
		\end{eqnarray*}
		\item Ta  có
		\begin{eqnarray*}
			&y'&=\dfrac{\sqrt{x^2+x+1}-x\cdot \dfrac{2x+1}{2\sqrt{x^2+x+1}}}{\left(\sqrt{x^2+x+1}\right)^2}\\ 
			& &=\dfrac{2x^2+2x+2-2x^2-x}{2\sqrt{x^2+x+1}\left(x^2+x+1\right)}=\dfrac{x+2}{2\sqrt{x^2+x+1}\left(x^2+x+1\right)}. 
		\end{eqnarray*}
\end{enumerate}}
\end{ex}

% nguồn: Ok_TuanLuan_Toan11_Lan2
\begin{ex}%[1D7V2-1]%[ma:3fceba1b]
	Tính đạo hàm  của hàm số $y=\cos x \cdot \mathrm{e}^{\tan x}$.
	 \loigiai{ Ta có
	 $$
	y=\cos x \cdot \mathrm{e}^{\tan x} \Rightarrow y'=-\sin x \cdot \mathrm{e}^{\tan x}+\cos x \cdot \dfrac{1}{\cos ^2 x} \cdot \mathrm{e}^{\tan x}=\mathrm{e}^{\tan x}\left(\dfrac{1}{\cos x}-\sin  x\right).
	$$ 
	 }
\end{ex}

% nguồn: Ok_TuanLuan_Toan11_Lan2
\begin{ex}%[1D7V2-1]%[ma:d8384c54]
	 Tính đạo hàm của hàm số
	$ y=\sqrt{\dfrac{1}{2}+\dfrac{1}{2}\sqrt{\dfrac{1}{2}+\dfrac{1}{2}\sqrt{\dfrac{1}{2}+\dfrac{1}{2}\cos x}}}$ với  $x\in (0;\pi)$.
	 \loigiai{
		 Biến đổi hàm số, ta được
		\begin{eqnarray*}
			&y&=\sqrt{\dfrac{1}{2}+\dfrac{1}{2}\sqrt{\dfrac{1}{2}+\dfrac{1}{2}\sqrt{\dfrac{1}{2}+\dfrac{1}{2}\cos x}}}\\ 
			 & &=\sqrt{\dfrac{1}{2}+\dfrac{1}{2}\sqrt{\dfrac{1}{2}+\dfrac{1}{2}\sqrt{\cos^2\dfrac{x}{2}}}}\\
			& &= \sqrt{\dfrac{1}{2}+\dfrac{1}{2}\sqrt{\dfrac{1}{2}+\dfrac{1}{2}\cos\dfrac{x}{2}}}\\ 
			&  &=\sqrt{\dfrac{1}{2}+\dfrac{1}{2}\sqrt{\cos^2\dfrac{x}{4}}}\\
			&  &=\sqrt{\dfrac{1}{2}+\dfrac{1}{2}\cos\dfrac{x}{4}}=\sqrt{\cos^2\dfrac{x}{8}}=\cos\dfrac{x}{8}
		\end{eqnarray*}
	Do đó  $ y'=\left(\cos\dfrac{x}{8}\right)'=-\dfrac{1}{8}\sin\dfrac{x}{8}$.}
\end{ex}

% nguồn: Ok_TuanLuan_Toan11_Lan2
\begin{ex}%[1D7V2-1]%[ma:4e74d819]
	Cho hàm  số $y=\ln \left(x^2+\sqrt{x^2+1}\right)$. Tìm đạo hàm của hàm số đã cho.
	\loigiai{Ta có 
	$y'=\dfrac{\left(x^2+\sqrt{x^2+1}\right)'}{x^2+\sqrt{x^2+1}}=\dfrac{2 x+\frac{\left(x^2+1\right)'}{2 \sqrt{x^2+1}}}{x^2+\sqrt{x^2+1}}=\dfrac{2 x+\frac{x}{\sqrt{x^2+1}}}{x^2+\sqrt{x^2+1}}=\dfrac{x\left(2  \sqrt{x^2+1}+1\right)}{\left(x^2+\sqrt{x^2+1}\right) \sqrt{x^2+1}}$.
	} 
\end{ex}
